% lecture 1-2

\newpage

\pagenumbering{arabic}
\setcounter{page}{1}

\chapter{Arithmetic and Geometry of Complex Numbers}


%Historically, complex numbers arose from solving polynomial equations. Very often we refer to them as {\em imaginary} numbers, because the roots of polynomial may not be real numbers. However, we can construct complex numbers from real numbers explicitly, and the computation of complex numbers is very concrete. There is nothing imaginary.


\section{Introduction: from natural numbers to complex numbers} 

We begin with a famous remark of the German mathematician Leopold Kronecker:

\begin{quote}
God created the natural numbers, the rest is the work of man.
\end{quote}

The natural numbers $\mathbb{N}$ allow us to perform addition and multiplication, but they are not closed under subtraction or division. For instance, the equations $x+5=3$ and $2x=3$ have no solutions in $\mathbb{N}$. To overcome this limitation, we enlarge $\mathbb{N}$ first to the integers $\mathbb{Z}$, and then to the rational numbers $\mathbb{Q}$. Let us briefly recall how $\mathbb{Q}$ can be constructed from~$\mathbb{Z}$.


A rational number is usually written as $a/b$, where $a$ and $b$ are integers with $b\neq 0$. Formally, we represent such a number by an ordered pair $(a,b)$ of integers. Two pairs $(a,b)$ and $(c,d)$ represent the same rational number if and only if $ad = bc$. Thus a rational number is, strictly speaking, an equivalence class of such pairs.

Addition and multiplication are defined by
\[
(a,b) + (c,d) \triangleq (ad + bc,\, bd), \qquad
(a,b)\cdot(c,d) \triangleq (ac,\, bd).
\]
For example, the rational numbers $1/2$ and $1/3$ are represented by $(1,2)$ and $(1,3)$, and their sum is
\[
(1,2) + (1,3) = (1\cdot 3 + 2\cdot 1,\, 2\cdot 3) = (5,6).
\]
Pairs of the form $(k,1)$ behave exactly like the integer $k$, so $\mathbb{Z}$ naturally embeds into this system. Moreover, every nonzero rational number has a multiplicative inverse: if $a\neq 0\neq b$, then
\[
(a,b)\cdot(b,a) = (ab,ab),
\]
which is equivalent to $(1,1)$.

Although $\mathbb{Q}$ is closed under the usual algebraic operations, it is not complete in an analytic sense. In particular, it is not closed under taking supremum and infimum. Recall that the \emph{supremum} of a set $A$ is the least upper bound of $A$, and the \emph{infimum} is the greatest lower bound. A classical example shows the failure of completeness: the set
\[
\{x\in \mathbb{Q} : x^2 < 2\}
\]
has no smallest upper bound in $\mathbb{Q}$.

To remedy this, we complete $\mathbb{Q}$ with respect to the Euclidean metric. This can be done using Dedekind cuts or Cauchy sequences, and the resulting number system is the field of real numbers $\mathbb{R}$. By construction, $\mathbb{R}$ has the property that every nonempty subset bounded above has a supremum, every Cauchy sequence converges, and $\mathbb{Q}$ sits densely inside $\mathbb{R}$.

However, even $\mathbb{R}$ is not the end of the story. The real numbers are not algebraically closed: a polynomial
\[
c_0 + c_1 x + c_2 x^2 + \cdots + c_d x^n
\]
with real coefficients need not have a real root. The simplest example is $x^2 + 1$, which has no solution in $\mathbb{R}$. The remarkable fact is that adjoining a single new element $i$, defined formally by $i^2 = -1$,
is enough to remedy this deficiency. The resulting number system,
\[
\mathbb{C} = \{\, a + bi : a,b \in \mathbb{R} \,\},
\]
is called the \emph{complex numbers}. It contains $\mathbb{R}$ as a subfield 
(via the identification $a \leftrightarrow a + 0i$), and it turns out to be 
\emph{algebraically closed}: every nonconstant polynomial with complex coefficients 
has a root in $\mathbb{C}$.

This property is known as the \emph{Fundamental Theorem of Algebra}. It is not 
obvious from the definition of $\mathbb{C}$, and we will give a proof using complex analysis. Nevertheless, the theorem justifies the introduction of complex numbers as the natural setting for algebraic equations.

%In the next section, we will develop the algebraic and geometric structure of 
%$\mathbb{C}$, and see how complex numbers can be viewed both as ordered pairs of 
%real numbers and as points in the plane.





\section{Models of complex field}



The rational numbers $\mathbb{Q}$ and the real numbers $\mathbb{R}$ share a common algebraic structure. In abstract algebra, this structure is called a \emph{field}.



\begin{svgraybox}
	\begin{definition} \index{Field} \index{Field!Extension}
		A number system $(F,+,\cdot)$ is called a \df{field} if it satisfies the following axioms:
\begin{enumerate}
	\item (closure) $a+b\in F$ for all $a,b\in F$.
	\item (associativity) $(a+b)+c = a+(b+c)$ for all $a,b,c\in F$.
	\item (commutativity) $a+b = b+a$ for all $a,b\in F$.
	\item (existence of zero) There exists $0\in F$ such that $0+a = a+0 = a$ for all $a\in F$.
	\item (additive inverse) For every $a\in F$, there exists $a'\in F$ such that $a+a'=0$.
	\item (closure) $a\cdot b\in F$ for all $a,b\in F$.
	\item (associativity) $(a\cdot b)\cdot c = a\cdot (b\cdot c)$ for all $a,b,c\in F$.
	\item (commutativity) $a\cdot b = b\cdot a$ for all $a,b\in F$.
	\item (existence of one) There exists $1\in F$ with $1\neq 0$ such that $1\cdot a = a\cdot 1 = a$ for all $a\in F$.
	\item (multiplicative inverse) For every $a\in F\setminus\{0\}$, there exists $a''\in F$ such that $a\cdot a''=1$.
	\item (distributivity) $a\cdot (b+c) = a\cdot b + a\cdot c$ for all $a,b,c\in F$.
\end{enumerate}

A subset $K$ of a field $F$ is called a \df{subfield} if the elements of $K$ satisfy all the field axioms. In this case, $F$ is called an \df{extension} of $K$.
\label{def:field}
\end{definition}
\end{svgraybox}


We require a “number system’’ to be a field so that addition, subtraction, multiplication, and division behave according to these axioms.



\begin{svgraybox}
	\begin{definition} \index{Complex ! Field}
		A \df{complex field} is a field that extends $\mathbb{R}$ and contains a special element $i$ satisfying $i^2=-1$, and is generated by $\mathbb{R}$ and $i$. Such a field is denoted by~$\mathbb{C}$.
	\end{definition}
\end{svgraybox}

The special number $i$ that satisfies $i^2=-1$ is called the \df{imaginary unit}.


\begin{remark}
The phrase “generated by $\mathbb{R}$ and $i$’’ means that there is no proper subfield containing all real numbers together with the element $i$. For example, consider the field of rational functions $f(T)/g(T)$, where $f(T)$ and $g(T)$ are polynomials in $T$ with complex coefficients. This field contains $\mathbb{R}$ and $i$, but it is far too large to be considered the complex field: mixing real numbers with $i$ never produces the transcendental variable $T$.
\end{remark}




\subsection{Complex numbers as points in the plane}
\label{sec:complex_number_as_vector}
\index{Complex ! Plane} \index{Argand diagram}

One way to construct a complex field is to regard each ordered pair $(x,y)$ of real numbers as a number. The set of all such pairs forms a plane called the \df{complex plane} or \df{Argand diagram}. The point $(x,y)$ may also be viewed as a vector, and addition is defined by vector addition (Fig.~\ref{fig:argang_diagram}).

Define
\[
F_1 \triangleq \{(a,b):\, a,b\in\mathbb{R}\}.
\]
Two pairs $(a,b)$ and $(c,d)$ are \df{equal} if $a=c$ and $b=d$.

Addition and multiplication in $F_1$ are defined by
\[
(a,b)+(c,d) \triangleq  (a+c,\, b+d), \qquad
(a,b)\cdot(c,d) \triangleq  (ac-bd,\, ad+bc).
\]

The real numbers embed into $F_1$ via $x \mapsto (x,0)$. Under this identification,
\[
(x_1,0)+(x_2,0) = (x_1+x_2,0), \qquad
(x_1,0)\cdot(x_2,0) = (x_1x_2,0),
\]
so the horizontal axis of the complex plane corresponds to the real line.

The element $(1,0)$ is the multiplicative identity. The element $(0,1)$ satisfies
\[
(0,1)\cdot(0,1) = (-1,0),
\]
so it plays the role of the imaginary unit $i$.

Every element $(x,y)$ can be written as
\[
(x,y) = x(1,0) + y(0,1) = x + iy.
\]
This shows that $1=(1,0)$ and $i=(0,1)$ form a basis of this number system over $\mathbb{R}$.

\begin{figure}
	\begin{center}
\begin{tikzpicture}[>=Latex, scale=1.2]
	% Axes
	\draw[->] (-0.5,0) -- (3.5,0) node[below right] {$\real$};
	\draw[->] (0,-0.5) -- (0,2.6) node[above left] {$\imag$};
	
	% Tick marks and labels (optional)
	\foreach \x in {1,2,3}
	\draw (\x,0.05) -- (\x,-0.05) node[below] {\scriptsize $\x$};
	\foreach \y in {1,2}
	\draw (0.05,\y) -- (-0.05,\y) node[left] {\scriptsize $\y$};
	
	% Point z = a + bi
	\coordinate (O) at (0,0);
	\coordinate (Z) at (2.5,1.6); % choose some coordinates for illustration
	
	% Vector from origin to z
	\draw[->, thick, blue] (O) -- (Z) node[above right] {$z = a + bi$};
	
	% Projections to axes (dotted)
	\draw[dashed] (Z) -- (2.5,0) node[below] {$a$};
	\draw[dashed] (Z) -- (0,1.6) node[left] {$b$};
	
	% Mark the point
%	\fill[blue] (Z) circle (1.5pt);
\end{tikzpicture}
	\end{center}
\caption{Argand diagram}
\label{fig:argang_diagram}	
\end{figure}

Multiplication by $i$ has the geometric meaning of rotation by 90 degrees (Fig.~\ref{fig:multiply_by_i}).


\begin{figure}
	\begin{center}
		
		
		\begin{tikzpicture}[>=Latex, scale=1.3]
			% Axes
			\draw[->] (-0.5,0) -- (3.5,0) node[below right] {$\real$};
			\draw[->] (0,-0.5) -- (0,3.5) node[above left] {$\imag$};
			
			% Origin
			\coordinate (O) at (0,0);
			
			% z and iz
			\coordinate (Z) at (3,1);   % z
			\coordinate (IZ) at (-1,3); % iz = (-b, a) if z = a+bi
			
			% Vector for z
			\draw[->, thick, blue] (O) -- (Z) node[below right] {$z$};
%			\fill[blue] (Z) circle (1.4pt);
			
			% Vector for iz
			\draw[->, thick, red] (O) -- (IZ) node[above left] {$iz$};
%			\fill[red] (IZ) circle (1.4pt);
			
			% Angle marking (z)
			\draw[->, thick] (1,0) arc[start angle=0, end angle=18.4, radius=1cm];
			\node at (1.3,0.25) {\scriptsize $\theta$};
			
			% Angle marking (iz)
			\draw[->, thick] (1,0) arc[start angle=0, end angle=120, radius=0.8cm];
			\node at (0.4, 1.0) {\scriptsize $\theta + \frac{\pi}{2}$};
		\end{tikzpicture}
	\end{center}
	\caption{Multiply $z=a+bi$ by $i$ gives $iz=-b+ai$.}
	\label{fig:multiply_by_i}
\end{figure}



\subsection{Complex numbers as matrices}
\label{sec:complex_numbers_matrices}
In the second model, a complex number is represented by a $2\times 2$ real matrix:
\[
F_2 \triangleq \left\{
\begin{bmatrix} a & -b \\ b & a \end{bmatrix} : a,b\in\mathbb{R}
\right\}.
\]
Addition and multiplication are the usual matrix operations. The arithmetic mirrors that of the first model:
\[
\begin{bmatrix} a & -b \\ b & a \end{bmatrix}
\begin{bmatrix} c & -d \\ d & c \end{bmatrix}
=
\begin{bmatrix} ac-bd & -(ad+bc) \\ ad+bc & ac-bd \end{bmatrix}.
\]

A real number $x$ corresponds to the diagonal matrix
\[
\begin{bmatrix} x & 0 \\ 0 & x \end{bmatrix}.
\]

The matrix
\[
\begin{bmatrix} 0 & -1 \\ 1 & 0 \end{bmatrix}
\]
satisfies
\[
\begin{bmatrix} 0 & -1 \\ 1 & 0 \end{bmatrix}^2
=
\begin{bmatrix} -1 & 0 \\ 0 & -1 \end{bmatrix},
\]
so it plays the role of the imaginary unit. Every matrix in $F_2$ decomposes as
\[
\begin{bmatrix} x & -y \\ y & x \end{bmatrix}
=
x\begin{bmatrix}1&0\\0&1\end{bmatrix}
+
y\begin{bmatrix}0&-1\\1&0\end{bmatrix}.
\]
The first part $\begin{bmatrix}x & 0 \\0 & x \end{bmatrix}$ is a diagonal matrix, while the second part $\begin{bmatrix}0 & -y \\ y & 0 \end{bmatrix}$ is a skew symmetric matrix.



\subsection{Complex numbers as polynomials}

Let $\mathbb{R}[T]$ denote the set of all polynomials in the form
$a_0 + a_1T + \cdots + a_d T^d$,
where with coefficients $a_0$, $a_1\dots,a_d$ in~$\mathbb{R}$.
Define an equivalence relation on $\mathbb{R}[T]$ by declaring $f(T)\sim g(T)$ if $f(T)-g(T)$ is divisible by $T^2+1$. Each equivalence class has a unique representative of the form $a+bT$. 
For example, the imaginary unit is represented by the equivalence class
$$
\{T + p(T)(T^2+1): \, p(T) \in R[T]\},
$$
a real number $a$ is represented by 
$$
\{a + p(T)(T^2+1): \, p(T) \in R[T]\}.
$$

In each equivalence class, we select a representative in the form $a+bT$, for some real numbers $a$ and~$b$.  Addition and subtraction is performed modulo $T^2+1$. The quotient ring is denoted
$\mathbb{R}[T]/(T^2+1)$.
This quotient ring has the structure of a field, and is a model of the complex field.

\bigskip

All three constructions are isomorphic. The correspondences are:
\[
(a,b)
\longleftrightarrow
\begin{bmatrix} a & -b \\ b & a \end{bmatrix}
\longleftrightarrow
a + bT.
\]

The real number 1 can be represented by
\[
(1,0),\quad
\begin{bmatrix} 1 & 0 \\ 0 & 1 \end{bmatrix},\quad
1.
\]
The vector $(1,0)$ corresponds ot the unit vector in the $x$-axis, the matrix $\begin{bmatrix} 1 & 0 \\ 0 & 1 \end{bmatrix}$ is the identity matrix, and ``1'' is the polynomial consisting of the constant term~1.


The imaginary unit corresponds to
\[
(0,1),\quad
\begin{bmatrix} 0 & -1 \\ 1 & 0 \end{bmatrix},\quad
T.
\]

Usually, we adopt the notation
\[
a+bi
\]
for a complex number. The elements $1$ and $i$ form a basis over $\mathbb{R}$, so $\mathbb{C}$ is a $2$‑dimensional real vector space. Multiplication is given by
\[
(a+bi)(c+di) = ac - bd + i(ad+bc).
\]

Real numbers embed into $\mathbb{C}$ via $r\mapsto r+0i$, so $\mathbb{C}$ is indeed an extension of $\mathbb{R}$.



\begin{svgraybox}
	\begin{theorem}
		Let
		\[
		\mathbb{C} \triangleq \{a+bi : a,b\in\mathbb{R}\},
		\]
		where $i$ is a symbol satisfying $i^2=-1$. Then $\mathbb{C}$ is a field and, simultaneously, a real vector space of dimension $2$.
	\end{theorem}
\end{svgraybox}

\begin{remark}
In the vector representation of complex numbers, the element $(0,1)$ is chosen to represent the imaginary unit $i$. However, this choice is not unique. We could instead take $(0,-1)$ to be the imaginary unit. If we do so, all definitions and operations can be reformulated consistently, and the resulting number system is isomorphic to the usual complex numbers. In other words, choosing $i$ or $-i$ is an arbitrary convention in the construction.
\end{remark}

\begin{greenbox}
	\begin{definition} \index{Real part} \index{Imaginary part} \index{Imaginary unit} \label{def:real_imaginary_part}
		For a complex number $z = a + bi$ in $\mathbb{C}$, the \df{real part} and \df{imaginary part} of $z$ are defined by
		\[
		\real(z) \triangleq a, \qquad \imag(z) \triangleq b.
		\]
A complex number $z$ is said to be \df{purely imaginary} if $\real(z)=0$, and \df{real} if $\imag(z)=0$. 
	\end{definition}
\end{greenbox}



\section{Complex conjugate, and modulus}



\begin{greenbox}
	\begin{definition} \index{Complex conjugate}
		The \df{complex conjugate} of a complex number $z = a + bi$ is defined as
		\[
		\bar{z} \triangleq z^* \triangleq a - bi.
		\]
	\end{definition}
\end{greenbox}

Geometrically, the real and imaginary parts of $z$ are simply the $x$- and $y$-coordinates of the point representing $z$ in the complex plane.  
The complex conjugate $\bar{z}$ is the reflection of $z$ across the real axis (See Fig.~\ref{fig:complex_conjugation}).  



\begin{figure}
	\begin{center}
\begin{tikzpicture}[>=Latex, scale=1.2]
	% Axes
	\draw[->] (-0.5,0) -- (4.0,0) node[below right] {$\real$};
	\draw[->] (0,-2.0) -- (0,2.0) node[above left] {$\imag$};
	
	% Origin
	\coordinate (O) at (0,0);
	
	% Point z = a+bi and its conjugate
	\coordinate (Z) at (3,1.2);   % z
	\coordinate (Zbar) at (3,-1.2); % \bar z
	
	% Vector to z
	\draw[->, thick, blue] (O) -- (Z) node[above right] {$z = a+bi$};
%	\fill[blue] (Z) circle (1.4pt);
	
	% Vector to \bar z
	\draw[->, thick, red] (O) -- (Zbar) node[below right] {$\bar z = a-bi$};
%	\fill[red] (Zbar) circle (1.4pt);
	
	% Vertical dashed line showing reflection
	\draw[dashed] (Z) -- (Zbar);
	
	% Projections on real axis
	\draw[dashed] (Z) -- (3,0) node[below right] {$a$};
	\draw[dashed] (Zbar) -- (3,0);
	
	% Label real axis tick at a (optional)
	\draw (3,0.08) -- (3,-0.08);
\end{tikzpicture}
	\end{center}
\caption{Complex conjugation}
\label{fig:complex_conjugation}
\end{figure}

We list some elementary properties of complex conjugation in the next proposition. The proof is straightforward. 


\renewcommand{\labelenumi}{(\alph{enumi}).}
\begin{greenbox}
	\begin{proposition} \label{prop:conjugate} 
		For complex numbers $z\in \mathbb{C}$:
		\begin{enumerate}
			\item $(z^*)^* = z$. 
		
			\item $z^* = z$ if and only if $z$ is real. 	
			\item $z^* = -z$ if and only if $z$ is purely imaginary.

		\end{enumerate}
	\end{proposition}
\end{greenbox}

\begin{comment}
Part (i)
import Mathlib.Tactic

open Complex

-- Define the usual notation for complex conjugate
notation "conj" => (starRingEnd ℂ)

/-
conj conj z = z
-/
lemma conj_conj (z : ℂ) : conj (conj z) = z := by
-- Use extensionality on real and imaginary parts
apply Complex.ext
· -- real parts are equal
simp 
· -- imaginary parts are equal
simp 

-- Prove using existing theorem in Mathlib
-- `star_star z` is a general theorem about involution in
-- a semi-ring
@[simp] lemma conj_conj' (z : ℂ) : conj (conj z) = z :=
by exact star_star z

	
	
part (ii)
import Mathlib.Tactic

open Complex

-- Define the usual notation for complex conjugate
notation "conj" => (starRingEnd ℂ)

/-
conj z = z if and only if imag(z) = 0
-/

-- detailed proof 
example (z : ℂ) : conj z = z ↔ z.im = 0 := by
constructor
· -- conj z = z → z.im = 0
intro h
-- Two complex numbers are equal iff their components are equal
rw [Complex.ext_iff] at h 
-- Break h into real (h_re) and imaginary (h_im) parts
obtain ⟨h_re, h_im⟩ := h
-- Simplify (conj z).im to -z.im
simp at h_im 
-- We now have -z.im = z.im. Arithmetic solves this.
linarith 
· -- z.im = 0 → conj z = z
intro h
-- To prove two complex numbers are equal, check components
-- and use the hypothesis that Im(z)=0
simp [Complex.ext_iff, h]

-- prove using existing theorem in Mathlib
@[simp] lemma conj_eq_iff_im' (z : ℂ) : conj z = z ↔ z.im = 0 := by
exact conj_eq_iff_im
\end{comment}


\begin{proof}
Part (i) says that the conjugate of a conjugate is the same as itself \iflean \ 
\href{https://live.lean-lang.org/#codez=JYWwDg9gTgLgBAWQIYwBYBtgCMB0AVJAYxmEICgyIwBTAOzgGEJx1qAPCgWk7gBFqAZsFrU4aUQFcAzhKTo4tCDBTAI9AdDiFmYVmy1qAVhIDmKamUXKSauACJttQ3bgBeAHxwAFFOVQASsImAKK0ACZwgECEAJQUAPScZHDJBk6phnAAXm5ZZJxxZKwgIEjpAPqOGV7ZAFxR0XB1ld7NmQ2uWY0dWACeSXDccACqUqLsMHRSqrRywDA9cLZQ1HJwSOFwoEgmwkhQC2B7MFL9SGC6C0ws7Djj/QDtAzzLq4ewUmvLcNQAjrLo/WSU3AcAeT02JR2M32cDex0+Yz+ckBcGBYFBeR4AAUoBAAG6SKa0EzfNjAXxBMSoajQaggTb0ZBoTBYTFwAAGvj2ZS5UCy7M2H1KJjo1Cgq3EtPpSCwEAk8GEeIg6Hl0wZbNKoxAwE4UCChTpJXKlQA5N5avVGukWkYsu1OjVXP1eqSiPBeTy/FkgA}{(\LEAN proof for part (i))}\fi. \index{Involution} Parts (ii) and (iii) give the relation between complex conjugate and the real and imaginary part of a complex number. \iflean
\href{https://live.lean-lang.org/#codez=JYWwDg9gTgLgBAWQIYwBYBtgCMB0AVJAYxmEICgyIwBTAOzgGEJx1qAPCgWk7gBFqAZsFrU4aUQFcAzhKTo4tCDBTAI9AdDiFmYVmy1qAVhIDmKamUXKSauACJttQ3bgBeAHxwAFFOVQASsImAKK0ACZwgECEAJQUAPScZHDJBk5wAF5uGXDAAnBI4XBq6ACeOXmgSCZe6dFZAAxknHFcPGHUysCsEWBQEBB5ZOxILKI1cABcUXVTjobZrtmAKYQZOKANk4tYJUmpvlASxNC7AO1w3KnzmYuZgEmEq+uLjSk5tDB9cKi7yRd4AO4QAyjfS0CQgLDUKBSfJQUTUACOsnkuTy4mAUCBkBEb2hSFhcARSO+cCgfzgAG0mMCcOwYAB9FEAXXy8FQcGJFwAQrCkABrT6vGCAnnyLyoOmwuoFCKVEzCPFlMUMkB1MB4mBSYkQLCdeiAC/JxbCADSfZWAS/JNp8OTwAMqgXS5RVzDLRNYgMSAzjpN3EqT2lmm9bWuAAdVEijJqCQADdRF63VlvaAcHAAIJQYBoEAdUhwKQQdCx6FoYBSHDEzC0PGZtmnc48JPup5we7O64ZYnCd6Ar4vX6A3oQWNiAGYvQKMEQqEwuGIuQmwioaiEfnacBqOga4PSuDSUTiT4lSDiP3FqPwACSIBq0VczxSfvAFKpunYNLY9JRJtQjPiiRe6QAFRZF65T5IUxRlLkJLUHIN4NE0LTDKM3iZFMMSTJcCznJkKzevizYTFsOzJI4+yHEKUB1hcbYgXcqwEXA97JF2Hy9ikvyjmur4gpOkK4vihJyOUqJLuimIbjiM4EnO6DEqSz46Ho76fgIAjMigArBtysH8myXbCrBoqGtQUqFLK8pQIq4qgKq6qai82q6nABoStQ37mpa7E/La9qYAITpGC6OD4kKOEhRYLyPmAAYmcGYYKBAkYxqI+GiIs8b4hWlk1uyyRnBcaUbK2QXtl6nZvGxwZ4AOfTDjAXFKewE7gvx0lCegJp/KIi7LvyB7cZJW59jwO57mIS6HseS6nhNmkijed6+v6lJNWwKkMmp36/q0cCDsO0hBASbCliQtAmBN1DQNQ7rCIgKAYNgZAAALktFzKsCAIBIJcdIIptAjKgA5KhmEYbMpWJnAeEJoRxG7MMxC/f9KLKkAA}{(\LEAN proof for part (ii) and (iii))}
\fi
\end{proof}


\begin{svgraybox}
	\begin{theorem}\  \label{thm:complex_automorphism} 
\begin{enumerate}
	\item $(z_1 + z_2)^* = z_1^* + z_2^*$, for $z_1, z_2\in\mathbb{C}$,
	\item $(z_1 z_2)^* = z_1^* z_2^*$,  for $z_1, z_2\in\mathbb{C}$,
	\item $(1/z)^* = 1/(z^*)$ for $z\neq 0$.
\end{enumerate}
\end{theorem}
\end{svgraybox}

This theorem says that the conjugate of a sum is equal to the sum of the conjugates, and the conjugate of a product is equal to the product of the conjugates.

\begin{proof}
Suppose $z_1=a+bi$ and $z_2=c+di$. For the conjugate of sum, we calculate
\begin{align*}
z_1^*+z_2^* &= (a+bi)^* + (c+di)^* \\
&= a-bi + (c-di) \\
&= (a+c) - i(b+d) \\
&= (z_1+z_2)^*.
\end{align*}	
Likewise, we prove the second part by
\begin{align*}
z_1^* z_2^* &= (a-bi)(c-di) \\
&= ac-bd - i(ac+bd) \\
&= ((ac-bd) +i(ac+bd) )^* \\
&= (z_1z_2)^*.
\end{align*}

For the last part, we can apply part (ii) and check that for $z\neq 0$,
$$
(1/z)^* z^* = ((1/z)z)^* = 1^* = 1.
$$
Since the product of $(1/z)^*$ and $z^*$ is equal to 1, the inverse of $z^*$ is equal to $(1/z)^*$.
\iflean \href{https://live.lean-lang.org/#codez=JYWwDg9gTgLgBAWQIYwBYBtgCMB0AVJAYxmEICgyIwBTAOzgGEJx1qAPCgWk7gBFqAZsFrU4aUQFcAzhKTo4tCDBTAI9AdDiFmYVmy1qAVhIDmKamUXKSauACJttQ3bgBeAHxwAFFOVQASsImAKK0ACZwgECEAJQUZAD0nGR4qNTQ1CBwAIw4AMx5ZHDeAF4A+llwANRwZQBM0QB6AFRuNeXNVW21HUjhJeVdjS2ubVkddc1knPEU4umZjoalSGFhpYslgIEENYBBBHAAXFHRBwZONdvVG8V7IxtexRe7x/sjWACehXDccGGCwqJIAA0cCwwMIwIin1Y8CQBxGDxwUGoAG44NCQXDzjhQFDqPBCJjroiUWi8T9CTtsSBPp9CHJyEUit9gmwwL0Ig84OzdnBhDAIAZwGo6DApIcvLDqlhgENOl4CdUwjLJozToZzp0rntPkyeABBMC6N5q0zmH5/WjAGy0cWwnjS47VeVfH4ynVwUqtJgsdg4EAAay53k4WEdjB0ej9gYJXk4YWer2NUAE6Hd3wA4lAIBIwHAkXIuX1QEgTMIkFBjWzYGLvJKtMceMAvFhOvH3Z6Rt7dL6A7XNcdY83W9EEyCPqrGVAgqVaAJ3ZOU2meAAhY2/ISW63Aq1wagAR1k6Ckaq2lWusVVHZP90e58x7zzi9mqXmatKIAk6HWRi2PMOMROK5thaLVWjuTkWjvF4x0+b513+LlAVBcEwlxGFCU2YlUXRFtoIRHEinRAk8MpJFsLJCISKpGkijpdAGUZZlWW5TkWJ2d0gLgECf2uOAlzgA0jRNMwYFEeDN1UWh2y9CMe0DWFY1DLjwx9Ngoy0YN43vJNF1Vb4AAVqGTaAFlk/QP3QEhdFIFQ1H4gB1UQAHdengfk82oIhUFrAl7S0xsJQiKVCAvRieAAVSkUQWVE2gpEkuQrWNdypGs+A+QFfxPPkbkAEkQBLMsKzgEwIDkKRpM7Mz1IlZTfIxFotMHWFGs6FsQJHbS+InVUkENdBjS7SN2HgedVQAdrgeLwDGxlJumsBZtVKdaBMbqwrgFJRCRXxeWPGAoAkUQH3Eq1JMq68IKeIpoIfZNUwoRIdSuQBvAgATtaeUeOid6phmdgkB9EoTgArxUGKE4IcAAyI4AABmeE9inez6rh+j7bvHKa3hAakij6oThAAN1KfdSggAR30/Em9zJkRSinExUBgd1YJ4AA5CAnLgJzRBoYyoEycQtHpT9bJtGomi8LJ4mKY4RhwBXWiyWl6Q4ni6p45H1t1FTpREYTzBrKWZeiSDhm8Y3Zdlpp+I50TDgKsBKfkYBjy8NgmjeU3Wg95S3iaYEpAFHm4GkURAATCMQBVKkEiH9FyoDCCrLzAn97iaJG3tHO6uYAbUjx3nYAXRZvSeAAZVAayBGNYpLeOdzlZT24fwqHqbsTPM86GnsqaJ9ZekIah5HBkv+KYJxTVEuByeyPbsgutv24xru4FzwvhTHqYeH0rNCeO418YGoJdzYV2SFWsQXyREBj2ERAUAwbAyABoGbz/I5AJ4x5QJbs53+qFBRMnwAbEDgMEPcOApDY0yF4QIq0AASzA/R9WWKsbwvhyzwJCH0ACnJzwUFft2X8vF/wI04txM4vE/7qnfpBHY2djQgLYEQeAECoEwO8NgpBIAUFOwshgvw2DQgRDwdsAhL8WFvwhmQr+VDkYjC+lQtG2lmGsJUt2NSixShExqEAA}{(\LEAN proof)}\fi
\end{proof}

\begin{comment}
import Mathlib.Tactic

open Complex

-- Define the usual notation for complex conjugate
notation "conj" => (starRingEnd ℂ)


/-
Theorem 1.3.3
(z_1 + z_2)^* = z_1^* + z_2^* and (z_1 z_2)^* = z_1^* z_2^*
-/

theorem conj_add_conj (z₁ z₂ : ℂ) : conj z₁ + conj z₂ = conj (z₁ + z₂) := by
-- define a, b, c, d 
let a := z₁.re; let b := z₁.im
let c := z₂.re; let d := z₂.im

calc
-- Expand z₁ and z₂ into components: (a + bi)^* + (c + di)^*
conj z₁ + conj z₂ 
-- Apply conjugate definition: (a - bi) + (c - di)
_ = Complex.mk a (-b) + Complex.mk c (-d) := by rfl
-- Group real and imaginary parts: (a + c) - i(b + d)
_ = Complex.mk (a + c) (-(b + d)) := by
ring_nf
rfl
-- By definition, it equals conj (z₁+z₂)
_ = conj (z₁ + z₂) := by rfl


theorem conj_mul_conj (z₁ z₂ : ℂ) : conj z₁ * conj z₂ = conj (z₁ * z₂) := by
-- define a,b,c,d
let a := z₁.re; let b := z₁.im
let c := z₂.re; let d := z₂.im

calc
-- Expand z₁ and z₂
conj z₁ * conj z₂ 
-- Apply conjugate definition
_ = Complex.mk a (-b) * Complex.mk c (-d) := by rfl
-- Perform complex multiplication
-- We want to reach (ac - bd) - i(ad + bc)
-- Use Extensionality to split into Real and Imaginary goals
_ = Complex.mk (a * c - b * d) (-(a * d + b * c)) := by 
apply Complex.ext 
· simp
· simp
ring 
-- The rest is true by definition
_ = conj (z₁ * z₂)  := by rfl

example (z : ℂ) (hz : z ≠ 0) : conj z⁻¹ = (conj z)⁻¹ := by
symm
apply inv_eq_of_mul_eq_one_right

-- Now we perform the calculation: z*(1/z) = ... = 1
calc
conj z * conj z⁻¹  
-- Combine conjugates: (1/z)* z* = ((1/z)z)*
-- Note: map_mul is (x*y)* = x* * y*, so we use ← to go backwards
_ = conj (z*z⁻¹) := by rw [← map_mul]

-- Simplify z(1/z) to 1
_ = conj 1         := by rw [Complex.mul_inv_cancel hz]

-- Conjugate of 1 is 1
_ = 1         := by rw [map_one]


-- Prove by applying an existing theorem in Mathlib
example (z₁ z₂ : ℂ) : conj z₁ + conj z₂ = conj (z₁ + z₂) := by
exact Eq.symm (RingHom.map_add (starRingEnd ℂ) z₁ z₂)

example (z₁ z₂ : ℂ) : conj z₁ * conj z₂ = conj (z₁ * z₂) := by 
exact Eq.symm (RingHom.map_mul (starRingEnd ℂ) z₁ z₂)

example (z : ℂ) : conj z⁻¹ = (conj z)⁻¹ := by
exact Complex.conj_inv z
\end{comment}

By Theorem~\ref{thm:complex_automorphism}, we see that complex conjugation is an {\em automorphism} of the complex field. \index{Automorphism} It fixes all the real numbers and maps $i$ to $-i$, preserving the algebraic structure. (cf.\ the remark before Definition~\ref{def:real_imaginary_part})

\smallskip

We now illustrate how complex conjugation is used to perform division. Suppose we want to 
find the complex number $z$ satisfying $(c+di)z = a+bi$, assuming $c+di \neq 0$. The solution is the quotient $(a+bi)/(c+di)$. Multiply numerator and denominator by the conjugate of the denominator:
\begin{equation}
	\frac{a+bi}{c+di} \cdot \frac{c-di}{c-di}
	= \frac{ac+bd}{c^2+d^2} + i\,\frac{bc-ad}{c^2+d^2}.
	\label{eq:complex_division}
\end{equation}

Complex division can also be performed using the $2\times 2$ matrix representation. Division corresponds to matrix inversion:
\[
\begin{bmatrix} a & -b \\ b & a \end{bmatrix}
\begin{bmatrix} c & -d \\ d & c \end{bmatrix}^{-1}
=
\frac{1}{c^2+d^2}
\begin{bmatrix} ac+bd & ad-bc \\ bc-ad & ac+bd \end{bmatrix}.
\]

\begin{remark}
	The quantity $c^2 + d^2$ is the determinant of the matrix
	$\begin{bmatrix} c & -d \\ d & c \end{bmatrix}$.  
	In the matrix model, complex conjugation corresponds to taking the transpose.
\end{remark}

\begin{example}
	Solve the equation
	\[
	3z + iz + 2 = 1 - i.
	\]
	
	First isolate $z$:
	\begin{align*}
		3z + iz + 2 &= 1 - i, \\
		3z + iz &= -1 - i, \\
		z &= \frac{-1 - i}{3 + i}.
	\end{align*}
	
	Now compute the inverse of $3+i$:
	\begin{align*}
		z &= \frac{-1 - i}{3 + i} \cdot \frac{3 - i}{3 - i} \\
		&= \frac{(-1 - i)(3 - i)}{10} \\
		&= -0.4 - 0.2i.
	\end{align*}
\end{example}


\begin{figure}
	\begin{center}
\begin{tikzpicture}[>=Latex, scale=1.2]
	% Axes
	\draw[->] (-0.5,0) -- (4.0,0) node[below right] {$\real$};
	\draw[->] (0,-0.5) -- (0,2.6) node[above left] {$\imag$};
	
	% Origin
	\coordinate (O) at (0,0);
	
	% Point z = a+bi
	\coordinate (Z) at (3,2);  % choose some coordinates for illustration
	
	% Vector representing z and its modulus
	\draw[->, thick, blue] (O) -- (Z) node[above right] {$z = a+bi$};
%	\fill[blue] (Z) circle (1.4pt);
	
	% Projections to axes
	\draw[dashed] (Z) -- (3,0) node[below] {$a$};
	\draw[dashed] (Z) -- (0,2) node[left] {$b$};
	
	% Right angle marker
%	\draw (0.3,0) -- (0.3,0.3) -- (0,0.3);
	
	% Label modulus
	\node[blue] at (1.8,0.3) {$|z| = \sqrt{a^2+b^2}$};
\end{tikzpicture}
\end{center}
\caption{Geometric meaning of the modulus of a complex number.}
\label{fig:complex_modulus}
\end{figure}

The product of $z$ and its conjugate is real.

\begin{svgraybox}
	\begin{theorem} \label{thm:norm_squared}
	If $z=a+bi$ for $a,b\in\mathbb{R}$, then 
\begin{enumerate}	
\item	$z z^* = a^2+b^2$. 
\item 	 
\begin{equation}
	a = \frac{z + z^*}{2}, \qquad
	b = \frac{z - z^*}{2i}.
	\label{eq:real_imaginary_conjugate}
\end{equation}
\end{enumerate}
	\end{theorem}
\end{svgraybox}

\begin{proof}
We prove part (i) by calculation,
$$zz^* = (a+bi)(a-bi) = a^2 - (bi)^2 = a^2 - b^2(i^2) = a^2+b^2.
$$
The proof of Part (ii) is straightforward. \iflean \href{https://live.lean-lang.org/#codez=JYWwDg9gTgLgBAWQIYwBYBtgCMB0AVJAYxmEICgyIwBTAOzgGEJx1qAPCgWk7gBFqAZsFrU4aUQFcAzhKTo4tCDBTAI9AdDiFmYVmy1qAVhIDmKamUXKSauACJttQ3bgBeAHxwAFFOVQASsImAKK0ACZwgECEAJQUZAD0nGR4qNTQ1CBwAIw4AMw4ACwANHAACkiw3sDRJQBecLUAegBUbnBIjQBMcADUcFhdZJzxXDxh1MrArBFgUBAQAmTsSCyiXvUAXFHRcBtkcAcNcK2Ohkeu3rU4UNTRXb2XOKB33RsXWACe+3Cs8Ei7FyuNwA3D8Jv0AQ0niBvt9CHJyIcjicjEdvgduHAAIJgXQfAxOUzmLZef59LDVUlwHgU2JIgD6bS8TFWbBwIAA1u1+jtWsydHp2Vz/l5OFhojsDm9+vioAJ0Oi4IrMfwhLRgDZ6As4CAJOgSLpSCo1CSkP8aaLxTs+sBSZbrf0kHTDoyLizdOwhd5/q1zRC+WKJd7jt5Aw8sCGnZDPnA5QqkcqeABlUCGgT4rHoEzULBQJBbDrdcn3G1eAAMzoOrsYAs9nODjTgRYhjc6OzLUve+MVB0cJigPbgAHbY0FqTx/NQ5HAwBUYAXmn7LV0dhdC+HBkiDiOoGPMQBJEBIEzCCr42ewLacJARvo3tplihIzFMWgAN2olXdgvr+g7MAgOBjF8WMp3kWgJBALBPzgfQvHhXxqAiACa1ZSs4GrKlW3DOBsK2GJDmlGVBwOJBcXQfFv09dgYBI4c4CkVMlS3LcoAAdzgABtQAEwlQj02QWSc5HpSAOP+ToSl4qiBIEIT0BEiAOIjToAF06IOG4aBQWMOM46ScEEsD6RuFTmJY+jGPAMzzJ0ripNrGS5IUsSm0kvjBUM4TRIhVT1NArT4HYrjPPk0A1JsjSglGGc5m1aQxwAA11eTTgSpY2BWD1Ll2bZdkVeoUScc5HhuF4HnWaEyulT5viCvSHPZPV6VOcK4AkWgNHQCJFCgEAkwAR2+SywFqqKKESZJUnSTIcnyYoyjnKpgCjIFRAudYHlOBodniJslWGaLxkmaYYvmRZllWbL8J2LZVqZeo+i22odr26NuwOX5uWlVbvk+iNvuhWEDikD4QBhXsEUVDbHtRZ64F27pEzgYI2FncIjhQ/T63+LBFUwrHhQhGGiv5VkvRxoMEbe2N5UVJGcTxOBxjVDVVC1AQCWMMwYAsBkmQJ7lb3cutCftF7uiRarZVphNnx4UpPw0XqxFSdowjCVmTW9Hoo1LLBAzx/mGvrKk7x2LwhcDcWIslrsafjQ4kZTFhgHTFXRAgkBPxQaAtk6Mk4GAR8+fWgWvG6X122tzsZWsw4+ygPbMQAMQgQhpED2hGPGd24AAck5PPBx3PcJzAmc5wLB5/guCP2mL0daBMcc4EPY9Tygc9K6J6kIwuYOtyR18P0qAuuXDqMyzzsRAPELQkEQiIbmnT3oMTvOJ5ymIi5D7w6/+a74demOYzosjGf0mi6JHYbr4Y1NDbXanb6GHhZggD9Y8cbPPzHZf5CQOjUA7daBngrrAKQ7R4Bz2gieWg6om4sXSplVgV1cq3WuGtbKxMzhwyplLb458KLCzZFfb4mIADiEBpxZC2HJcBMBIFHhgIQVA5CeAABkAASSYME3HYXAfwPCSTQ05ttI+bZMHfBvg/ARVDpydC2G3OBYCLyMJ1CgVhAjuG8LgB2HwwhCCiDusASB/9nSYiEborwoinri2iIDbc99wBxAmikNINwZp5EKCUcolQvDAGwDsK4oB7rUjEXgve1R9ojAullDah8MGhPWvUHgdij7h2aPuSUUszIgzBgIhAeoDTEKwEoVA99xiQJjJ0LJM92hvggMACIAg8zEDZnAVISAwhEFSFIWqukNZvnpNQfq9JXYCFapiAAqlIUQKMeZZzZnIDU+IUKsOoIQLkcl4jKI7l3CB0j2jkUog1K+SIZFWUseXNRTDNEVPLA+dClyRosV3E3EiB4jwqM7gwu5LCHmdD7k2cUT4W6lDmJ/cYigQCngAonRQtBOC1E/BAI5w0uL7npCIekKK5hqVfqdT+8VEGAIiOwUxJBEHiGmpnRAKAMDYGQZdBJ6CoQ3DCTg8RVM3rfGWMQUCIyxndLCM1WG0V37EsYqS9GFLfBjhpZ4ulyA0CYFxnE1BrKCJJMyCk8J6TdqZOySfL4Bx+XwFAEK+kMgsBiqKrUIAA}{(\LEAN proof)} \fi
\end{proof}

\begin{comment}
import Mathlib.Tactic

open Complex

-- Define the usual notation for complex conjugate
notation "conj" => (starRingEnd ℂ)


/-
Theorem 1.3.4, Part (i), z z^* = a^2 + b^2
-/

-- detailed proof
example (z : ℂ) :
z * conj z = (z.re)^2 + (z.im)^2 := by
let a := z.re; let b := z.im

calc
z * conj z 
-- Apply conjugate: (a + bi)(a - bi)
_ = (Complex.mk a b) * (Complex.mk a (-b))   := by rfl

-- Definition of multiplication: (aa - b(-b)) + i(a(-b) + ba)
_ = Complex.mk (a * a - b * (-b)) (a * (-b) + b * a) := by rfl

-- Simplify Algebra: a^2 + b^2 + i(0)
_ = Complex.mk (a ^ 2 + b ^ 2) 0  := by 
congr
· ring -- Real part: a*a - (-b^2) = a^2 + b^2
· ring -- Imaginary part: -ab + ab = 0

-- Convert Complex.mk x 0 to just real number x (casted to Complex)
_ = (a ^ 2 + b ^ 2 : ℂ)   :=  by
apply Complex.ext
· simp 
rw [← Complex.ofReal_pow a 2, ← Complex.ofReal_pow b 2]
repeat rw [Complex.ofReal_re] 
· simp 
rw [← Complex.ofReal_pow a 2, ← Complex.ofReal_pow b 2]
repeat rw [ofReal_im]
ring

-- proof using `mul_conj`
example (z : ℂ) :
z * conj z = (z.re)^2 + (z.im)^2 := by
rw [Complex.mul_conj]
unfold normSq
simp
ring


/-
Theorem 1.3.4, Part (iia) z.re = (z + conj z) / 2 
-/

-- detailed proof
example (z : ℂ) : z.re = (z + conj z) / 2  := by 
let a := z.re
let b := z.im

symm
calc
(z + conj z) / 2 
-- Expand z to Complex.mk a b
_ = (Complex.mk a b + conj (Complex.mk a b)) / 2 := by rfl

-- Apply definition of conjugate
_ = (Complex.mk a b + Complex.mk a (-b)) / 2     := by rfl

-- Perform the addition: (a+a) + i(b-b)
_ = (Complex.mk (a + a) (b + -b)) / 2             := by rfl

-- Simplify the numerator: 2a + i0
_ = (Complex.mk (2 * a) 0) / 2   := by 
congr 2 -- Focus inside the 'mk'
· ring -- Real part: a + a = 2 * a
· ring -- Imaginary part: b + -b = 0

-- Convert 'mk (2a) 0' to the casted real number '(2a : ℂ)'
_ = (2 * a : ℂ) / 2    := by 
apply Complex.ext
· simp
· simp
_ = a := by simp

-- prove by considering real and imaginary parts at the beginning      
example (z : ℂ) : z.re = (z + conj z) / 2 := by
apply Complex.ext

-- Goal 1: Real parts match
-- LHS: z.re
-- RHS: ((z + conj z) / 2).re
· simp

-- Goal 2: Imaginary parts match
-- LHS: 0 (since z.re is real)
-- RHS: ((z + conj z) / 2).im
· simp


/-
Theorem 1.3.4, Part (iib) z.im = (z - conj z) / (2i) 
-/
example (z : ℂ) : z.im = (z - conj z) / (2*I)  := by 
symm
-- Multiply both sides by 2*I to avoid fraction headaches
rw [div_eq_iff]
-- Use Extensionality to check Real/Imaginary parts
· apply Complex.ext
· simp -- Real parts match (0 = 0)
· simp
ring
-- Imaginary parts match (2b = 2b)

-- Prove denominator non-zero
· simp [I_ne_zero]

-- prove using and existing theorem in Mathlib
example (z : ℂ) : z.re = (z + conj z) / 2  := by
exact re_eq_add_conj z

-- prove using and existing theorem in Mathlib
example (z : ℂ) : z.im = (z - conj z) / (2*I)  := by
exact im_eq_sub_conj z
\end{comment}

\begin{svgraybox}
	\begin{definition}
		The \df{modulus} of $z = a + bi$ is
		\[
		|z| = \sqrt{a^2 + b^2}.
		\]
		It is also called the \df{absolute value}, \df{magnitude}, or \df{radius}.
	\end{definition}
\end{svgraybox}

The geometric meaning of the modulus is depicted in Fig.~\ref{fig:complex_modulus}.


\begin{svgraybox}
	\begin{proposition}\ 
		
		\begin{enumerate}
			\item $|z|^2 = z z^*$ for all $z \in \mathbb{C}$.
			\item For all $z_1, z_2 \in \mathbb{C}$,
			\[
			|z_1 z_2| = |z_1|\,|z_2|.
			\]
		\end{enumerate}
	\end{proposition}
\end{svgraybox}

\begin{proof} Part (i) is a re-statement of part (i) of Theorem~\ref{thm:norm_squared}.
	Using part (i) and commutativity of multiplication, we can obtain part (ii) by
	\[
	|z_1 z_2|^2
	= (z_1 z_2)(z_1 z_2)^*
	= z_1 z_2 z_1^* z_2^*
	= (z_1 z_1^*)(z_2 z_2^*)
	= |z_1|^2 |z_2|^2.
	\]
\iflean \href{https://live.lean-lang.org/#codez=JYWwDg9gTgLgBAWQIYwBYBtgCMB0AVJAYxmEICgyIwBTAOzgGEJx1qAPCgWk7gBFqAZsFrU4aUQFcAzhKTo4tCDBTAI9AdDiFmYVmy1qAVhIDmKamUXKSauACJttQ3bgBeAHxwAFFOVQASsImAKK0ACZwgECEAJQUZAD0nGRwcHio1NDUIHAAjDgAzDgAbAA0cGBIsN7A0WUAXnB1AHoAVG5wAD51HU0ATGSc8WTsSCyiXg0AXFHRcNMNbY6Gje2AaAR1q3BNcL1zrnBYAJ7JcFAA7nAA2kxjbDggEugA+ksAuincK15LjbP7ilAQABlACOjRO5yugATCRg6PQ4CACfzUORPSBnV4nAEgF5IXwnEbEBTQYEgp7UMnYp5SMENT5QahhCSEUQwCBwKSoCBnabEwGglb7dabbb9CiJMhpDIM7J5QqlcqVeBeYA1MhdQCBBI1AEEEHXamr1XV1A3iKQGPDC1GUwFYETAUAgiOGbFGunGdS1dW1cxmPvWWraXs2aw9mza621m0m+yOJ0+QOELL5ICkcEqokUtE4IjMJAAbtQyvp9h02B0cHA4zwAOqiBlnKDAGCs9JwEwQORwACC7QAQmJ2R0u3qS729VJ2eI4KwQCAkHA5zBCOkpDgIRdLjCkFgpE9EU9MznvFTD9QTHAntFMSlIZu0zu9wID2oj14Hs9T+evCeX2eL7NvxJZ9aCPS8rwoD4eC7MBdEOMRWy7AAFABJaYhxHToxzgQAUwm7Po+z6ddoQ5MkKWpMlgAEJ9t13MiaMxKt7Dwdl7QgAtThReRizgQ4yibDkJCo0hqFTNlygddjtFuOBAETCbiZMOOwiLvRFkVRYRDHeCC4HjM4m2XAdYWk2gJBALBqCgUT2WkUQllMcxxKoCySBEk4wGkVBcXxRiADFhGATl4NEVjEQOODCDkQhHhUNQpBOCL0HIFIUi8f04EDSMfRiLYdkrZKL3aCYAx1WZFiMbwPXSkr8pSaMwtODcYXfF4jAAcgqjUWi9cD8qeQrKoy0rvB+AaDCcaq6qOBqrjnMAnnfa9kr6/YiqqkaNSGiZvTK8bur2erG1oEwThSZbvH9KMZnwtpUqDLLonw/LJrg29mqWdrKrKN62p1RbTsKi6qojS7stFGrwee6bLma9FXgAbkrc1HPYqbpCCOB2ACkgjrTMBgDgYREBQDBsGdV1WA6nV7r9Ua7pDDUwzgYH9qmwl4BuN07ipd9Gk9bUgA}{(\LEAN proof)}\fi
\end{proof}


\begin{comment}
	
import Mathlib.Tactic

open Complex

-- Define the usual notation for complex conjugate
notation "conj" => (starRingEnd ℂ)


/-
Theorem 1.3.6, part (i), z z^* = |z|^2
-/
example (z : ℂ) : z * conj z = ‖z‖ ^ 2 := by
rw [Complex.mul_conj]  -- z (conj z) = normSq z
rw [← Complex.ofReal_pow]
norm_cast
exact normSq_eq_norm_sq z -- reduce to show:  normSq z = ‖z‖ ^ 2


/-
Theorem 1.3.6, part (ii)
|z₁ z₂| = |z₁| |z₂|
-/  
-- detailed proof
example (z₁ z₂ : ℂ) : ‖z₁ * z₂‖  = ‖z₁‖ * ‖z₂‖ := by
-- Since norms are non-negative, x = |x|. 
-- We rewrite the goal A = B to |A| = |B| so the lemma matches.
rw [← abs_of_nonneg (norm_nonneg _)]
rw [← abs_of_nonneg (mul_nonneg (norm_nonneg _) (norm_nonneg _))]

-- Apply the API: |A| = |B| ↔ A^2 = B^2
rw [← sq_eq_sq_iff_abs_eq_abs]

-- "To prove real x = y, it suffices to prove complex ↑x = ↑y"
rw [← ofReal_inj] 

-- Switch to Complex numbers to use conjugate properties
push_cast

-- Finish the proof by calculations
calc
(‖z₁ * z₂‖ : ℂ) ^ 2 
_ = (z₁ * z₂) * conj (z₁ * z₂)      := by rw [← mul_conj' (z₁*z₂)]
_ = (z₁ * z₂) * (conj z₁ * conj z₂) := by rw [map_mul]
_ = (z₁ * conj z₁) * (z₂ * conj z₂) := by ring
_ = (‖z₁‖ : ℂ)^2 * (‖z₂‖ : ℂ)^2     := by rw [mul_conj' z₁ , mul_conj' z₂]
_ = (‖z₁‖ * ‖z₂‖ : ℂ) ^ 2           := by rw [mul_pow]; 

-- prove by using existing api in Mathlib
example (z₁ z₂ : ℂ) : ‖z₁ * z₂‖  = ‖z₁‖ * ‖z₂‖ := by exact Complex.norm_mul z₁ z₂
\end{comment}
Substituting $z_1 = a+bi$ and $z_2 = c+di$ into part (ii) yields the classical identity
\begin{equation}
	(ac - bd)^2 + (ad + bc)^2 = (a^2 + b^2)(c^2 + d^2).
	\label{eq:sum_of_square}
\end{equation}

The Euclidean distance between two complex numbers can be expressed in terms of the complex absolute value:
\[
d(z_1, z_2) \triangleq |z_1 - z_2|.
\]
It is easy to check that $d$ is nonnegative, symmetric, and vanishes only when $z_1 = z_2$.  
The remaining property is the triangle inequality.

\begin{svgraybox}
	\begin{proposition}[Triangle inequality] 		\label{prop:triangle_inequality} \index{Triangle inequality}
		For any $z_1, z_2 \in \mathbb{C}$,
		\[
		|z_1 + z_2| \le |z_1| + |z_2|.
		\]
	\end{proposition}
\end{svgraybox}

Although the triangle inequality is geometrically clear, we give an algebraic proof to show its compatibility with complex arithmetic.

\begin{proof}
	Compute the square of the left-hand side:
	\begin{align}
		|z_1 + z_2|^2
		&= (z_1 + z_2)(z_1^* + z_2^*) \notag \\
		&= |z_1|^2 + 2\real(z_1 z_2^*) + |z_2|^2 \notag \\
		&\le |z_1|^2 + 2|\real(z_1 z_2^*)| + |z_2|^2.
		\label{eq:triangle_inequality_proof1}
	\end{align}
	
	Since $|\real(w)| \le |w|$ for any complex number $w$,
	\begin{equation}
		|\real(z_1 z_2^*)|
		\le |z_1 z_2^*|
		= |z_1|\,|z_2|.
		\label{eq:triangle_inequality_proof2}
	\end{equation}
	
	Substituting \eqref{eq:triangle_inequality_proof2} into \eqref{eq:triangle_inequality_proof1} gives
	\[
	|z_1 + z_2|^2
	\le (|z_1| + |z_2|)^2.
	\]
The proof is complete after taken square of both sides.	\iflean \href{https://live.lean-lang.org/#codez=JYWwDg9gTgLgBAWQIYwBYBtgCMB0AVJAYxmEICgyIwBTAOzgGEJx1qAPCgWk7gBFqAZsFrU4aUQFcAzhKTo4tCDBTAI9AdDiFmYVmy1qAVhIDmKamUXKSauACJttQ3bgBeAHxwAFFOVQASsImAKK0ACZwgECEAJQUAPScZHB4qNTQ1CBwAIw4AMw4AOwANMlQwEi0JqxwwtQAjrKYMACecAA+AF6AgQQA1B2AQQRtcIAmRO3dQz1jg2SccVw80qLswL5BYqnpIFI19MhomFhk7Egsol7dcANwAFxR0TdwgGgEF5MDjyNP3e+Tz/3v1644FhWsdiApoCAAPpIMJhSHVC4DeZwMLUZTAVgRMBQCAQARHNgnXRnRH9B4xB7PLpwV5/D5U76fOkAoHNJJwbhwADKwkIokUUC2cCQUH5ak4IjMJAAbtQSvpAW02G0cHB2ZyAOqiUUAdzKMFE4jgJggcjgAEE3HAAEJiCDtc1DRXWoZSe1G1ggEBIODemCEVJSHDsqA6uAAbUACYTCrBSSF4yGKWiS7wCqFJlOQ6IAXRDYajMbjCYz1BM3hhcJLZa8acTakz9xrELrydLcCzOYocA5PHNYF0rSN5oACgBJW5tR1WtouuCAFMILQA9ABMVutK7zEcjUjq8OokJ3ue77NQSFlcFQlJelzpi7gq8BDJXNKZj2fk1XACpvBdv45DDe0Q4KKNyAsCardt2nJMLQsqwHAAByELPjA9r+NQZrDiK8B4j+cB/kYlyxJBcChlu4KClyu71HWgoHnUR6QWRBZ4uhciQiBTaUXUP7UrS0SdiRzHRiAEjoJC/6Md2zGBJUAASzA4N6YDQrCUmkUEiYCFBPDBGwYAVBERpyCY1BYFASCkHA2IQGEEjELcHRZE536vM5y6uZcy4uS+HTeR57I6XwKwwGUWASAacAAAailFwqEDiUjbMZsLADYtCBaR+YVhx1DqdBahwfA5z4QYThEcBohYEQADWdqfG+y6ZZqohhO6qArHAGhQJcWTCuEXnBiRnIpKIUj1BIdB8jUUi3CVBHlR0QEgZwnhcSAVEVStnhpptHQcp4zyNZlUigGAcBqOgrThqJ4n/iUrEYeJoolLt1G7rWh6ZcxOUvcKsK5flPAANLUK0vjUGAjlfr136dUa/6mOYF3aU5pV+Z+zW9gNPraKc+i0BIIBYNQPWGTUMDbIjUqiKe55GlIJyiGx8gGbAQ2Qad4AXbQV0Ro97GikD3JnZgAiDqkZXGDTYikyAtz/ucMP/r19yrb1pUq35335tG/65T+ZILQB3Q5hBkGcjyLDAOLcDoV4gCJhM8WSNfcqGfC7G5CfmymQrdJT6/dvpiRJzAgMLRqir4M3CugpnmUg31BF2WhyOQJFUr5t73ub3aQlaT6rj8bzvjn37nNSxuAZVoGsllEaoMLfYDnAo4iA0chpc0tz2zq9wADyKjqQyZfnoyFy+vyNS+MOZ1Xby1+BmUkSYjgmD1nI5NyoXAGA2xIDqSAS6IeMgLYBqCtsPgd1H/VGZLAhEKhPXLsRJGQUg/Z86Ke60ZknLLlVM3TA1Bti93uOPHUjwsY2lFEgOqNk7Jgk6H1DGToxhZCGJ0ZcbRR4Fy+KXKepcYZeAZKVR4Ws/j3BZOBZitZbrqXzo+AhRdXzEPIV8UqpDKGPGoWBVoZRKh4OYV0aexc/jsO/GQ8uU8+F1zoc2SSMCABiT9NA0CgAIagYIdyyFFMI7wZDxG8LgHeVcNCBFBCAA}{(\LEAN proof)} \fi
\end{proof}

\begin{comment}
import Mathlib.Tactic

open Complex

-- Define the usual notation for complex conjugate
notation "conj" => (starRingEnd ℂ)

/-
Theorem 1.3.7, Triangle inequality |z₁+z₂| ≤ |z₁| + |z₂|
-/

-- use existing theorems in Mathlib
example (z₁ z₂ : ℂ) : ‖z₁ + z₂‖ ≤ ‖z₁‖ + ‖z₂‖ := by exact norm_add_le z₁ z₂

-- detailed proof
example (z₁ z₂ : ℂ) : ‖z₁ + z₂‖ ≤ ‖z₁‖ + ‖z₂‖ := by
-- Since norms are non-negative, x = |x|. 
-- We rewrite the goal A = B to |A| = |B| so the lemma matches.
rw [← abs_of_nonneg (norm_nonneg _)]
rw [← abs_of_nonneg (add_nonneg (norm_nonneg _) (norm_nonneg _))]

-- Apply the API: |A| = |B| ↔ A^2 = B^2
rw [←sq_le_sq]

have h : ‖z₁ + z₂‖ ^ 2 = ‖z₁‖^2 + ‖z₂‖^2 + 2 * (z₁ * conj z₂).re := by 
-- Convert Norm^2 to Real Part of (z * conj z)
rw [← normSq_eq_norm_sq]
rw [← ofReal_re (normSq (z₁ + z₂))]
rw [← mul_conj]
rw [RingHom.map_add]
ring_nf  -- Expand the algebraic product: z1z1* + z1z2* + z2z1* + z2z2*
-- Distribute `re` across the addition
rw [add_re]
-- Convert (z * conj z).re back to ‖z‖^2
-- We do this for z1 and z2.
-- The sequence is: (z * conj z).re -> (normSq z).re -> normSq z -> ‖z‖^2
simp only [mul_conj, ofReal_re, normSq_eq_norm_sq]
rw [add_re, add_re]
-- Key step: z2 * z1* is the conjugate of z1 * z2*
-- And a complex number and its conjugate have the same Real part.
simp only [ofReal_re]
-- Simplify the conjugate term: conj(z2 * conj z1) -> z1 * conj z2
rw [← conj_re (z₂ * conj z₁)] 
-- Simplify Re(↑‖z1‖^2) to ‖z1‖^2
rw [map_mul, conj_conj, mul_comm]
-- the rest is algebra
ring

calc
‖z₁ + z₂‖ ^ 2 
_ = ‖z₁‖^2 + ‖z₂‖^2 + 2 * (z₁ * conj z₂).re := by rw [h]
-- Apply Inequality: Re(w) <= |w| 
_ ≤ ‖z₁‖^2 + ‖z₂‖^2 + 2 * ‖z₁ * conj z₂‖ := by 
gcongr -- 1. Strips away the common terms (squares and the factor 2)
apply re_le_norm -- 2. Applies Re(w) ≤ ‖w‖
-- Break product |z1 z2*| = |z1| |z2|
_ = ‖z₁‖^2 + ‖z₂‖^2 + 2 * (‖z₁‖ * ‖conj z₂‖) := by rw [norm_mul]
_ = ‖z₁‖^2 + ‖z₂‖^2 + 2 * ‖z₁‖ * (‖conj z₂‖) := by ring
_ = ‖z₁‖^2 + ‖z₂‖^2 + 2 * ‖z₁‖ * (‖z₂‖) := by rw [norm_conj]
-- Factor perfect square
_ = (‖z₁‖ + ‖z₂‖) ^ 2 := by ring
\end{comment}

\section{Polar form and DeMoivre formula}

A nonzero complex number, like a vector in $\mathbb{R}^2$, can be written in \df{polar form}.  
Given any nonzero complex number $x + iy$, we may write
\[
x + iy = r(\cos\theta + i\sin\theta),
\]
where $r > 0$ and $\theta$ is a real number.  
The number $r$ is the modulus of $x+iy$, representing the length of the corresponding vector, and $\theta$ is the angle from the positive real axis to that vector (Fig.~\ref{fig:polar_form}). In contrast to polar form, we will call $x+iy$ the Cartesian form of a complex number. 

\begin{figure}
	\begin{center}
		
		\begin{tikzpicture}[>=Latex, scale=1.4]
			% Axes
			\draw[->] (-0.5,0) -- (4.0,0) node[below right] {$\real$};
			\draw[->] (0,-0.5) -- (0,2.4) node[above left] {$\imag$};
			
			% Origin
			\coordinate (O) at (0,0);
			
			% Point z in polar form
			\coordinate (Z) at (3,1.8); % example coordinates
			
			% Vector for z
			\draw[->, thick, blue] (O) -- (Z) node[above right] {$z = r(\cos\theta + i\sin\theta)$};
%			\fill[blue] (Z) circle (1.5pt);
			
			% Projections
			\draw[dashed] (Z) -- (3,0) node[below] {$r\cos\theta$};
			\draw[dashed] (Z) -- (0,1.8) node[left] {$r\sin\theta$};
			
			% Angle arc
			\draw[thick] (1,0) arc[start angle=0, end angle=31, radius=1cm];
			\node [right] at (0.9,0.35) {$\theta$};
			
			% Label r
			\node[blue, below] at (1.7,1.0) {$r$};
		\end{tikzpicture}
	\end{center}
	\caption{Polar form diagram}
	\label{fig:polar_form}
\end{figure}



Using the matrix representation of complex numbers (Section~\ref{sec:complex_numbers_matrices}), the $2\times 2$ matrix corresponding to $a+bi$ can be written as
\[
\begin{bmatrix} a & -b \\ b & a \end{bmatrix}
=
r \begin{bmatrix} \cos\theta & -\sin\theta \\ \sin\theta & \cos\theta \end{bmatrix}.
\]
The matrix
\[
\begin{bmatrix} \cos\theta & -\sin\theta \\ \sin\theta & \cos\theta \end{bmatrix}
\]
is the standard rotation matrix, representing counterclockwise rotation by $\theta$.

\begin{svgraybox}
	\begin{proposition} \label{prop:polar_mult}
		Let
		\[
		z_1 = r_1(\cos\theta_1 + i\sin\theta_1), \qquad
		z_2 = r_2(\cos\theta_2 + i\sin\theta_2)
		\]
		be two complex numbers in polar form. Then
		\[
		z_1 z_2 = r_1 r_2 \big[\cos(\theta_1 + \theta_2) + i\sin(\theta_1 + \theta_2)\big].
		\]
	\end{proposition}
\end{svgraybox}

\begin{proof}
	Using the definition of complex multiplication and the angle addition formulas,
	\begin{align*}
		z_1 z_2
		&= r_1 r_2 (\cos\theta_1 + i\sin\theta_1)(\cos\theta_2 + i\sin\theta_2) \\
		&= r_1 r_2\big[\cos\theta_1\cos\theta_2 - \sin\theta_1\sin\theta_2
		+ i(\cos\theta_1\sin\theta_2 + \sin\theta_1\cos\theta_2)\big] \\
		&= r_1 r_2\big[\cos(\theta_1 + \theta_2) + i\sin(\theta_1 + \theta_2)\big].
	\end{align*}
\end{proof}

Thus, multiplication by a complex number
\[
a+bi = r(\cos\theta + i\sin\theta)
\]
has a clear geometric interpretation:  
(i) rotate by $\theta$ counterclockwise, and  
(ii) scale by a factor of $r$.  
In particular, multiplication by $i$ corresponds to a $90^\circ$ counterclockwise rotation.

\begin{example} \label{ex:Steward}
	This example is adapted from \cite{StewardTall}.  
	In Fig.~\ref{fig:Steward}, the circle has radius~$1$ and is centered at the origin.  
	The points $u$, $v$, and $w$ have modulus $1$, less than $1$, and greater than $1$, respectively.
	
	The sums and products of various combinations of $u$, $v$, and $w$ are shown as points $a$ through $h$.  
	Locate the points
	\[
	u+v,\quad u+w,\quad v+w,\quad u+v+w,
	\]
	\[
	uv,\quad uw,\quad vw,\quad uvw.
	\]
\end{example}


\begin{figure}
	\begin{center}
		\begin{tikzpicture}[scale=3]
			
			% angle parameter (in degrees)
			\def\theta{20}
			\def\xu{cos(\theta)}
			\def\yu{sin(\theta)}
			%\def\xv{0.55}
			%\def\yv{-0.15}
			%\def\xw{-1.6}
			%\def\yw{1.1}
			\def\xv{-0.55}
			\def\yv{-0.15}
			\def\xw{1.6}
			\def\yw{1.1}
			
			
			% axes
			\draw[->] (-2,0) -- (1.7,0) node[right] {$x$};
			\draw[->] (0,-1.3) -- (0,1.7) node[above] {$y$};
			
			% label the point (0,1)
			\filldraw (1,0) circle (1pt);
			\node[above right] at (1,0) {$1$};
			
			
			% unit circle
			\draw[thick] (0,0) circle (1);
			
			% point on the circle at angle \theta
			\coordinate (u) at ({\xu}, {\yu});
			\coordinate (v) at (\xv, \yv);
			\coordinate (w) at (\xw, \yw);
			\filldraw[blue] (u) circle (1.2pt) node[above right] {$u$};
			\filldraw[blue] (v) circle (1.2pt) node[below right] {$v$};
			\filldraw[blue] (w) circle (1.2pt) node[above left] {$w$};
			
			\coordinate (a) at ({\xu+\xv}, {\yu+\yv});
			\coordinate (b) at ({\xu+\xw}, {\yu+\yw});
			\coordinate (c) at ({\xv+\xw}, {\yv+\yw});
			\coordinate (d) at ({\xu+\xv+\xw}, {\yu+\yv+\yw});
			\filldraw[red] (a) circle (1.2pt) node[above right] {$b$};
			\filldraw[red] (b) circle (1.2pt) node[below left] {$h$};
			\filldraw[red] (c) circle (1.2pt) node[above right] {$f$};
			\filldraw[red] (d) circle (1.2pt) node[above left] {$g$};
			
			
			\coordinate (e) at ({\xu*\xv-\yu*\yv}, {\yu*\xv+\yv*\xu});
			\coordinate (f) at ({\xu*\xw-\yu*\yw}, {\yu*\xw+\yw*\xu});
			\coordinate (g) at ({\xv*\xw-\yv*\yw}, {\yv*\xw+\yw*\xv});
			
			\def\xp{(\xu*\xv-\yu*\yv)}
			\def\yp{(\yu*\xv+\yv*\xu)}
			\coordinate (h) at ({\xp*\xw-\yp*\yw}, {\yp*\xw+\yw*\xp});
			\filldraw[red] (e) circle (1.2pt) node[above right] {$a$};
			\filldraw[red] (f) circle (1.2pt) node[below left] {$e$};
			\filldraw[red] (g) circle (1.2pt) node[below right] {$c$};
			\filldraw[red] (h) circle (1.2pt) node[below right] {$d$};
			
			
			% optional: draw radius to the point
			%\draw[dashed] (0,0) -- (P);
			
			% optional: mark the angle
			%\draw[->] (0.4,0) arc (0:\theta:0.4);
			%\node at ({0.55*cos(\theta/2)},{0.55*sin(\theta/2)}) {$\theta$};
		\end{tikzpicture}
	\end{center}
	\caption{Example~\ref{ex:Steward}. The points $a$ to $h$ are obtained by adding or multiplying some combinations of $u$, $v$, and $w$.}
	\label{fig:Steward}
\end{figure}



\begin{svgraybox}
	\begin{definition} \index{Argument} \index{Phase} \index{Angle}
		The \df{argument} of a nonzero complex number $z$ is the angle from the positive real axis to the line segment joining $0$ to $z$.  
		We denote it by $\arg(z)$.  
		The argument is also called the \df{phase} or \df{angle} of $z$.
	\end{definition}
\end{svgraybox}

The argument of a complex number is defined only up to integer multiples of $2\pi$, since rotating by $2\pi$ radians returns to the same point.   The argument of $0$ is not defined.


\begin{example}
	The number $-1$ can be written in polar form as
	\[
	-1 = 1\cdot(\cos(\pi + 2\pi k) + i\sin(\pi + 2\pi k)),
	\]
	where $k \in \mathbb{Z}$.  
	Thus
	\[
	\arg(-1) = \pi + 2\pi k, \qquad k\in\mathbb{Z}.
	\]
\end{example}

The function $\arg(z)$ is a \emph{multi-valued} function, meaning that there are more than one possible values for $\arg(z)$. To be more rigorous, we may interpret the argument function as a {\em set-valued function}. For instance, we may define the $\arg(-1)$ as the set $\{\pi+2\pi k:\, k\in \mathbb{Z}\}$. However, for simplicity, we will just write $\pi+2\pi k$, and understand that the variable $k$ can assume any integral value.

To remove the ambiguity in the argument, one often restricts $\theta$ to a specific interval.  
Common choices are  
(i) $(-\pi, \pi]$, or  
(ii) $[0, 2\pi)$.  
The unique value of the argument lying in the chosen interval is called the {\em principal argument}.  
This choice is purely conventional and serves to simplify notation.

\begin{svgraybox}
	\begin{definition} \label{def:principal_argument} \index{Argument!principal}
		If we choose the interval $(-\pi, \pi]$ as the range, then for any nonzero complex number $z$ the \df{principal argument} $\theta_0$ is the unique angle in $(-\pi, \pi]$ satisfying
		\[
		\cos\theta_0 = \frac{\real(z)}{|z|}, \qquad
		\sin\theta_0 = \frac{\imag(z)}{|z|}.
		\]
We write \df{$\Arg(z)$} for the principal argument.
	\end{definition}
\end{svgraybox}

The choice of the interval $(-\pi,\pi]$ as the range of $\Arg(z)$ agrees with convention of the argument function in Python.








\begin{example}
	If we take $(\pi, \pi]$ as the range, the principal argument of $-1-i$ is $-3\pi/4$.
\end{example}



The DeMoivre formula describes how the argument of a complex number behaves when the number is raised to an integer power.

\begin{svgraybox}
	\begin{theorem}[DeMoivre formula] \index{DeMoivre formula}
		For any $n \in \mathbb{Z}$ and $\theta \in \mathbb{R}$,
		\begin{equation}
			(\cos\theta + i\sin\theta)^n
			= \cos(n\theta) + i\sin(n\theta).
			\label{eq:DeMoivre}
		\end{equation}
	\end{theorem}
\end{svgraybox}

\begin{proof} 
	The formula is trivial when $n=0$ and $n=1$.  
	Assume it holds for some positive integer~$n$. Then
	\begin{align*}
		(\cos\theta + i\sin\theta)^{n+1}
		&= (\cos\theta + i\sin\theta)^n(\cos\theta + i\sin\theta) \\
		&= (\cos n\theta + i\sin n\theta)(\cos\theta + i\sin\theta) \\
		&= (\cos n\theta \cos\theta - \sin n\theta \sin\theta) \\
		&\qquad {} + i(\cos n\theta \sin\theta + \sin n\theta \cos\theta) \\
		&= \cos((n+1)\theta) + i\sin((n+1)\theta).
	\end{align*}
	Thus the formula holds for all nonnegative integers.
	
	For negative $n$, write $n=-m$ with $m>0$. Then
	\begin{align*}
		(\cos\theta + i\sin\theta)^n
		&= \frac{1}{(\cos\theta + i\sin\theta)^m} \\
		&= \frac{1}{\cos m\theta + i\sin m\theta}
		\cdot \frac{\cos m\theta - i\sin m\theta}{\cos m\theta - i\sin m\theta} \\
		&= \cos m\theta - i\sin m\theta \\
		&= \cos(n\theta) + i\sin(n\theta).
	\end{align*}
	This proves the formula for all integers $n$.  \iflean \href{https://live.lean-lang.org/#codez=JYWwDg9gTgLgBAWQIYwBYBtgCMB0AVJAYxmEICgIwBTAOzgGEJx0qAPMsgegFoy9Uq0KiDgBGHABZJcACJUEEYADcoVOADNoIAK7okZAGLQ4wGjCoBzKlDg0ANGQAUhCAGc4gDuA4AahMAqOFdTTwBKAD06AF44F3dHGg8QnxM4AKC6eMSybk4OLl4AdTV1YChXeDAoCCU1OQVlVQ0tXSQmmxoIGhpLFGU1U3MrMuzc7O44AAUqmrEALknptSwATxMaABNtYmBOsjRBVRE6xRU1Ry95wFxCJPi4ecBUQiTZp1jPZPS3gIBJJLDbOGir1uAUS72CwNCqTgXzu0RWYzEODgAGUYEhYNCNlsSJ04LiaGRTJttvi4AB3YBoMgAHzgdIAXtYIACAHwIgBCSFcakIXKo8yicAADAiADIACWRALgjhwcvCQulogRACVJdLXorfB9Fd8lclFdFlUFwDS6YEtoR/sBUKyEV8sdsZuUqGB5gBBVyubQgNQwKDaYrGexwSrVIPtZLK7jjFVUMlQSlqelheLeURJLlwFN0AL0hHusBgdCrfZrYk4uioZaQfZBVxkKBkuAAbUgZIA+t7CIQ7CZUABdOAcGNwcVIDYsOBl9EWH20GDuUxBdY8txwCfrQKmea3XwZqFeaK5t6+DxkE1gPE0EutgByKBwvPKXctfaQ63WHZ06D7nSo366AOI7cAAqtyDBMMWbA4Gw8AwMyhACIQADWcBxkg6Abhs0IgEgFimOiqxgOii6BK66IoFQJZkEgRa3owzAwXBZAAO3buAraxB2H7rH26Q8Z+wF0omNAWGxHFXi2Am8X23G8cJcCieJCJTBAEDqHAABM8xqTMKzTgI05ECQVoAAZEtiOw0AA5GZewCEIRzyCcqg2TKFxwNcMp0A8TwvOuXjasEXjfL8/yAuuEKgsFGQnqCeqzHCyyEo6lbuXQFJoHAIiWSSNAdtWtYCPWEmjlgfIxHyZB0peZXjHllaBOYYA1UpzZthAnbdr25ZWZ0hU1hAdbAK4Q5tZe163i2D4wE+XIwK+Pbvp+gG/teAE/opdHFqsjHQawsGsDAbXsZNLbyZ+/GmIJ6yKSJpjiXSZ2gFJMlXTEbi3fdSmPSBCzqZpADMumLHA2jpBYcAAKK6NYNnuJoUA/voo78KNJjuBOcBsKNJBiYZBzCGsiAoBg2AOUTzn1KcNnuecdxeTcvlwI8dwBe4QW1SFUI/BE0rOFFNB+DFdKOB88QiyESRhbCcDwmwJmfa4t1djdP4dl8Hbtv8575GQd5wX2YYzGWxwNBGKNtLYPQkDMbCQN0Zg4GQRQ21QW4wGSzICOg1A2CwIB4Q2OQcDwY7Uf7cCB3hYiU2sNRlGoGk46mrB+D8mM4wAjtomHTsyDuONw6c/CMZAx60/4dusShfhrMKOKwjPefMoicI4DsyqX0vSl3xfdxnvdJfLKWoEgMyoPM/c93AgAGRMKcsGYrxA46wYAdt0HaMlUg9lyU1Ffudg4XssQcIgUE7wcyrioF1a9Xk3vMyw/MrcE/YVJEaZC8ug5Af9CEIR0wABGLgAn4QCu5kGiI4cBSRfBgObp/OAkD15L1WE2VsgAEwgfrdPeNx349xCMBDsfc0GOCFE8ZK7VWy8U3pYDsvIaCEGoiQ/UI8DKYJbA7beTJgL5Ajn7aw0dhCxy0vHUwicIIp0FnEVg8DtwZHkakTOVBc6YXcECEuCiJbaJUeXSusgXIW0KpHYRT8rhPGHO3WR3cdHgmUUPDUUU9GxTfo4mEHCUpcJwTwhucBWDAS4dXWu9ddCayCR1Xx681qa3ccQ2i9FVguDElAHiUAoZdwAPwdkbH9f65tThtCtkjG2FhegzAGJYYRDt/xmEpk5IxNNVD0IsB2GeLdmaMzZs8GxQJlFuKbjLPmgotHCwGTKXR4ze6JWShwcek8ZCMz6VFCZEsPHhFGULPwxd5EKMUT5HZ2iX6eLmSJKJG0a7KG1vfWxQywRKOGUkGgikuGFJab7f2ikV7wEukfdW4StY6wyNojgv9yArLkfYx5KjNnOLiMLXZvdfAHMlkik56CaEtlQDINhkU4jcGmckVFhK/COOhJi5Sm91Bh14O8y2LRraYSwlUoYIdcj7EafSjsSMrYM0sT5RmgASQn8rYrmBzQqAP5jAoEwtRaTPBHKmZgD0E/z5O4TKlJUBmg0rNHKdpRz0EqqIAUWdWgdBoIS22fQ1iDHMYKIUfZRB9i0nKHAIQL5rhoIneCRkHR2sjBVVCBc4CzW0FAfONAfRYGESAC8r0CxJMJhoUo5RCaNNKRa7o5S7b9DMNU4YPymmuTUF4ONZps3IktPqyIbJDWVR0taLGZSKl5oDYK6I3BnVwG4K6+UCJK3Vtyu4cwSEaCkGZasd+Ih9weoROISC3rrDwFZcI1wywzBIGbghJSVBI3RuEeiLVvpTLPFHOaCUyJ5j0jgH8ad6YkgxpcL6dwog4CcGzKmEAD6ETmjVFe/4IJ5ZUGfVQdw96DwgnjZxTo01/U4EHT2dp2c+z0nbK0vs8HnyLWzfw0crrQ33zLBYCA+c8IwCQmBuAZluXZvaevMyZlQzWHUCBmANFz10gAOKkfQG3d9Mo3XhUcN+g8+K34iYfVLOAbrGwdUkchjsITlDASLbRhhXcvAiajCEIAA}{(\LEAN proof)} \fi
\end{proof}



Using the matrix representation of complex numbers, DeMoivre's formula becomes
\[
\begin{bmatrix}
	\cos\theta & -\sin\theta \\
	\sin\theta & \cos\theta
\end{bmatrix}^n
=
\begin{bmatrix}
	\cos(n\theta) & -\sin(n\theta) \\
	\sin(n\theta) & \cos(n\theta)
\end{bmatrix}.
\]
Geometrically, this states that rotating by an angle $n$ times is equivalent to a single rotation by $n$ times the angle.

\begin{remark} Because the expression ``$\cos \theta + i \sin \theta$'' occurs very frequently,
  some people use the short-hand notation $$cis(\theta) \triangleq \cos \theta + i \sin \theta.$$
 Once we have proved Euler's identity, we will write it as $e^{i\theta}$.
\end{remark}



\begin{example}
	Compute $(-1 + i\sqrt{3})^8$.
	
	The number $-1 + i\sqrt{3}$ has polar form
	\[
	2\!\left(-\tfrac12 + i\tfrac{\sqrt{3}}{2}\right)
	= 2(\cos(2\pi/3) + i\sin(2\pi/3)).
	\]
	Thus
	\[
	(-1 + i\sqrt{3})^8
	= 2^8(\cos(8\cdot 2\pi/3) + i\sin(8\cdot 2\pi/3))
	= 256(\cos(4\pi/3) + i\sin(4\pi/3)).
	\]
	In Cartesian form, this is
	\[
	128(-1 - i\sqrt{3}).
	\]
\end{example}

\medskip

\begin{example}
	Compute $(-1+i)^{20}$.
	
	We write
	\[
	-1+i = \sqrt{2}\,(\cos(3\pi/4) + i\sin(3\pi/4)).
	\]
	Then
	\begin{align*}
		(-1+i)^{20}
		&= (\sqrt{2})^{20} \big(\cos(20\cdot 3\pi/4) + i\sin(20\cdot 3\pi/4)\big) \\
		&= 2^{10}(\cos\pi + i\sin\pi) \\
		&= -1024.
	\end{align*}
\end{example}

\medskip

\begin{example}
	Express $\sin(5\theta)$ as a polynomial in $\sin\theta$.
\end{example}

Using DeMoivre's formula,
\begin{align*}
	\cos(5\theta) + i\sin(5\theta)
	&= (\cos\theta + i\sin\theta)^5 \\
	&= \cos^5\theta
	+ 5i\cos^4\theta\sin\theta
	- 10\cos^3\theta\sin^2\theta \\
	&\quad {} - 10i\cos^2\theta\sin^3\theta
	+ 5\cos\theta\sin^4\theta
	+ i\sin^5\theta.
\end{align*}
Taking imaginary parts,
\begin{align*}
	\sin(5\theta)
	&= 5\cos^4\theta\sin\theta
	- 10\cos^2\theta\sin^3\theta
	+ \sin^5\theta \\
	&= 5(1-\sin^2\theta)^2\sin\theta
	- 10(1-\sin^2\theta)\sin^3\theta
	+ \sin^5\theta \\
	&= 5\sin\theta - 20\sin^3\theta + 16\sin^5\theta.
\end{align*}

\smallskip

We can use DeMoivre's formula to extract the $n$-th roots of a complex number. Let $w \neq 0$ be written in polar form
\[
w = |w|(\cos(\Arg w) + i\sin(\Arg w)).
\]
We seek all $z$ such that $z^n = w$.  
Write $z = r(\cos\theta + i\sin\theta)$. Then
\[
z^n = r^n(\cos(n\theta) + i\sin(n\theta)) = w.
\]
Thus
\[
r = |w|^{1/n},
\qquad
\theta = \frac{\Arg w + 2k\pi}{n},
\quad k = 0,1,\ldots,n-1.
\]
Therefore the $n$-th roots of $w$ are
\begin{equation}
	|w|^{1/n}
	\cos\!\left(\frac{\Arg w + 2k\pi}{n}\right)
	+ i\,|w|^{1/n}
	\sin\!\left(\frac{\Arg w + 2k\pi}{n}\right),
	\label{eq:nth_root}
\end{equation}
for $k = 0,1,\ldots,n-1$. The $n$-th root of a nonzero complex number $w$ are the vertices of a regular $n$-gon in the complex plane.

The case of the $n$-th roots of unity is of special importance.

\begin{svgraybox}
	\begin{definition} \index{Roots of unity}
		Given a positive integer $n$, a complex number $\zeta$ is called an \df{$n$-th root of unity} if $\zeta^n=1$.
	\end{definition}
\end{svgraybox}

An example of seventh roots of unity is shown in Fig.~\ref{fig:roots_of_unity}.

We can express the complex $n$-th roots of any nonzero complex number $w$ using $n$-th roots of unity. Suppose $z_0$ is one of the $n$-th roots of $w$. If $z_1$ is another $n$-th root of $w$, then the ratio $z_0/z_1$ satisfies
$$
\big( \frac{z_0}{z_1} \big)^n = \frac{z_0^n}{z_1^n} = \frac{w}{w} = 1.
$$
Hence, $z_0/z_1$ is a root of unity. We can write $z_1$ as
$$
z_1 = z_0 \zeta
$$
for some $n$-th root of unity $\zeta$. Conversely if $\zeta$ is an $n$-th root of unity, $z_0\zeta$ is also an $n$-th root of unity. Therefore, once we have obtain one particular $n$-th root $z_0$ of $w$, we can obtain the complete set of the $n$-th roots of $w$:
$$
z_0,\  z_0 \zeta,\  z_0 \zeta^2 ,\  \ldots,\ z_0 \zeta^{n-1}.
$$


\begin{example}
Let $n$ be a positive integer. To find all the $n$-th root of a positive real number $a$, we first note that $\sqrt[n]{a}$ is one of the solutions to $z^n=a$. The $n$ distinct roots are characterized by $\sqrt[n]{a} \zeta^k$, where $\zeta = \cos(2\pi/n) + \sin(2\pi/n)$, and $k=0,1,2,\ldots, n-1$. 
\end{example}




\begin{figure}
	\begin{center}

\begin{tikzpicture}[scale=2, >=Latex]
	
	% Number of roots
	\def\n{7}
	
	% Unit circle
	\draw[thick] (0,0) circle (1);
	
	% Axes
	\draw[->] (-1.3,0) -- (1.3,0) node[right] {$\real$};
	\draw[->] (0,-1.3) -- (0,1.3) node[above] {$\imag$};
	
	% Roots of unity
	\foreach \k in {0,...,\numexpr\n-1} {
		\coordinate (Z\k) at ({cos(360*\k/\n)},{sin(360*\k/\n)});
		\fill (Z\k) circle (0.03);
		\draw[thick, blue, ->] (0,0) -- (Z\k);
		
		% Label each root
%		\node[anchor=south east] at (Z\k) {$\omega^\k$};
	}
	
	\node[above right] at (Z1) {$\omega^1$};
	\node[above left] at (Z2) {$\omega^2$};
	\node[above left] at (Z3) {$\omega^3$};
	\node[below left] at (Z4) {$\omega^4$};
	\node[below left] at (Z5) {$\omega^5$};
	\node[below right] at (Z6) {$\omega^6$};
	\node[above right] at (Z0) {$\omega^0$};		
	% Optional: label the unit circle
%	\node at (0.7,0.7) {$|z|=1$};
	
\end{tikzpicture}
	\end{center}
\caption{The powers of $\omega = (\cos (2\pi/7) + i \sin (2\pi/7))$ are the seventh roots of unity.}
\label{fig:roots_of_unity}
\end{figure}




\section{Applications in plane geometry}


\subsection{A criterion of collinear points}

We start with the simple but useful observation that
$$
\arg(z_1/z_2) = \arg(z_1) - \arg(z_2) \bmod 2\pi.
$$
This device allows us to compute the angle between two line segments in terms of the argument function.

We can test whether three points $z_1$, $z_2$ and $z_3$ lies on the same line in terms of complex operations, without going to the real domain. The idea is: if we divide $z_1-z_2$ by $z_3-z_2$, the argument of the quotient is the angle between the line segment between $z_1$ and $z_2$, and the line segment between $z_3$ and $z_2$. They are collinear if and only if the angle is $0$ or $\pi$,   
$$
\arg \Big( \frac{z_1-z_2}{z_3-z_2} \Big) = 0 \text{ or } \pi.
$$
Equivalently, this condition can be written as
\begin{equation}
\imag \Big( \frac{z_1-z_2}{z_3-z_2} \Big) = 0.
\label{eq:three_points_collinear}
\end{equation}





\subsection{Points lying on an equilateral triangle}

We can derive a condition to test whether three complex points form an equilateral triangle. If $z_1$, $z_2$ and $z_3$ are the vertices of an equilateral triangle, then the line segments $z_1-z_2$ and $z_3-z_2$ must have the same length, and the angle between them is equal to $\pm \pi/3$. Hence,
$$
\frac{z_1-z_2}{z_3-z_2} = \cos(\pm \pi/3) + i \sin(\pm \pi/3) = \frac{1\pm \sqrt{3}i}{2}.
$$
This condition is equivalent to
\begin{align*}
	\frac{z_1-z_2}{z_3-z_2}-\frac{1}{2} &= \frac{\pm \sqrt{3}i}{2} \\
	\Leftrightarrow \ \ (2z_1 -z_2-z_3) &= \pm \sqrt{3} i (z_3-z_2) \\
	\Leftrightarrow \ \ (2z_1 -z_2-z_3)^2  &=  -3  (z_3-z_2)^2 \\
	\Leftrightarrow \ \ z_1^2 + z_2^2 + z_3^2 &= z_1z_2 +  z_1z_3 + z_2z_3.
\end{align*}
This gives a concise criterion to determine whether $z_1$, $z_2$, and $z_3$ form an equilateral triangle.




\subsection{A proof of Napoleon's theorem}

Draw an arbitrary polygon with 4 sides. Label the vertices by $A$, $B$, $C$ and $D$. On each side, we draw a square with this side as the base. Label the midpoint of the four squares by $P$, $Q$, $R$ and $S$. Show that the line $PR$ is perpendicular to the line $QS$. (See Fig.~\ref{fig:Napoleon_problem})

\begin{figure}
	\begin{center}
		\begin{tikzpicture}[scale=1.2]
			
	\def\xA{0.6}
	\def\yA{0.2}
	\def\xB{3}
	\def\yB{0.6}
	\def\xC{4}
	\def\yC{2}
	\def\xD{0.5}
	\def\yD{2.5}
 %Define quadrilateral points
	\coordinate (A) at (\xA,\yA);
	\coordinate (B) at (\xB,\yB);
	\coordinate (C) at (\xC,\yC);
	\coordinate (D) at (\xD,\yD);
	
 %Draw quadrilateral
	\draw[black, thick] (A) -- (B) -- (C) -- (D) -- cycle;
	
 % Draw squares on each side
	\draw[blue, thick] (B) -- ++(-\yA+\yB, \xA-\xB) -- ++(-\xB+\xA, -\yB+\yA) -- (A) ; % Square on AB
	\draw[blue, thick] (C) -- ++(-\yB+\yC, \xB-\xC) -- ++(-\xC+\xB, -\yC+\yB) -- (B); % Square on BC
	\draw[blue, thick] (D) -- ++(-\yC+\yD, \xC-\xD) -- ++(-\xD+\xC, -\yD+\yC) -- (C); % Square on CD
	\draw[blue, thick] (A) -- ++(-\yD+\yA, \xD-\xA) -- ++(-\xA+\xD, -\yA+\yD) -- (D) ; % Square on DA
	
	
 %Compute centers of squares
	\coordinate (P) at ($(A)!0.5!(B)+(-\yA*0.5+\yB*0.5, \xA*0.5-\xB*0.5)$); % center of square on AB
	\coordinate (Q) at ($(B)!0.5!(C)+(-\yB*0.5+\yC*0.5, \xB*0.5-\xC*0.5)$); % center of square on BC
	\coordinate (R) at ($(C)!0.5!(D)+(-\yC*0.5+\yD*0.5, \xC*0.5-\xD*0.5)$); % center of square on CD
	\coordinate (S) at ($(D)!0.5!(A)+(-\yD*0.5+\yA*0.5, \xD*0.5-\xA*0.5)$); % center of square on BC
	
 %Draw center points
	\filldraw[purple] (P) circle (2pt) node[above left] {$P$};
	\filldraw[purple] (Q) circle (2pt) node[above right] {$Q$};
	\filldraw[purple] (R) circle (2pt) node[below right] {$R$};
	\filldraw[purple] (S) circle (2pt) node[below left] {$S$};
	
 %Mark vertices of quadrilateral
	\filldraw[purple] (A) circle (2pt) node[below left] {$A$};
	\filldraw[purple] (B) circle (2pt) node[below left] {$B$};
	\filldraw[purple] (C) circle (2pt) node[right] {$C$};
	\filldraw[purple] (D) circle (2pt) node[below left] {$D$};
	
	
 %Draw diagonals
	\draw[red, thick] (P) -- (R);
	\draw[red, thick] (Q) -- (S);
	
 %Mark intersection point
	\coordinate (O) at (intersection of P--R and Q--S);
			%	\filldraw[black] (O) rectangle ++(0.15,0.15); % right angle marker
		\end{tikzpicture}
	\end{center}
	\caption{Napoleon's theorem for quadrilateral: the lines $PR$ and $QS$ are perpendicular.}
	\label{fig:Napoleon_problem}
\end{figure}

We can use complex numbers to solve this puzzle. More precisely, we use the fact that multiplication by $i$ means rotation by 90 degrees counter-clockwise. Indeed, the matrix representation of $i$ is
$\begin{bmatrix} 0 & -1 \\ 1 & 0 \end{bmatrix}$, and the multiplication by this matrix from the left is the same as rotation by 90 degrees.

We first represent each points in the question by complex numbers. Suppose the point $A$ is represented by complex number $z_A$, point $B$ by $z_B$, etc. The complex number $z_P$ is related to $z_A$ and $z_B$ through the following equation
$$
z_P = \frac{z_A+z_B}{2} + i \frac{z_A-z_B}{2}.
$$
The first term $(z_A-z_B)/2$ is the midpoint of the line segment $AB$. The second term is the vector from the midpoint of $AB$ to the point~$P$. 

By the same argument, we can express $z_Q$, $z_R$ and $z_S$ in terms of the four vertices of the quadrilateral,
\begin{align*}
	z_Q &= \frac{z_B+z_C}{2} + i \frac{z_B-z_C}{2}, \\
	z_R &= \frac{z_C+z_D}{2} + i \frac{z_C-z_D}{2}, \\
	z_S &= \frac{z_D+z_A}{2} + i \frac{z_D-z_A}{2} 	.
\end{align*}
After some simplification, we get
\begin{align*}
	z_P &= \frac{1}{2} \Big[ (1+i)z_A + (1-i) z_B\Big], \\
	z_Q &= \frac{1}{2} \Big[ (1+i)z_B + (1-i) z_C\Big], \\
	z_R &= \frac{1}{2} \Big[ (1+i)z_C + (1-i) z_D\Big], \\
	z_S &= \frac{1}{2} \Big[ (1+i)z_D + (1-i) z_A\Big].			
\end{align*}

The two line segments $PR$ and $QS$ are represented by $z_P-z_R$ and $z_Q - z_S$, respectively,
\begin{align*}
	z_P-z_R &=  \frac{1}{2} \Big[ (1+i)(z_A-z_C) + (1-i) (z_B-z_D)\Big] \\
	z_Q-z_S &=  \frac{1}{2} \Big[ (1+i)(z_B-z_D) + (1-i) (z_C-z_A)\Big].
\end{align*}
We can now verify that 
\begin{align*}
	i(z_P-z_R) &= \frac{1}{2} \Big[ i(1+i)(z_A-z_C) + i(1-i) (z_B-z_D)\Big] \\
	&=  \frac{1}{2} \Big[ (i-1)(z_A-z_C) + (i+1) (z_B-z_D)\Big] \\
	&=  \frac{1}{2} \Big[ (1-i)(z_C-z_A) + (1+i) (z_B-z_D)\Big] \\
	&= z_Q-z_S.
\end{align*}
This proves not only that $PR$ is perpendicular to $QS$, it also proves that they have the same length.


\section{Appendix: Solving cubic equation using discrete Fourier transform}

\label{sec:cubic_equation} \index{Cubic equation}
We outline a method of finding the roots of a polynomial of degree 3
$$
a x^3 + b x^2 + cx + d
$$
using discrete Fourier transform (DFT). In classical literature, this is called the method of {\em Lagrange resolvent}. 
%Since discrete Fourier transform is implemented in Python, we will take advantage of this fact and write a short program that solve cube equations in general. 

The first ingredient in this method is the DFT of order 3. Let $\omega = (1/2)+i\sqrt{3}/2$ be a cube root of unity. We will use the following property
$$
1 + \omega + \omega^2 = 1
$$
about cube root of unity.
Let $W$ be  the $3\times 3$ matrix
$$
W \triangleq  \begin{bmatrix}
	1 & 1 & 1 \\ 1 & \omega &\omega^2 \\ 1 &\omega^2 & \omega
\end{bmatrix}
$$
The discrete Fourier transform of order 3 is the linear transformation that takes a column vector $\mathbf{v} = \begin{bmatrix} z_1 & z_2 & z_3 \end{bmatrix}^T$ as input and return $W \mathbf{v}$. We can show that
$$
W^H W = 3 \mathbf{I}_{3\times 3},
$$
where $W^H$ denotes the conjugate transpose of matrix $W$ (Exercise~\ref{exercise:DFT})).
The inverse transform is thus given by
$$
W^{-1} = \frac{1}{3}W^H.
$$

The second idea is to simplify the cubic polynomial by reducing the coefficients of $x^2$ to~0. We first assume without loss of generality that the leading coefficient is 1 and consider the a cubic equation in the form
$$
x^3 + bx^2 + cx + d=0.
$$
By letting $u= x + b/3$, we can transform this equation to
\begin{equation}
	u^3 + \big( c - \frac{b^2}{3}\big) u + \big( d - \frac{cb}{3} + \frac{2b^3}{27}\big)=0.
	\label{eq:A}
\end{equation}
Let $\alpha$, $\beta$ and $\gamma$ be the roots of \eqref{eq:A}. 

We re-define $c$ and $d$ by
\begin{align*}
	c &\gets c - b^2/3 \\
	d &\gets d - cb/3 + 2b^3/27,
\end{align*}
so that
$$
(u-\alpha)(u-\beta)(u-\gamma) = u^3 +cu + d.
$$
By taking the inverse DFT, we can find $y_0$, $y_1$ and $y_2$ such that
\begin{align*}
	\alpha &= y_0 + y_1 + y_2 \\
	\beta & =y_0 + y_1\omega + y_2 \omega^2 \\
	\gamma &=y_0 + y_1 \omega^2 + y_2 \omega.
\end{align*}

%In theory, the value of $y_0$, $y_1$ and $y_2$ can be obtained by
%$$
%\begin{bmatrix} y_0 \\ y_1 \\ y_2\end{bmatrix}
%= \frac{1}{3} W^H
%\begin{bmatrix} \alpha \\ \beta \\ \gamma \end{bmatrix}.
%$$
%By since we do not know $\alpha$, $\beta$ and $\gamma$ yet, we only know that such $y_0$, $y_1$ and $y_2$ certainly exists (and is uniquely determined by $\alpha$, $\beta$ and $\gamma$). 
Because the sum of roots $\alpha+\beta+\gamma$ is zero, it is guaranteed that $y_0=0$. Hence, we can write
\begin{align}
	\alpha &=  y_1 + y_2  \label{eq:1}\\
	\beta & = y_1\omega + y_2 \omega^2 \\
	\gamma &= y_1 \omega^2 + y_2 \omega \label{eq:3}
\end{align}
and our task is to solve for $y_1$ and $y_2$.

We derive the relationship between $y_1$, $y_2$, $c$ and $d$ by
\begin{align*}
	u^3 + cu + d &= (u-\alpha)(u-\beta)(u-\gamma) \\
	&= (u - y_1 - y_2)(u - y_1\omega  - y_2\omega^2)(u - y_1\omega^2 - y_2\omega) \\
	&= u^3 + \big[ (y_1+y_2)(y_1\omega+y_2\omega^2) + (y_1+y_2)(y_1\omega^2+y_2\omega) +  \\
	& \phantom{=}\ (y_1\omega+y_2\omega^2)(y_1\omega^2+y_2\omega) \big] u + (y_1+y_2)(y_1\omega+y_2\omega^2) (y_1\omega^2 + y_2 \omega).
\end{align*}
By comparing the coefficients on both sides, after some simplification,  we obtain
\begin{align*}
	c &= -3y_1y_2 \\
	d &= -(y_1^3 + y_2^3).
\end{align*}

We can solve for $y_1^3$ and $y_2^3$ by solving the following quadratic equation
$$
(y-y_1^3)(y-y_2^3) = y^2 + dy - \frac{c^3}{27}.
$$
The solutions are 
$$
\frac{1}{2}\big( -d \pm \sqrt{d^2 + 4c^3/27}\big) .
$$

We can let $y_1$ to be a cube root of one of the above solutions,
$$
y_1 := \Big( \frac{1}{2}\big( -d \pm \sqrt{d^2 + 4c^3/27}\big) \Big)^{1/3},
$$
and let $y_2$ be the complex number such that $c=-3y_1y_2$. (If $y_1$ is zero, then we cannot calculate $y_2$ by $-c/(3y_1)$ and we need to switch the role of $y_1$ and $y_2$, or assign another value to $y_1$.) 

Once we know $y_1$ and $y_2$, we can compute $\alpha$, $\beta$ and $\gamma$ by equations \eqref{eq:1} to \eqref{eq:3}. This can be done in one step by calling the DFT function. Then, the final solution is obtained by subtracting $b/3$.



\section{Appendix: Complex numbers in Python}

To emphasize the computation aspect of complex numbers, we describe how to calculate complex numbers in Python.
Complex number is a built-in type in Python. The imaginary unit is represented by the symbol $\mathsf{j}$.


For further references see \href{https://trinket.io/python3/2c41b2bfb1c0}{Sample Python code} 


\begin{lstlisting}[language=Python, frame=single]
# Basic arithemtic of complex numbers
z1 = 2 + 3j
z2 = -1 + 0.2j

print(f"The sum of {z1} and {z2} is {z1+z2}.")
print(f"The difference of {z1} and {z2} is {z1-z2}.")
print(f"The product of {z1} and {z2} is {z1*z2:.3}.")
print(f"The quotient of {z1} and {z2} is {z1/z2:.3}.")
\end{lstlisting}

The program output is:

\begin{lstlisting}[language=Python, frame=single]
The sum of (2+3j) and (-1+0.2j) is (1+3.2j).
The difference of (2+3j) and (-1+0.2j) is (3+2.8j).
The product of (2+3j) and (-1+0.2j) is (-2.6-2.6j).
The quotient of (2+3j) and (-1+0.2j) is (-1.35-3.27j).
\end{lstlisting}

\medskip

Commands for the real part, imaginary part and conjugate is part of the data structure.

\begin{lstlisting}[language=Python,frame=single]
# Real part, imaginary part, and conjugate	
	
z = 3.2-5.3j
print(f"Real Part of {z} = {z.real}")
print(f"Imaginary Part of {z} = {z.imag}")
print(f"Conjugate of {z} = {z.conjugate()}")
\end{lstlisting}

If we run this program, we get 

\begin{lstlisting}[language=Python, frame=single]
Real Part of (3.2-5.3j) = 3.2
Imaginary Part of (3.2-5.3j) = -5.3
Conjugate of (3.2-5.3j) = (3.2+5.3j)	
\end{lstlisting}


\medskip

The argument and modulus are called absolute value and phase in Python.

\begin{lstlisting}[language=Python, frame=single]
# Absolute value and phase
from cmath import phase
from numpy import abs

z = 12+34j
print (f"z = {z}")
print (f"Argument of z = {phase(z):.4} (rad)")
print (f"Magnitude of z = {abs(z):.4}")
\end{lstlisting}

The command \textsf{phase} and \textsf{abs} are imported from library \textsf{cmath} and \textsf{numpy}, respectively. The returned value of \textsf{phase()} is in radian.

\begin{lstlisting}[language=Python,frame=single]
z = (12+34j)
Argument of z = 1.232 (rad)
Magnitude of z = 36.06
\end{lstlisting}

We can check Proposition \ref{prop:polar_mult} by the Python program below.

\begin{lstlisting}[language=Python, frame=single]
from cmath import phase
from numpy import abs, sin, cos
from math import isclose

z1 = 12+34j
z2 = 3-6j
r1 = abs(z1)   # absolute value of z1
r2 = abs(z2)   # absolute value of z2
theta1 = phase(z1)   # argument of theta1
theta2 = phase(z2)   # argument of theta2
print (f"z1 = {z1}, z2 = {z2}")  
print (f"z1*z2 = {z1*z2}")   # compute the product directly

w = r1*r2*(cos(theta1+theta2) + sin(theta1+theta2)*1j)
assert isclose(abs(w- z1*z2), 0, abs_tol=1e-8)
\end{lstlisting}

The variables \textsf{z1} and \textsf{z2} can be any complex numbers. The result of \textsf{z1*z2} and the value of \textsf{w} are the same, up to numerical error.

\begin{lstlisting}[language=Python, frame=single]
z1 = (12+34j), z2 = (3-6j)
z1*z2 = (240+30j)
\end{lstlisting}


In Python, we can compute the the $n$-th root of a complex number by the command \textsf{z**(1/n)}. Raising to a rational power $m/n$ can be computed by \textsf{z**(m/n)}, for integers $m$ and $n$. Python will return one of the $n$ possible answers. Compare the following calculations in Python.




\begin{lstlisting}[language=Python, frame=single]
# Computation of roots of complex numbers in Python
z = -1+1j
w1 = (z**3)**(1/4) # raise to the power 3 and then take 4th root
w2 = z**(3/4)  # raise to the power 3/4
print(f"w1 = {w1:.5}")
print(f"w2 = {w2:.5}")

print(f"\nThey are both fourth roots of {z*z*z}.")
print(f"w1^4={w1**4:.5}")
print(f"w2^4={w2**4:.5}")
\end{lstlisting}

The number \textsf{w1} is obtained by  first raising $z$ to the third power and then taking the fourth root. The  number \textsf{w2} is computed by directly raising $z$ to the power $3/4$. Mathematically, they are supposed to be the same, but the actually results are not the same. It is because there are four possible answers when we take the fourth root. The fourth power of both of them are equal to $(-1+i)^3=2+2i$.

\begin{lstlisting}[language=Python,frame=single]
w1 = (1.2719+0.253j)
w2 = (-0.253+1.2719j)

They are both fourth roots of (2+2j).
w1^4=(2+2j)
w2^4=(2+2j)
\end{lstlisting}


Using complex numbers, we can implement the condition of collinearity in \eqref{eq:three_points_collinear} by Python.

\begin{lstlisting}[language=Python, frame=single]
# Determine whether three points lie on a straight line
# Input three tuples:  pt1= (x1,y1), pt2= (x2,y2), pt3= (x3,y3)
# return True if the three points pt1, pt2, pt3 are collienar
def is_collinear(pt1, pt2, pt3):
	z1 = complex(pt1[0],pt1[1])   # convert to complex numbers
	z2 = complex(pt2[0],pt2[1])
	z3 = complex(pt3[0],pt3[1])
	return abs(((z1-z2)/(z3-z2)).imag) <= 1e-8 
	# test whether the imaginary part equals 0, up to numerical error 
\end{lstlisting}


Samples are shown below. We checked that the points $(1,2)$, $(2,3)$, and $(3,4)$ are collinear, but $(1,2)$, $(2,3)$, and $(3,3)$ are not collinear.
\begin{lstlisting}[language=Python, frame=single]
In [1]: is_collinear((1,2),(2,3),(3,4))
Out[1]: True

In [2]: is_collinear((1,2),(2,3),(3,3))
Out[2]: False	
\end{lstlisting}



\section*{Summary}

\begin{itemize}
	
\item A complex number $a+bi$ can be represented by $2\times 2$ matrix $\begin{bmatrix} a & -b \\ b & a\end{bmatrix}$. When the determinant $a^2+b^2$ is equal to~1, this is a rotation matrix.
	
\item Geometric meaning of the operations of complex numbers:	\begin{itemize} 
	\item Addition of complex numbers is vector addition.
	
	\item Multiplying by a complex number is the same as rotation followed by an expansion.
	
	\item Complex conjugate is the mirror image with respect to the real axis.
	
	\item The modulus of a complex number is the distance to the
	 origin.
	 
\end{itemize} 

    \item $\real(z) = (z+z^*)/2$, $\imag(z) = (z-z^*)/(2i)$, for all $z\in\mathbb{C}$.
    
    \item Complex conjugate respects addition, multiplication, and taking inverse: $(z_1+z_2)^*=  z_1^*+z_2^*$, $(z_1z_2)^* = z_1^*z_2^*$, $(1/z)^* = 1/(z^*)$.

	\item $|z|^2 = z z^*$ for all $z \in \mathbb{C}$.

	\item $|z_1 z_2| = |z_1|\,|z_2|$, for any $z_1, z_2 \in \mathbb{C}$.
		
	\item Triangle inequality: $|z_1+z_2| \leq |z_1|+|z_2|$ for all $z_1, z_2\in\mathbb{C}$.

	\item DeMoivre formula:
$
	(\cos\theta + i\sin\theta)^n = \cos (n\theta) + i\sin(n\theta).
$
for $n\in\mathbb{Z}$ and $\theta\in\mathbb{R}$.

\item A complex number $z$ is a called an $n$-th root of unity if $z^n=1$.
\end{itemize}


\section*{Exercises}

\renewcommand{\labelenumi}{\arabic{enumi}.}
\begin{enumerate}
	
	\item Write the following complex numbers in Cartesian form
\begin{flalign*}
(a)\  & (2-i)(3+2i), & (b)\ & \frac{(4+i)^2}{i^3},
& (c)\ & \frac{1+2i+3i^2}{(3+4i)^2},  &&&&&& \\
(d)\ & (1+i)^{1024}, & (e)\ &  \frac{1}{i + \frac{1}{i + \frac{1}{i+1}}}, & (f)\ & \big(1+\frac{i}{5}\big)^5. &&&&&&
\end{flalign*}
	
	
	\item Solve the following linear equations.
\begin{flalign*}
(a)\ & (2+i)z = (-1+i/2), & (b)\ & 3-iz = 2z-i. &&&&&
\end{flalign*}	


\item Explain what is wrong in the proof of $1=-1$:
$$
1 = \sqrt{1} = \sqrt{(-1)(-1)} = \sqrt{-1} \sqrt{-1} = i \cdot i = -1
$$


\item By writing a complex number as a pair $(x,y)$ and following the computation rules in Section~\ref{sec:complex_number_as_vector}, find all solutions to the following quadratic equation,
$$
z^2 + 4iz - 3 = 0.
$$


\item Let $\Arg(z)$ denote the principal argument of nonzero complex number $z$. Determine whether the following equations hold for all nonzero complex numbers $z_1$ and $z_2$. If yes, provide a proof, otherwise, give a counter-example.
\begin{enumerate}
	\item $\Arg(z_1) + \Arg(z_2) = \Arg(z_1z_2)$.
	\item $\Arg(z_1) - \Arg(z_2) = \Arg(z_1/z_2)$.
\end{enumerate}


\item Determine the value of $z$ that satisfies the system of equations:
$$
\begin{cases}
	2z + 4\bar{z} = -2,&  \\
	z - 2\bar{z} = 1 .& 
	\end{cases}
$$

\item Solve the following equations.
\begin{enumerate}
	\item $z^2 = -i$.
	\item $z^2 = 1+\sqrt{3} i$.
	\item $z^8 = 2-2i$.
\end{enumerate}


	\item Use DeMoivre's formula to derive
	\begin{enumerate}
		\item $\cos 3\theta = 4 \cos^3\theta - 3 \cos \theta$.
		\item $\sin 3\theta = -4 \sin^3\theta + 3 \sin \theta$.
		\item $\cos 4\theta =  8\cos^4\theta - 8 \cos^2\theta +1$.
        \item $\sin 4\theta =  4 \cos\theta (\sin \theta - 2  \sin^3\theta)$.
	\end{enumerate}

\item Suppose $z$ lie on the unit circle. Show that for any complex number $w$ that is not equal to $z$, we have
$$
\Big| \frac{w-z}{w \bar{z} - 1} \Big| = 1. 
$$

\item Prove the reverse triangle inequality for complex numbers:
$$
 \big| |z_1| - |z_2| \big| \leq |z_1-z_2|,
$$
for $z_1, z_2\in \mathbb{C}$.

\item  Show that the following identity holds for any nonzero complex number $z$ and positive integer $n$.
\begin{align*}
	&(z-z^{-1}) (z^{2n} + z^{2n-2}+ \cdots +z^{4} + z^{2}+1+z^{-2} + z^{-4}+ \cdots +z^{-2n+2}+ z^{-2n}) \\
	&= z^{2n+1}-z^{-2n-1}.
\end{align*}
Substitute $z$ by $\cos\theta+ i \sin \theta$ and show that
$$
\frac{\sin (2n+1)\theta}{\sin \theta} = 1 + 2\cos 2\theta + 2\cos 4\theta + \cdots + 2\cos 2n\theta.
$$

\item \begin{enumerate}
	\item Find all the square roots of $3i$.
	\item Find all the fourth roots of $-1+i$.
	\item Find all the sixth roots of unity.
\end{enumerate}

% Complex matrices
\item (Complex matrices) Many theorems in linear algebra can be formulated over the complex field. A major difference with the real vectors and matrices is that the transpose operator must be replaced by the Hermitian operator, and the dot product of two column vectors $\mathbf{v}$ and $\mathbf{u}$ with complex components is defined by $\mathbf{v}^H \cdot \mathbf{u}$.


{\bf Definition.} The {\em conjugate} of a vector or matrix with complex entries is computed by taking the complex conjugate of each entry in the vector or matrix. The
{\em Hermitian} of a matrix is defined as the conjugate of the transpose of the matrix. The Hermitian of a matrix $M$ is denoted by $M^H$.

For example, the Hermitian of $\begin{bmatrix}i & 2-i \\ -3 & 4+3i \end{bmatrix}$ is
$\begin{bmatrix}-i & -3 \\ 2+i & 4-3i \end{bmatrix}$.


{\bf Definition.} An $n\times n$ complex matrix $A$ is called 
\begin{enumerate}
	\item {\em Hermitian} if $A=A^H$, 
	\item {\em skew-Hermitian} if $A=-A^H$,
	\item {\em unitary} if $A\cdot A^H =I$,
\end{enumerate}
where $I$ is the $n\times n$ identity matrix.

Show that 
\begin{enumerate}
	\item the eigenvalues of Hermitian matrix is real, 
	\item the eigenvalues of skew-Hermitian matrix is purely imaginary (the eigenvalue is in the form $\alpha i$ for some $\alpha\in\mathbb{R}$), 
	\item the eigenvalues of unitary matrix has modulus 1.
\end{enumerate}

\item \label{exercise:DFT} (Discrete Fourier transform) Fix a positive integer $n$ and let $\omega$ be an $n$-th root of unity defined by
$$
\omega = \cos(2\pi/n)-i\sin(2\pi/n).
$$
The complex number $\omega$ has modulus 1 and argument $-2\pi/n$.
Define an $n\times n$ matrix
$$
W_n \triangleq  \begin{bmatrix}
	1  & 1 & 1 & 1& \cdots & 1 \\
	1 & \omega & \omega^2 & \omega^3 & \cdots & \omega^{n-1} \\
	1 & \omega^2 & \omega^4 & \omega^6 & \cdots & \omega^{2(n-1)} \\
	\vdots & \vdots & \vdots& \vdots & \ddots & \vdots \\
	1 & \omega^{n-1} & \omega^{2(n-1)} & \omega^{3(n-1)} & \cdots & \omega^{(n-1)^2} \\
\end{bmatrix}.
$$
The $(i,j)$-entry of $W_n$ is $\omega^{(i-1)(j-1)}$, for $i,j=1,2,\ldots, n$.


Given a column vector $\mathbf{v}$ of length $n$ (whose entries can be real or complex), the discrete Fourier transform (DFT) $\hat{\mathbf{v}}$ of $\mathbf{v}$ is computed by a matrix multiplication
$$
\hat{\mathbf{v}} \triangleq  W_n \mathbf{v}.
$$


\begin{enumerate}
	\item Write down the the matrix $W_n$ for $n=3$. Find the DFT of $\begin{bmatrix}1 \\ -1 \\ 1 \end{bmatrix}$.
	
	\item  Let $W_n^H$ denote the conjugate transpose of $W_n$. Show that  $$\frac{1}{n}W_n^H W_n = \frac{1}{n} I. $$
	
	Hint: for any integers $a$ and $b$ between 0 and $n-1$, show that
	$$
	\frac{1}{n} \sum_{k=0}^{n-1} {\omega^{ak}(\omega^{bk})^*} = \begin{cases}
		1 & \text{ if } a=b\\
		0 & \text{ if } a \neq b.
	\end{cases}
	$$
\end{enumerate} 


\item (Quaternions as $2\times 2$ complex matrices) In this exercise we construct the quaternions as complex matrices. Let $M_2(\mathbb{C})$ denote the set of $2\times 2$ matrices over the complex numbers. Consider the set
\[
\mathcal{A}
=
\Big\{
\begin{pmatrix}
	z & w \\
	-\bar{w} & \bar{z}
\end{pmatrix}
:\, z,w \in \mathbb{C}
\Big\}
\subseteq M_2(\mathbb{C}).
\]

\begin{enumerate}
	\item Show that $\mathcal{A}$ is closed under matrix addition and multiplication, and contains the identity matrix.  
	Conclude that $\mathcal{A}$ is a subalgebra of $M_2(\mathbb{C})$.
	
	\item Define the matrices
	\[
	I =
	\begin{pmatrix}
		1 & 0 \\ 0 & 1
	\end{pmatrix},\qquad
	\mathbf{i} =
	\begin{pmatrix}
		i & 0 \\ 0 & -i
	\end{pmatrix},\qquad
	\mathbf{j} =
	\begin{pmatrix}
		0 & 1 \\ -1 & 0
	\end{pmatrix},\qquad
	\mathbf{k} =
	\begin{pmatrix}
		0 & i \\ i & 0
	\end{pmatrix}.
	\]
	Verify that all four matrices lie in $\mathcal{A}$.
	
	\item Compute the products
	\[
	\mathbf{i}^2,\qquad \mathbf{j}^2,\qquad \mathbf{k}^2,
	\qquad
	\mathbf{i}\mathbf{j},\qquad \mathbf{j}\mathbf{i},
	\]
	and show that
	\[
	\mathbf{i}^2 = \mathbf{j}^2 = \mathbf{k}^2 = -I,
	\qquad
	\mathbf{i}\mathbf{j} = \mathbf{k},
	\qquad
	\mathbf{j}\mathbf{i} = -\mathbf{k}.
	\]
	
	\item Prove that every element of $\mathcal{A}$ can be written uniquely in the form
	\[
	aI + b\mathbf{i} + c\mathbf{j} + d\mathbf{k},
	\qquad a,b,c,d \in \mathbb{R}.
	\]
	
	\item Conclude that $\mathcal{A}$ is a $4$-dimensional real algebra with basis $I$, $\mathbf{i}$, $\mathbf{j}$, and $\mathbf{k}$
	satisfying the multiplication rules
	\[
	\mathbf{i}^2 = \mathbf{j}^2 = \mathbf{k}^2 = -I,
	\quad
	\mathbf{i}\mathbf{j} = \mathbf{k},\quad
	\mathbf{j}\mathbf{k} = \mathbf{i},\quad
	\mathbf{k}\mathbf{i} = \mathbf{j}, \quad
	\mathbf{j}\mathbf{i} = -\mathbf{k},\quad
\mathbf{k}\mathbf{j} = -\mathbf{i},\quad
\mathbf{i}\mathbf{k} = -\mathbf{j}.
	\]
\end{enumerate}

This algebra is (by definition) isomorphic to the quaternion algebra.



\item (Pirate and treasure chest) 
	A pirate arrived at an island and he found three trees on the island: an alma fig tree (A), a breadfruit tree (B), and a coconut palm (C). The pirate buried a treasure chest according to the following rule:
	
	\begin{enumerate}
		\item Starting at tree $A$, he walked to tree $B$, counting the number of steps.
		He then turned $90^\circ$ to the left and walked the same number of steps
		from tree $A$. He placed a stake in the ground at that point.
		
		\item Starting again at tree $A$, he walked to tree $C$, counting the number of steps.
		He then turned $90^\circ$ to the right and walked the same number of steps
		from tree $A$. He placed another stake at that point.
		
		\item The treasure is buried at the midpoint of the two stakes.
	\end{enumerate}
	
	The pirate then removed the two stakes and left the island.
	
	After some years, the pirate returned to the island and wants to dig up the treasure.
	However, tree $A$ was gone. Could the pirate still find the exact location of the treasure? 
	
	(Optional: Suppose we modify the rule as follows. Fix a positive real number $r$. In steps (a) (resp. step~b), after turning left (resp.~right), we walk for a distance of $r$ times the distance from tree $A$ to tree $B$ (resp.~$C$). Can we locate the treasure without knowing the location of tree $A$?)


\item (Medians of a triangle) Use complex numbers to show that the three medians of a triangle intersects at the same point.

\item (Two-lighthouse problem)
	Two lighthouses stand at complex positions $a$ and $b$.  
	A ship at point $z$ sees the angle between the two lighthouses as exactly $60^\circ$.
	
	\begin{enumerate}
		\item Show that the condition can be written as
		\[
		\arg\!\left(\frac{z-a}{z-b}\right)=\pm \frac{\pi}{3}.
		\]
		\item Deduce that the locus of such points $z$ consists of two circular arcs.
		\item Identify the centers and radii of these circles.
	\end{enumerate}
	
\item (Napoleon's triangle via complex numbers)
	Given a triangle $ABC$, construct equilateral triangles externally on each side.  
	Let their outer vertices be $A'$, $B'$, and $C'$.
	
	\begin{enumerate}
		\item Using complex numbers, express $A'$, $B'$, and $C'$ in terms of $A,B,C$.
		\item Show that triangle $A'B'C'$ is equilateral.
	\end{enumerate}

\item (Area of triangle)  \begin{enumerate}
	\item  Consider a triangle in the complex plane with vertices 0, $z_1$ and $z_2$. Prove that the area of the triangle is equal to $\frac{1}{2} |\imag(z_2 \bar{z}_1)|$

	\item More generally, let $z_1$, $z_2$ and $z_3$ be the vertices of a triangle in the complex plane. Suppose that the paths $z_1 \rightarrow z_2 \rightarrow z_3 \rightarrow z_1$ is in the counter-clockwise direction. Prove that the area of the triangle can be computed by the determinant
$$\frac{i}{4}
\begin{vmatrix}
	1 & 1 & 1 \\
	z_1& z_2 & z_3 \\
	\bar{z}_1 & \bar{z}_2 & \bar{z}_3
\end{vmatrix}.
$$


\end{enumerate}


\item Implement the algorithm in Section~\ref{sec:cubic_equation} to calculate the roots of a cubic equation with complex coefficients.

\end{enumerate}



%%%%%%%%%%%%%%%%%%%%%%%%%%%%%%%



\chapter{Linear Fractional Transformation and Riemann sphere}




\section{Complex inversion}

In Euclidean geometry, {\em circle inversion} (a.k.a.\ geometric inversion or circular inversion) is the transformation
$$(x,y) \mapsto (\frac{x}{x^2+y^2},\frac{y}{x^2+y^2}).
$$
The origin is excluded from the domain of this mapping because $x^2+y^2$ has no inverse when $x=y=0$. 
Let the inversion be taken in the circle with center $O$ and radius $R$.  
For a point $P\neq O$, its inverse of $P$ is defined as the point $P'$ such that
\[
|OP|\cdot |OP'| = R^{2}.
\]
In particular, every point on the circle is mapped to itself under the circle inversion.  In terms of polar coordinates, the point $z = r(\cos\theta + i\sin\theta)$ is mapped to the point $(1/r)(\cos\theta + i\sin\theta)$. Although this transformation is nonlinear in coordinates, the point $P'$ can be constructed using only straightedge and compass.





\begin{figure}
	\begin{center}
		
		\begin{tikzpicture}[scale=1.2]
			
			% Inversion circle
			\coordinate (O) at (0,0);
			\def\R{2}
			\draw[thick, blue ] (O) circle (\R);
			\fill (O) circle (1.5pt) node[below left] {$O$};
			
			% Point P outside the circle
			\coordinate (P) at (3,1);
			\fill [red] (P) circle (1.5pt) node[right] {$P$};
			
			% Inverse point P' on OP such that OP * OP' = R^2
			% Compute P' explicitly for the picture
			\coordinate (Pprime) at ($(O)!{(\R*\R)/( (3)^2 + (1)^2 )}!(P)$);
			\fill [red] (Pprime) circle (1.5pt) node[above left] {$P'$};
			
			% Draw line OP
			\draw[thick] (O) -- (P);
			
			% Label R
			\draw[dotted] (O) -- ++(\R,0);
			
			% Label r
			\node[below] at (1,0) {$R$};	
		\end{tikzpicture}
	\end{center}
	\caption{Circle inversion. The point $P$ outside is mapped to the point $P'$ inside the circle.}
	\label{fig:circle_inversion}
\end{figure}


\index{Inversion function}
The circle inversion is closely related to complex division. Geometrically, the complex inversion function $f(z) = 1/z$ can be interpreted as a circle inversion function followed by a reflection along the $x$-axis.  We can take the complex conjugate of $(1/r)(\cos\theta + i\sin\theta)$ to obtain
$$
(1/r)(\cos\theta - i \sin \theta) = 
\frac{1}{r(\cos\theta + i \sin \theta)} =
\frac{1}{z}.
$$
Hence, the function $z\mapsto 1/z$ and the circle inversion with radius $R=1$ differ only by a mirror reflection. We will see that the complex inversion map preserves orientation, but the circle inversion reverses orientation.

\begin{figure}
	\begin{center}
		
		\begin{tikzpicture}[scale=1.2]
			
			% Inversion circle
			\coordinate (O) at (0,0);
			\def\R{1}
			\draw[thick, blue ] (O) circle (\R);
			\fill (O) circle (1.5pt) node[below left] {$O$};
			
			% Point P outside the circle
			\coordinate (P) at (2,1.5);
			\fill [red] (P) circle (1.5pt) node[right] {$z$};
			
			% Inverse point P' on OP such that OP * OP' = R^2
			% Compute P' explicitly for the picture
			\coordinate (Q) at (2,-1.5);	
			\coordinate (Pprime) at ($(O)!{(\R*\R)/( (2)^2 + (1.5)^2 )}!(Q)$);
			\fill [red] (Pprime) circle (1.5pt) node[below] {$z^{-1}$};
			
			% Draw line OP
			\draw[thick] (O) -- (P);
			
			% Label R
			\draw[dotted] (O) -- ++(\R,0);
			\node[right] at (1,0) {$1$};	
		\end{tikzpicture}
	\end{center}
	\caption{Complex inversion. The point $z$ is mapped to the point $z^{-1}$.}
	\label{fig:complex_inversion}
\end{figure}



The inversion function has a very special property: it transforms circles and straight lines to circles and straight lines.  
Consider a quadratic equation in two variables,
\begin{equation}
	A(x^2+y^2)+Bx+Cy +D = 0,
	\label{eq:cline}
\end{equation}
where $A$, $B$, $C$ and $D$ are real numbers such that $A\neq 0$. 
By re-writing it as
\begin{align*}
	x^2 + y^2 + \frac{B}{A}x + \frac{C}{A}y &= -\frac{D}{A}\\
	(x + \frac{B}{2A})^2 + (y + \frac{C}{2A})^2 &= \frac{B^2 + C^2-4DA}{4A^2},
\end{align*}
it is the equation of a circle if $(B^2+C^2-4DA)>0$.

When $A$ is zero but $B$ or $C$ is nonzero, \eqref{eq:cline} is an equation of a straight line. When $A$, $B$, and $C$ are all zero but $D$ is nonzero, the solution set is empty. Finally, when all coefficients are zero, the the solution set is the whole $x$-$y$ plane.
Hence, in general, the equation in \eqref{eq:cline} defines a circle, a straight line, an empty set, or the plane.




\begin{svgraybox}
	\begin{theorem}
		The inversion function $f(z)=1/z$ transforms circles and straight lines to circles and straight lines. \label{thm:circle_straight_line1}
	\end{theorem}
\end{svgraybox}

\begin{proof} Suppose the coefficients in \eqref{eq:cline} are not all zero, so that the solution set is  a circle, a straight line or empty set.
	Represent $x$ and $y$ by complex variable $z=x+iy$, and re-write \eqref{eq:cline} as
$$	Azz^* + B \frac{z+z^*}{2} + C \frac{z-z^*}{2i} + D	=0 ,
$$
and simplify it to
	\begin{equation}
			a z z^* + b z + b^* z^* + c  = 0,
		\label{eq:cirline_proof}
	\end{equation}
	where $a=A$, $b=(B-iC)/2$, and $c=D$. The coefficients $a$ and $c$ are real numbers, and the coefficient $b$ is a complex number.
	If we map $z$ to $w=1/z$, the equation becomes
	$$
	a \frac{1}{w w^*} + b \frac{1}{w} + b^* \frac{1}{w^*} + c = 0,
	$$
	which is the same as
	$$
	a + b w^* + b^* w + c w w^*= 0.
	$$
	This has the same form as in \eqref{eq:cirline_proof}, and hence it represents a circle or a straight line.
\end{proof}




Using complex numbers, We can specify a circle in the complex plane using the modulus. Given a point $z_0$ in the complex plane and a positive real number, the circle with radius $r$ and center $z_0$ has equation
$$
|z - z_0| = r.
$$
Using the fact that $|z|^2 = zz^*$, we can write the equation in an alternate form:
$$
(z-z_0)(z-z_0)^*  = r^2.
$$
If we write $z=x+iy$ and $z_0 = x_0+iy_0$, the above equation becomes the usual equation of circle:
$$
(x - x_0)^2 + (y-y_0)^2 = r^2.
$$

A straight line in the $x$-$y$ plane can be represented by equation
\begin{equation}
	ax +by = c,
	\label{eq:straight_line1}
\end{equation}
for some real numbers $a$, $b$ and $c$.  Using the relation $x=\real(z) = (z+z^*)/2$ and $y=\imag(z) = (z-z^*)/(2i)$ in \eqref{eq:real_imaginary_conjugate}, we can write this equation as
$$
a \frac{z+z^*}{2} + b \frac{z-z^*}{2i} = c.
$$
We can simplify this equation to
\begin{align}
	ai(z+z^*) + b(z-z^*) & = 2ic \notag \\
	(b+ ai)z - (b -ai)z^* &= 2ic. \label{eq:straight_line_proof}
\end{align}
If we define $\alpha=b+ai$, then we can further simplify it to
\begin{equation}
	\imag(\alpha z) = c. \label{eq:straight_line_1}
\end{equation}
%In words,  only the points on the straight line satisfies $\imag(\alpha z)=c$.

Alternately, we can write the equation of a straight line using $\real$ instead of $\imag$. We let $\beta=a-ib$ in \eqref{eq:straight_line_proof}, we obtain
\begin{equation}
	\real(\beta z) = c. \label{eq:straight_line_2}
\end{equation}
A straight line can thus be \eqref{eq:straight_line1} can be represented by \eqref{eq:straight_line_1} or \eqref{eq:straight_line_2}.


\begin{example}
By identifying a point $(x,y)$ with complex variable $z$, we can represent the straight line $x=y$ by
$$
\real(\frac{1+i}{2}z) = 0.
$$
In words, this equation says that if we rotate the points on the straight line $x=y$ by $45^\circ$, we obtain the vertical line passing through the origin. Because the denominator 2 is a real constant, we can simplify the equation to
$$
\real((1+i)z) = 0.
$$
\end{example}




\begin{example}
	Find the image of the circle $|z-2i|=1$ under the transformation $f(z)=1/z$.
	
	Let $w=1/z$.
	\begin{align*}
		\Big| \frac{1}{w} -2i\Big|  &= 1 \\
		\Leftrightarrow \ \ |1-2wi|^2 &= |w|^2 \\
		\Leftrightarrow \ \ 3 |w|^2 -2wi +2 w^*i &= -1.
	\end{align*}
	By completing square, the last equation is equivalent to
	$$
	\big| w + \frac{2i}{3}\big|^2 = \frac{1}{9},
	$$
	which is the equation of a circle with radius $1/3$ and center $-2i/3$. (See Fig.~\ref{fig:circle_inversion_transformation1})
\end{example}

\begin{figure}
	\begin{center}
\begin{tikzpicture}[>=Stealth, line cap=round, line join=round]
	
	% Common axis style
	\tikzset{
		axes/.style={blue!70!black, thick, -Stealth},
		lab/.style={font=\small},
	}
	
	%----------------------------
	% Left: z-plane
	%----------------------------
	\begin{scope}[xshift=0cm, yshift=0cm]
		% axes
		\draw[axes] (-2.3,0) -- (2.6,0);
		\draw[axes] (0,-2.3) -- (0,4.2);
		
		% circle: center 2i, radius 1
		\draw[red, thick] (0,2) circle (1);
		
		% points and labels
		\fill[red] (0,2) circle (1.2pt);
		\fill[red] (0,1) circle (1.2pt);
		\fill[red] (0,3) circle (1.2pt);
		
		\node[lab, above right] at (0,2) {$2i$};
		\node[lab, right]       at (0,1) {$i$};
		\node[lab, above right] at (0,3) {$3i$};
		
		\node[lab, below=10pt] at (0,-2.3) {$z$-plane};
	\end{scope}
	
	%----------------------------
	% Middle: mapping arrow + label
	%----------------------------
	\draw[thick, -Stealth] (3.2,2.3) .. controls (4.2,2.7) and (5.2,2.7) .. (6.2,2.3);
	\node[lab] at (4.7,3.1) {$f(z)=\dfrac{1}{z}=w$};
	
	%----------------------------
	% Right: w-plane
	%----------------------------
	\begin{scope}[xshift=9.0cm, yshift=0cm]
		% axes
		\draw[axes] (-2.6,0) -- (2.6,0);
		\draw[axes] (0,-2.6) -- (0,3.6);
		
		% image circle (center -2i/3, radius 1/3)
		\draw[green!60!black, thick] (0,-2/3) circle (1/3);
		
		% points
		\fill[green!60!black] (0,-2/3) circle (1.1pt);
		\fill[green!60!black] (0,-1)   circle (1.1pt);
		
		% labels with small callout arrows (to reduce crowding)
		\draw[green!60!black, -Stealth] (0.55,-1.55) -- (0.06,-1.00);
		\node[lab, green!50!black] at (0.85,-1.50) {$-i$};
		
		\draw[green!60!black, -Stealth] (0.85,-0.6666) -- (0.06,-0.6666);
		\node[lab, green!50!black] at (1.25,-0.6666) {$-\dfrac{2}{3}i$};
		
		\node[lab, below=10pt] at (0,-2.6) {$w$-plane};
	\end{scope}
	
\end{tikzpicture}
\caption{The image of a circle under the inversion transformation.}
\label{fig:circle_inversion_transformation1}
	\end{center}
\end{figure}



\begin{example}
	Equation $|z-z_0|=|z_0|$ describes a circle that passes through the origin. Find the image of this circle under the transformation $f(z)=1/z$. 
	
	Define a variable $w=1/z$.
	\begin{align*}
		\Big| \frac{1}{w} - z_0 \Big| &= |z_0| \\
		\Leftrightarrow \ \ |1-wz_0|^2 &= |wz_0|^2 \\
		\Leftrightarrow \ \ 1 - wz_0 - {w}^* {z_0}^* & = 0 \\
		\Leftrightarrow \real(wz_0) &= 1/2.
	\end{align*}
	This equation defines a straight line in the $w$-plane. Indeed, if we write $w=u+iv$ and $z_0 = x_0+iy_0$, then
	$$
	\real(wz_0) = 1/2\ \  \Longleftrightarrow\ \  ux_0-vy_0=0.5.
	$$
	Hence, the image of the circle $|z-z_0|=|z_0|$ is a straight line with equation $ux_0-vy_0=0.5$.
\end{example}

\begin{figure}
	\begin{center}

\begin{tikzpicture}[>=Stealth, line cap=round, line join=round, scale=0.8]
	\tikzset{
		axes/.style={blue!70!black, thick, -Stealth},
		lab/.style={font=\small},
	}
	
	%----------------------------
	% Parameter (visual scale only)
	%----------------------------
	\def\z0{1.8} % controls circle center/radius in the z-plane (drawing units)
	
	%----------------------------
	% Left: z-plane
	%----------------------------
	\begin{scope}[xshift=0cm]
		% axes
		\draw[axes] (-2.6,0) -- (4.2,0);
		\draw[axes] (0,-2.6) -- (0,2.8);
		
		% circle centered at z0 on real axis with radius z0 (passes through origin)
		\draw[red, thick] (\z0,0) circle (\z0);
		
		% mark origin and center
		\fill[black] (0,0) circle (1.2pt);
		\fill[black] (\z0,0) circle (1.2pt);
		
		% labels
		\node[lab, below] at (\z0,0) {$z_0$};
		
		\node[lab, below=10pt] at (0,-2.6) {$z$-plane};
	\end{scope}
	
	%----------------------------
	% Middle: mapping arrow + label
	%----------------------------
	\draw[thick, -Stealth] (3.8,1.6) .. controls (4.9,2.1) and (6.0,2.1) .. (7.1,1.6);
	\node[lab] at (5.45,2.45) {$f(z)=\dfrac{1}{z}=w$};
	
	%----------------------------
	% Right: w-plane
	%----------------------------
	\begin{scope}[xshift=10.0cm]
		% axes
		\draw[axes] (-2.6,0) -- (4.2,0);
		\draw[axes] (0,-2.6) -- (0,2.8);
		
		% vertical line (schematic placement)
		\def\xa{1.7}
		\draw[green!60!black, thick] (\xa,-2.4) -- (\xa,2.6);
		
		% x-intercept point and label (no arrow)
		\fill[black] (\xa,0) circle (1.2pt);
		\node[lab] at (\xa+0.45,0.55) {$\dfrac{1}{2z_0}$};
		
		\node[lab, below=10pt] at (0,-2.6) {$w$-plane};
	\end{scope}
	
\end{tikzpicture}		
	\end{center}
\caption{Mapping a circle to a straight line.}
\label{fig:circle_inversion_transformation2}
\end{figure}



\section{Spherical representation of complex numbers}


The complex inversion $z\mapsto z^{-1}$ function is not defined when $z=0$. However, it is tempting to claim that ``$1/0$ is infinity''. 
We extend the domain of the complex inverse function by the following procedure.

%We identify a complex number $x+iy$ on the horizontal plane in a 3D space by $x+iy \mapsto (x,y,0)$. Write a general point in the 3D space as $(\alpha,\beta, \gamma)$ (See Fig.~\ref{fig:stereographic_projection}). Place a sphere $S$ of radius $r$ centered at $(0,0,r)$ in the 3D space. The equation of $S$ is
%\begin{equation}
%	S:\ \ \alpha^2 + \beta^2 + (\gamma - r)^2 = r^2.
%	\label{eq:sterographic_projection_proof1}
%\end{equation}
%The sphere $S$ is called the Riemann sphere. 
%We call $N = (0,0,2r)$ the north pole of  $S$. Given any point $(x,y,0)$, let $P(x,y)$ be the point of $S$ such that $(0,0,2r)$, $(x,y,0)$ and $P(x,y)$ are collinear. 

\index{Stereographic projection}
We identify a complex number $x+iy$ with the point $(x,y,0)$ in $\mathbb{R}^3$. Consider the sphere 
\begin{equation}
	S = \{(x,y,z)\in\mathbb{R}^3 :\, \alpha^2 + \beta^2 + (\gamma - r)^2 = r^2\}
	\label{eq:sterographic_projection_proof1}
\end{equation}
which has radius $r$ and is tangent to the $x$-$y$ plane at the origin. The \df{stereographic projection} $P(x,y)$ maps a complex number $x+iy$  to the point $P(x,y)$ on this sphere obtained by drawing the line through $(x,y,0)$ and the north pole $(0,0,2r)$. The point $P(x,y)$ is defined to be the intersection of this line with the sphere. The mapping is one-to-one; a proof is given in Section~\ref{sec:stereographic_projection}.

\begin{svgraybox}
	\begin{definition} \index{Riemann sphere} \index{Point at infinity} The north pole of the sphere is called the \df{point at infinity},  denoted by the symbol~$\infty$. The sphere itself is referred to as the \df{Riemann sphere}.
\end{definition}
\end{svgraybox}

Motivated by the stereographic projection, we enlarge the complex number system by adjoining a single extra point, denoted by $\infty$, which does not belong to $\mathbb{C}$. After adding this point, the resulting set becomes bijection with the Riemann sphere under stereographic projection.



\begin{svgraybox}
\begin{definition} \index{Extended complex number system} \index{Point at infinity}
		The \df{extended complex number system} is the set $\mathbb{C}\cup\{ \infty\}$, also called the \df{extended complex plane}.  We will denote it by $\bar{\mathbb{C}}$. Other common notation include includes $\hat{\mathbb{C}}$ and $\mathbb{C}_\infty$.	A number $z$ in the extended number plane is said to be $\df{finite}$ if  $z\neq \infty$.
	\end{definition}
\end{svgraybox}


\begin{figure}
\begin{center}
\begin{tikzpicture}[scale=1, line cap=round, line join=round]
	
	%----------------------------
	% Parameters
	%----------------------------
	\def\R{1.65}                   % sphere radius (drawn)
	\def\eq{0.35}                  % equator "flattening" (smaller -> more 3D look)
	\coordinate (C) at (0,1.8);    % sphere center (2D placement)
	\coordinate (X) at (6.6,-1.6); % point (x,y) on plane (2D placement)
	
	%----------------------------
	% Plane (more perspective)
	%----------------------------
	\coordinate (A) at (-6.0,-2.6);
	\coordinate (B) at ( 5.2,-4.2);
	\coordinate (D) at (-2.2, 1.7);
	\coordinate (E) at ( 9, 0.6);
	
	\fill[gray!10] (A)--(B)--(E)--(D)--cycle;
	\draw[gray!60] (A)--(B)--(E)--(D)--cycle;
	
	% a couple of faint "grid" lines to enhance 3D
	\draw[gray!30] ($(A)!0.33!(B)$) -- ($(D)!0.33!(E)$);
	\draw[gray!30] ($(A)!0.66!(B)$) -- ($(D)!0.66!(E)$);
	\draw[gray!30] ($(A)!0.35!(D)$) -- ($(B)!0.35!(E)$);
	\draw[gray!30] ($(A)!0.70!(D)$) -- ($(B)!0.70!(E)$);
	
	%----------------------------
	% Sphere: shadow + shaded ball + outline + equator
	%----------------------------
	\fill[black!20, opacity=0.25] ($(C)+(0,-\R-0.35)$) ellipse ({1.05*\R} and {0.28*\R});
	
	\shade[ball color=white] (C) circle (\R);
	\path[name path=sphere] (C) circle (\R);
	\draw[black] (C) circle (\R);
	
	% equator (front solid, back dashed)
	\draw[gray!55]       ($(C)+(\R,0)$) arc (0:-180:{\R} and {\eq*\R});
	\draw[gray!55,dashed]($(C)+(\R,0)$) arc (0: 180:{\R} and {\eq*\R});
	
	% North and south poles
	\coordinate (N) at ($(C)+(0,\R)$);
	\coordinate (O) at ($(C)+(0,-\R)$);
	
	%----------------------------
	% Projection line N -> (x,y), intersect sphere at P
	%----------------------------
	\path[name path=projline] (N) -- (X);
	\path[name intersections={of=sphere and projline, by={P,Pb}}];
	
	%----------------------------
	% Draw triangle/segments (blue)
	%----------------------------
	\draw[blue!70!black, thick] (O) -- (N);
	\draw[blue!70!black, thick] (O) -- (X);
	\draw[blue!70!black, thick] (N) -- (X);
	
	%----------------------------
	% Mark + labels
	%----------------------------
	\fill[red] (N) circle (1.6pt);
	\fill[red] (O) circle (1.6pt);
	\fill[red] (P) circle (1.6pt);
	\fill[blue!70!black] (X) circle (1.4pt);
	
	\node[red, above left=2pt] at (N) {$N=(0,0,2r)$};
	\node[red, below left=2pt] at (O) {$O$};
	\node[red, right=2pt]      at (P) {$P(x,y)=(\alpha,\beta,\gamma)$};
	\node[blue!70!black, below right=2pt] at (X) {$(x,y)$};
	
\end{tikzpicture}
\end{center}

	\caption{Riemann sphere and stereographic projection}
\label{fig:stereographic_projection}
\end{figure}





There is no easy way to perform addition and multiplication on the Riemann sphere. The main purpose of introducing the Riemann sphere is to bring in the point of infinity, which can be treated as an ordinary
point like the other complex numbers, and making the resulting Riemann sphere a compact set. Hence, the Riemann sphere is often called the {\em one-point compactification} of the complex plane.



\begin{remark}
	In contrast to the extended complex number system, the {\em extended real number system} is obtained by adding two special points to the real number system, i.e., $\mathbb{R}\cup\{\infty, -\infty\}$. In the real case, we have both $+\infty$ and $-\infty$ due to the total ordering of the real number line. However, since there is no ordering in the complex case, we only have one point at infinity.
\end{remark}







\section{Linear fractional transformations}

We first treat the special case of affine transformations, which can be classified into three types.

A \df{translation} is a function $f:\mathbb{C}\rightarrow \mathbb{C}$ in the form
$$
f_w(z) = z + w
$$
for some complex constant $w$. Geometrically it is just moving horizontally by $\real(w)$ and vertically by $\imag(w)$.

A \df{scaling function} is simply expanding by a constant real factor. It can be written as
$$
g_a(z) = az
$$
for some real constant $a$. This is also called  a \df{dilation}.

A \df{rotation} by angle $\theta$ can be realized by multiplication by $(\cos\theta + i \sin\theta)$,
$$
h_\theta(z) = (\cos\theta + i \sin\theta)z.
$$
%We can understand rotation in matrix form. If the real and imaginary part of $z$ are $x$ and $y$, respectively, the R.H.S. in the above equation is the same as a matrix multiplication,
%$$
%\begin{bmatrix}
%	\cos\theta & -\sin\theta \\
%	\sin\theta & \cos\theta
%\end{bmatrix}
%\begin{bmatrix} x &-y \\ y & x \end{bmatrix}
%$$
%The $2\times 2$ matrix on the left is a rotation matrix. Multiplication by this matrix represents an action of rotation by $\theta$ radian counter-clockwise.

All three types of mappings preserve the shape of a geometric object in the domain. For example, if we draw a triangle on the complex plane and apply a translation or a dilation to it, we will get a triangle that is similar to the original one. In particular, straight line is mapped to a straight line. If we have two intersecting lines, the angle between them is invariant under translation, scaling and rotation.

We define the class of \df{affine transformations}/\df{affine functions} as the set of functions that can be obtained by composing translation, scaling, dilation and rotation.
A general form of an affine function is
$$
z \mapsto az+b,
$$
where $b$ is a complex constant and $a$ is a nonzero complex number. For example, when $a=1$, we obtain a translation. When $b=0$, the function $f(z)=az$ is a scaling followed by a rotation. This class of affine functions form a set that is closed under composition.


Moreover, the inverse function of an affine function is also affine. In algebraic terms, the affine functions form a group under composition.


\index{Group}
In general, a \df{group} is a set of elements $G$ on which we can define a binary operation $\circ$, satisfying the following axioms:
\renewcommand{\labelenumi}{\arabic{enumi}.}
\begin{enumerate}
	\item (closedness) For all $g,h\in G$, $g\circ h\in G$.
	\item (associative) For all $f,g,h\in G$, $(f\circ g)\circ h = f \circ (g \circ h)$.
	\item (identity) There exists an element $e\in G$ such that $e\circ g = g\circ e = g$ for all $g\in G$.
	\item (inverse) For all $g\in G$, there exists an element $g'$ such that $g\circ g' = g' \circ g = e$.
\end{enumerate}

It is straightforward to check that the composition of two affine transformation is also an affine transformation. Given two affine functions $f(z)=az+b$ and $g(z)=a'z+b$, the composition $g(f(z))$ is equal to
$$
g(az+b) = a'(az+b)+b' = a'a z + a' b + b'.
$$
Since $a$ and $a'$ are both nonzero, the coefficient $a' a$ is also nonzero.

\smallskip 

It is obvious that a translation, dilation and rotation preserve shape. From Theorem~\ref{thm:circle_straight_line1}, we know that the inversion function also transforms circle/straight line to circle/straight line. 
This motivate the study the class of function obtained by mixing and composing affine functions and the inversion function $z\mapsto 1/z$. The resulting class of function is called {\em linear fractional transformation}. By construction, this class of complex functions has the following property.

\begin{svgraybox}
	\begin{theorem}
		A linear fractional transformation $f(z)=(az+b)/(cz+d)$ transforms circles and straight lines to circles and straight lines. \label{thm:circle_straight_line2}
	\end{theorem}
\end{svgraybox}



However, in contrast to affine transformations, the composition of functions is not well defined, as the inverse function is not defined at $z=0$. We cannot take $\mathbb{C}$ as the domain for a general linear fractional transformation. To remedy this problem, we may define $1/0$ as the point at infinity in the extended complex number system, and treat a linear fractional transformation as a function from $\bar{\mathbb{C}}$ to $\bar{\mathbb{C}}$.


We adopt the following convention for the computation with the point at infinity. 
\begin{itemize}
	\item For any extended complex number $z$, which may be a finite complex number or infinity, we define $z+\infty = \infty + z = \infty$.
	\item For any $b \in \bar{\mathbb{C}}\setminus\{0\}$, we define $b\cdot \infty = \infty \cdot b = \infty$
	\item For any $c\in \mathbb{C}\setminus\{0\}$, we define $c/\infty = 0$ and $c/0 = \infty$. 
	\item $\infty-\infty$ and $0\cdot \infty$ are undefined, and they will not occur in any calculations.
\end{itemize}

\index{Linear fractional transformation} \index{M\"{o}bius transformation}
We now know how to deal with the reciprocal of 0, and how to do arithmetic with $\infty$.
A \df{linear fractional transformation} (or \df{M\"obius transformation}, or \df{bilinear transformation}) has the following from:
		\begin{equation}
			f(z) = \frac{az+b}{cz+d}
			\label{eq:Mobius}
		\end{equation}
where $a$, $b$, $c$ and $d$ are complex numbers with $ad-bc\neq 0$. The number $ad-bc$ is called the {\em determinant}.


%We remark that even though a linear factional transformation is defined by four complex constants $a$, $b$, $c$, and $d$, actually there are only three degrees of freedom. We can see this by noting that, if we multiply all the coefficients $a$, $b$, $c$ and $d$ by a nonzero constant $k$,  the resulting expression induces the same function, 
%$$
%\frac{kaz+kb}{kcz+kd} = \frac{az+b}{cz+d}.
%$$
%


If $b=0$, $c=0$ and $d=1$, a linear fractional transformation reduces to a rotation followed by a dilation, i.e., $f(z)=az$.

If $a=1$, $c=0$ and $d=1$, a linear fractional transformation is a translation $f(z) = z+b$.

If $a=0$, $b=1$, $c=1$ and $d=0$, a linear fractional transformation becomes the inversion function $f(z)=1/z$.


We can verify that any function in the form in \eqref{eq:Mobius} can be expressed as a composition of translation, scaling, rotation, and inversion. If $c=0$ and $d\neq0$, then $(az+b)/d$ is just an affine transformation, and hence can be computed by first multiplying by $a/d$ and then adding the constant $b/d$. When $c\neq 0$, by applying a partial fraction expansion, we can see that a linear fractional transformation can be written as:
\begin{equation}
	\frac{az+b}{cz+d} = \frac{a}{c} + \frac{bc-ad}{c} \frac{1}{cz+d}.
	\label{eq:bilinear_injective}
\end{equation}
The condition $ad-bc\neq 0$ in the definition of linear fractional transformation ensures that it is a non-constant function. In fact, we may scale all four parameters $a$, $b$, $c$ and $d$ by a nonzero complex constant such that $ad-bc=1$.

This decomposition suggests that a proper definition of $f(\infty)$ is $a/c$. When $z$ has very large modulus, the modulus of the fraction $1/(cz+d)$ becomes negligible. The value  of $(az+b)/(cz+d)$ tends to $a/c$ when $z$ tends to the point at infinity.

The function induced by a linear fractional transformation in~\eqref{eq:Mobius} is given as follows.

\begin{svgraybox}
	\begin{definition}
		Given complex numbers $a$, $b$, $c$ and $d$ with $ad-bc\neq 0$, we define a function $f(z) = (az+b)/(cz+d)$ on the extended complex plane. 
		
		\noindent When $c=0$, 
		$$
		f(z) \triangleq  \begin{cases}
			(a/d)z + (b/d) & \text{ if } z \neq \infty,\\
			\infty & \text{ if } z = \infty.
		\end{cases}
		$$
		When $c\neq 0$, 
		$$
		f(z) \triangleq  \begin{cases} 
			\frac{az+b}{cz+d} & \text{ if } z\neq -d/c,\ z \neq \infty\\
			\infty & \text{ if } z = -d/c \\
			a/c & \text{ if } z = \infty.
		\end{cases}
		$$
		\label{def:LFT}
	\end{definition}
\end{svgraybox}

We record several remarks about this computational formula. When $c=0$, the condition $ad-bc\neq 0$ forces $a\neq 0$ and $d\neq 0$. In this case, the LFT reduces to an affine map $az+b$, which sends the point at infinity to itself.


When $c\neq 0$. There are two exceptional  input values: $z=-d/c$ and $z=\infty$.
At $z=-d/c$, the denominator $cz+d$ vanishes, so we define $f(-d/c)=\infty$. For any finite complex number $z \neq -d/c$,
the value $f(z)$ cannot equal $a/c$. Indeed, if
$$
\frac{az+b}{cz+d} = \frac{a}{c} ,
$$
then
$$
c(az+b) = a(cz+d)\ \Longleftrightarrow\ 
caz+cb = acz+ad\ \Longleftrightarrow \ 
 ad-bc=0,
$$
contradicting the assumption that $ad-bc\neq 0$. 

This observation justifies defining  $f(\infty) = a/c$. For all other finite values $z\neq -d/c$, we set $f(z)=(az+b)/(cz+d)$. With these conventions, $f$ becomes a bijection on the extended complex plane.


\begin{svgraybox}
	\begin{theorem}
Let $f(z)=(az+b)/(cz+d)$ be a linear fractional transformation. Then exactly one of the following holds: 

\begin{enumerate} \label{thm:LFT_decomposition}
	\item If $f(\infty) = \infty$, then $c=0$ and $f(z)$ is an affine transformation.
	\item If $f(\infty)$ is finite, then $c\neq 0$, and $f(z)$ can be expressed as a composition of a translation, a scaling, and an inversion. 
\end{enumerate}		
	\end{theorem}
\end{svgraybox}

\begin{proof}
 By the rule of evaluating an LFT in \eqref{def:LFT}, we have 
 $f(\infty)=\infty$ if and only if the coefficient $c=0$. When $c=0$, $f(z)$ reduces to an affine transformation. When $c\neq 0$, the decomposition in \eqref{eq:bilinear_injective} shows that $f(z)$ can be computed by first multiplying $z$ by  $c$, translate by $d$, then multiply by the $\frac{bc-ad}{c}$, and finally translate by $a/c$.
\end{proof}

Next, we study the group properties of linear fractional transformations.

\begin{svgraybox}
	\begin{theorem}
		The composition of two linear fractional transformations is a linear fractional transformation. 
	\end{theorem}
\end{svgraybox}


\begin{proof} (sketch)
	In the followings, $f(z)$ is a linear fractional transformation $f(z) = (az+b)/(cz+d)$. It suffices to consider three cases.
	
	\noindent (i) The function $f(z)+\alpha$ is a linear fractional transformation for any complex constant $\alpha$.
	
	\noindent (ii) The function $\beta f(z)$ is a linear fractional transformation for any nonzero complex constant $\beta$.
	
	\noindent (iii) The function $1/f(z)$ is a linear fractional transformation.
\end{proof}

We next consider the inverse of a linear fractional transformation. We can use the fact that a linear fractional transformation can be written as the composition of  four types of functions --- translation, rotation, dilation, and inversion. Since all of them are invertible function, we obtain


\begin{svgraybox}
	\begin{theorem} The inverse of a linear fractional function exists, and is a linear fractional transformation.
	\end{theorem}
\end{svgraybox}

We give a more direct proof  by computing the inverse of a linear fractional transformation, and showing that the inverse is a linear fractional transformation.

\begin{proof}   Let $f(z)$ be a linear fractional transformation defined by $f(z) = (az+b)/(cz+d)$. We consider two cases.
	
	Suppose that $c=0$. In this case, we must have $a\neq 0$ and $d\neq 0$ (because we assume $ad-bc\neq 0$). The function $f(z)$ can be written in the form $az+b$, by writing $a/d$ as $a$ and $b/d$ as $b$. The inverse of $f(z)$ is $g(w) = (w-b)/a$. We also have $f(\infty)=\infty$ in this case. The inverse of $f(z)$ is
	$$
	g(w) = \begin{cases}
		(w-b)/a & \text{ if } w \neq \infty, \\
		\infty & \text{ if } w = \infty.
	\end{cases}
	$$
	We can check that $g\circ f$ and $f\circ g$ are both equal to the identity function.
	
	Next suppose $c\neq 0$. Using~\eqref{eq:bilinear_injective}, we write the equation $f(z)=w$ as 
$$	
	 \frac{a}{c} + \frac{bc-ad}{c} \frac{1}{cz+d} = w,
$$	
and try to express $z$ in terms of $w$.  We first simplify it to

	$$
	\frac{cw-a}{bc-ad} = \frac{1}{cz+d}.
	$$
	If $w \neq a/c$, we can solve for $z$ in $\mathbb{C}$
	$$
	z = \frac{1}{c} \Big( \frac{bc-ad}{cw-a} - d \Big) = \frac{b-dw}{cw-a}.
	$$
	If $w = a/c$, we have $f(\infty) = w = a/c$. The inverse function of $f(z)$ is
	$$
	g(w) = \begin{cases}
		\frac{b-dw}{cw-a} & \text{ if } \infty \neq w \neq a/c \\
		\infty & \text{ if } w = a/c \\
		-d/c & \text{ if } w = \infty.
	\end{cases}
	$$
	Since  $g(f(z))$ and $f(g(w))$ are identity functions from Riemann sphere to itself, we conclude that $f$ is a bijection.
\end{proof}

In general, if $f(z) =  \frac{az+b}{cz+d}$ is a linear fractional transformation, the inverse exists and is given by
$$
f^{-1}(z) = \frac{d z - b}{-c z + a}.
$$


\begin{example} \label{ex:mapping_upperhalfplane_to_circle}
	Consider the function $f(z) = (z+1)/(z-i)$. When $z=i$, we define $f(i)=\infty$ and when $z=\infty$, we define $f(\infty)=1$. To compute the inverse function of $f(z)$, we solve for $z$ in the equation $w=f(z)$,
	\begin{align*}
		w & = \frac{z+1}{z-i} \\
		wz-iw &= z+1 \\
		(w-1)z &= 1+iw \\
		z &= \frac{1+iw}{w-1}.
	\end{align*}
	The last line requires $w\neq 1$. The inverse of $f(z)$ is thus
	$$
	g(w) = \begin{cases}
		\frac{1+iw}{w-1} & \text{ if } w\neq 1,\ w \neq \infty,\\
		\infty & \text{ if } w = 1,\\
		i & \text{ if } w = \infty.
	\end{cases}
	$$
\end{example}

Consider two LFTs $f(z)$ and $g(z)$ represented by $\frac{az+b}{cz+d}$ and $\frac{\alpha z + \beta}{\gamma z + \delta}$, respectively. The composition $g\circ f$ is the linear fractional transformation given by
\begin{align*}
g(f(z)) &= 
g\Big(\frac{az+b}{cz+d} \Big) \\
&=\frac{(\alpha a + \beta c)z + (\alpha b+ \beta d)}{(\gamma a + \delta c) z + (\gamma b + \delta d)}.
\end{align*}
We note the similarity to matrix multiplication 
$$
\begin{bmatrix} \alpha & \beta \\ \gamma & \delta \end{bmatrix}
\begin{bmatrix} a & b \\ c & d \end{bmatrix} =
\begin{bmatrix} \alpha a + \beta c & \alpha b + \beta d \\
	\gamma a + \delta c & \gamma b + \delta d \end{bmatrix}.
$$
We may represent a linear fractional transformation $\frac{az+b}{cz+d}$ by $2\times 2$ complex matrix $\begin{bmatrix} a & b \\ c & d \end{bmatrix}$, and ``define'' the composition of two LFT's by matrix multiplication. Since the determinant $ad-bc$ is nonzero by assumption, inverse of the matrix $\begin{bmatrix} a & b \\ c & d \end{bmatrix}$ exists. Moreover, matrix multiplication is associative, and identity exists. Hence, the set of nonsingular $2\times 2$ complex matrices form a group under matrix multiplication.



\index{Group action} Linear fractional transformations can be formulated naturally as a group action on the extended complex plane.  In general,
an \df{action} of a group $G$ with identity element $e$ on a set $X$ is a function $$\alpha: G\times X\rightarrow X$$ 
satisfying the axioms:
\begin{enumerate}
	\item $\alpha(e,x) = x$ for all $x\in X$.
	\item $\alpha(g\circ h, x) = \alpha(g, \alpha(h,x))$ for all $g,h\in G$ and $x\in X$.
\end{enumerate}


Using this terminology, we can state the main theorem of this section.

\begin{svgraybox}
	\begin{theorem}
Let $G$ be the group of nonsingular $2\times 2$ complex  matrix, with composition $\circ$ given by matrix multiplication. For a matrix $\begin{bmatrix}a & b \\ c & d \end{bmatrix}$ and a point $z$ in the extended complex number system, define $\alpha$ by the linear fractional transformation in Definition~\ref{def:LFT}. Then $\alpha$ defines a group action of $G$ on the extended complex number plane.
	\end{theorem}
\end{svgraybox}


A direct verification of this theorem requires checking several cases. A formal LEAN proof is available on \href{https://github.com/wkshum/ComplexVariables/tree/main/MAT3253/Chapter2}{this github page}.
A more conceptual approach is to view the extended complex numbers as the complex projective line; we outline this perspective in Section~\ref{sec:complex_projective_geometry}.
 

 
Since multiplying all entries of a matrix by a nonzero complex scalar does not change the associated linear fractional transformation, we obtain an equivalence relation on the set of $2\times 2$ nonsingular matrices: two matrices are equivalent if one is a nonzero scalar multiple of the other. Two linear fractional transformations represent the same function precisely if their corresponding matrices are equivalent.

\begin{remark}
The quotient fo the group of $2\times 2$ matrices by  this equivalence relation is the projective  general linear group $PGL(2,\mathbb{C})$.  The group $PGL(2,\mathbb{C})$ acts faithfully on the Riemann sphere. 
A discrete subgroup of $PGL(2,\mathbb{C})$ is called a {\em Kleinian group}. The next example illustrates a finite subgroup of linear fractional transformations.
\end{remark}



\begin{example}
	Consider the linear fractional transformations $f_1(z) = 1/z$ and $f_2(z) = 1-z$. Both of them induce a permutation on the three points $0,1,\infty$ on the Riemann sphere. 
	\begin{flalign*}
		&&	f_1(0) &= \infty, & f_1(1)&=1, & f_1(\infty)=0, &&&&
		\\
		&&	f_2(0) &= 1, & f_2(1)&=0, & f_2(\infty)=\infty. &&&&
	\end{flalign*}
	In this example we want to exhaust all function that can be obtained by composing $f_1$ and $f_2$ recursively. For example, using $f_1$ and $f_2$, we can obtain functions in the form $f_1(f_2(z))$, $f_2(f_1(z))$, $f_1(f_1(f_2(z)))$, $f_2(f_1(f_2(z)))$, etc. How many functions can we get?
	
	First, we see that we can get the identity function from 
	$$f_1(f_1(z))=f_2(f_2(z))=z.$$ 
	
	We can discover two new functions from
	\begin{align*}
		f_1(f_2(z)) &= \frac{1}{1-z} ,\\
		f_2(f_1(z)) &= 1 - \frac{1}{z} = \frac{z-1}{z} . 
	\end{align*} 
	If we apply the inversion function to $f_2(f_1(z))$, we get
	$$
	f_1(f_2(f_1(z))) = \frac{z}{z-1}.
	$$
	
	We claim that these six functions
	\begin{flalign*}
		&f_0(z) = z, 
		&&f_1(z) = 1/z, 
		&&f_2(z) = 1-z, 
		&&f_3(z) = \frac{1}{1-z},
		&&f_4(z) = \frac{z-1}{z},
		&&f_5(z) = \frac{z}{z-1}
	\end{flalign*}
	are all we can get. We can organize the relationship among them in a group composition table. 
%	$$
%	\begin{array}{c|cccccc} 
%		\circ & f_0 & f_1 & f_2 & f_3 & f_4 & f_5 \\ \hline
%		f_0 &    \\
%		f_1 &  & f_0 & f_3 && f_5 \\
%		f_2 &  & f_4 & f_0 &  \\
%		f_3 & \\
%		f_4 &  &\\
%		f_5 &  &
%	\end{array}
%	$$
%	The non-empty cells correspond to the equations we know so far: $f_0 = f_1\circ f_1= f_2\circ f_2$, $f_3 = f1\circ f_2$, $f_4 = f_2 \circ f_1$, $f_5 = f_1\circ f_4 = f_1 \circ f_2 \circ f_1$. 
%	
%	We complete the table, and note that the six functions are closed under composition.
	$$
	\begin{array}{c||c|c|c|c|c|c} 
		& f_0(z) = & f_1(z) = &  f_2(z) = &  f_3(z) = & f_4(z) = &  f_5(z) = \\ 
		\circ & z &  \frac{1}{z} & 1-z & \frac{1}{1-z} &  \frac{z-1}{z} & \frac{z}{z-1} \\ \hline \hline
		f_0(z) = z &  z &  \frac{1}{z} & 1-z & \frac{1}{1-z} & \frac{z-1}{z} & \frac{z}{z-1} \\ \hline
		f_1(z) = \frac{1}{z} & \frac{1}{z}  & z & \frac{1}{1-z} & 1-z& \frac{z}{z-1} & \frac{z-1}{z}  \\ \hline
		f_2(z) = 1-z & 1-z & \frac{z-1}{z} & z &  \frac{z}{z-1} & \frac{1}{z} & \frac{1}{1-z} \\ \hline
		f_3(z) = \frac{1}{1-z} & \frac{1}{1-z} & \frac{z}{z-1}  & \frac{1}{z} & \frac{z-1}{z} & z & 1-z\\ \hline 
		f_4(z) = \frac{z-1}{z} & \frac{z-1}{z} & 1-z & \frac{z}{z-1} & z & \frac{1}{1-z} & \frac{1}{z}  \\ \hline
		f_5(z) = \frac{z}{z-1} & \frac{z}{z-1}  & \frac{1}{1-z} &   \frac{z-1}{z} & \frac{1}{z}& 1 - z & z\\ 
	\end{array}
	$$
	In every column and every row, each of the six functions appears 
	once and only once. 
\end{example}

A notable feature in this example is that the group generated by $f_1(z)$ and $f_2(z)$ is finite. In fact, each LFT in this example is determined by its action on three special points 0, 1 and $\infty$.	The functions correspond to the six  permutations of  0, 1 and $\infty$ on the Riemann sphere. We will see in the next section that an LFT is determined by its action on three distinct points.




\section{3-Transitivity of linear fractional transformation}

In this section we consider the question: Given two circles (or straight lines) on the complex plane, can we fine a linear fractional transformation that maps the first circle to the second circle. The answer is yes, and we can establish a stronger result. It is known that a circle is uniquely determined by three points. We prove in the next theorem that we can find a linear fractional transformation that maps any three distinct points to any three other distinct points. Moreover, the LFT is unique up to multiplying all coefficients $a$, $b$, $c$ and $d$ by a nonzero constant.

We first establish a special property of linear fractional transformation. We will use the notation of ``fixed point.'' Given a function $f$, we say that a number $z_0\in\bar{\mathbb{C}}$ is a \df{fixed point} of $f$ if $f(z_0)= z_0$.

\begin{svgraybox}
	\begin{theorem}
		If linear fractional transformation $f$ satisfies $f(0)=0$, $f(1)=1$ and $f(\infty)=\infty$, then $f(z) = z$ for all $z\in\bar{\mathbb{C}}$. \label{thm:LFT_3_fixed_point}
	\end{theorem}
\end{svgraybox}

\begin{proof}
Since $f(\infty)=\infty$, we know that $f(z)$ is an affine transformation by Theorem~\ref{thm:LFT_decomposition}. We can find complex numbers $a$ and $b$ such that $f(z) = az+b$ for all $z\in\mathbb{C}$. Substituting $f(0)$ by 0 and $f(1)$ by 1, the only possibility is $a=1$ and $b=0$. This proves that $f(z) = z$ for all $z\in\bar{\mathbb{C}}$.	
\end{proof}


\begin{svgraybox}
	\begin{theorem}
		Given any three distinct points $z_1$, $z_2$ and $z_3$  all in the extended complex plane $\bar{\mathbb{C}}$, we can find a linear fractional transformation that maps $z_1$ to $0$, $z_2$ to $1$, and $z_3$ to $\infty$. Moreover, this linear fractional transformation is unique.
		\label{thm:LFT3points0}
	\end{theorem}
\end{svgraybox}

\begin{proof} We first assume that $z_1$, $z_2$ and $z_3$ are finite complex numbers. The function
	\begin{equation}
		f(z) = \frac{z - z_1}{z-z_3} \cdot \frac{z_2-z_3}{z_2-z_1}
		\label{eq:LFT1}
	\end{equation}
	is a linear fractional transformation that satisfies the requirements, mapping $z_1$ to $0$, $z_2$ to $1$, and $z_3$ to $\infty$. 
	

	There are three boundary cases to consider. When $z_1=\infty$ and $z_2$ and $z_3$ are finite, we can map $\infty$ to 0, $z_2$ to 1, and $z_3$ to $\infty$ by
	\begin{equation}
	f(z) = \frac{z_2-z_3}{z-z_3}.
	\label{eq:LFT2}
	\end{equation}
	When $z_2=\infty$ and $z_1$ and $z_3$ are finite, we can map $z_1$ to 0, $\infty$ to 1, and $z_3$ to $\infty$ by
	\begin{equation}
	f(z) = \frac{z-z_1}{z-z_3}.
	\label{eq:LFT3}
	\end{equation}
	Finally, when $z_3=\infty$ and $z_1$ and $z_2$ are finite, we can map $z_1$ to 0, $z_2$ to 1, and $\infty$ to $\infty$ by
	\begin{equation}
	f(z) = \frac{z-z_1}{z_2-z_1}.
	\label{eq:LFT4}
	\end{equation}
	
	To prove uniqueness, suppose $g$ is another linear fractional transformation that satisfies the requirements in the theorem. Then $f\circ g^{-1}$ is a linear fractional transformation that fixed 0, 1 and $\infty$. By Theorem~\ref{thm:LFT_3_fixed_point}, $f\circ g^{-1}$ is the identity function. By the group property of linear fractional transformation, we have $g= f$.
\end{proof}



\begin{example}
	Find a linear fractional transformation that maps $0$ to $0$, $1$ to $1$, and $z_0$ to $\infty$, where $z_0$ is a complex number different from 0 and 1.
	
	Because we want to take $z_0$ to $\infty$, the linear fractional transformation has the form
	$$
	f(z)  = \frac{az+b}{z-z_0}
	$$
	for some choice of complex constants $a$ and $b$. Because we want $f(0)=0$, the only choice of constant $b$ is $b=0$. From the requirement $f(1)=1$, we derive that $a$ is equal to $1-z_0$. The answer is
	$$
	f(z) = (1-z_0)\frac{z}{z-z_0}.
	$$ 
	
	Consider the case that $z_0$ is equal to $i$. The three points 0, 1 and $i$ in the domain of the transformation lie on the circle with equation
	$|z - (1+i)/2|^2 = 1/2$.
	The three points 0, 1, $\infty$ in the range lie on the line $\imag(z)=0$ (the real axis plus the point at infinity). In this case, the fractional linear transformation
	$$
	z \mapsto \frac{(1-i)z}{z-i}
	$$ 
	maps the circle with and radius $1/\sqrt{2}$ and center at $(1+i)/2$ to the real axis.
\end{example}

%To prove the existence of linear fractional transformation in the theorem, we can consider a linear fractional transformation $f(z)$ that maps $z_1$ to 0, $z_2$ to 1, and $z_3$ to $\infty$, and another linear fractional transformation $g(w)$ that maps $w_1$ to 0, $w_2$ to 1, and $w_3$ to $\infty$. Then the composite function $g^{-1}\circ f$ is a linear fractional transformation that satisfies the required mapping properties. 



\begin{svgraybox}
	\begin{theorem}
		Given any three distinct points $z_1$, $z_2$ and $z_3$ and any three distinct points $w_1$, $w_2$ and $w_3$, all in the extended complex plane $\bar{\mathbb{C}}$, there is a unique linear fractional transformation that maps $z_1$ to $w_1$, $z_2$ to $w_2$, and $z_3$ to $w_3$. \label{thm:LFT3points}
	\end{theorem}
\end{svgraybox}

\begin{proof}
By Theorem~\ref{thm:LFT3points0}, there exists a unique LFT $f_1(z)$ that maps $z_1$, $z_2$ and $z_3$ to $0$, $1$, and $\infty$, and another unique LFT $f_2(z)$  that maps $w_1$, $w_2$ and $w_3$ to $0$, $1$, and $\infty$. The function $f_2^{-1}\circ f_1$ is the required LFT. 

To prove uniqueness, suppose $g(z)$ is an LFT that satisfies the conditions in the theorem. The composed function $f_2 \circ g \circ f_1^{-1}$ is an LFT that fixed 0, and 1 and $\infty$, and hence is the identity function. This implies that $g(z)$ must be equal to $f_2^{-1} \circ f_1$ by Theorem~\ref{thm:LFT_3_fixed_point}.
\end{proof}

\section{Cross ratio}

Given four distinct numbers $z_0$, $z_1$, $z_2$ and $z_3$ in $\bar{\mathbb{C}}$, we define the cross ratio as $f(z_0)$, where $f(z)$ is the unique function that maps $z_1$ to 0, $z_2$ to 1, and $z_3$ to $\infty$.

\begin{svgraybox}
	\begin{definition} \index{Cross ratio}
		Given four  complex numbers $z_0$, $z_1$, $z_2$ and $z_3$, we define the \df{cross ratio} as
		$$
		[z_0,z_1,z_2,z_3] \triangleq 
\begin{cases}		
	\frac{z_0 - z_1}{z_0-z_3} \cdot \frac{z_2-z_3}{z_2-z_1} & \text{ if } z_0 \neq \infty, \\
	\frac{z_2-z_3}{z_2-z_1} & \text{ if } z_0 = \infty,
\end{cases}
 		$$
 		when $z_1$, $z_2$ and $z_3$ are all finite. If one of $z_1$, $z_2$ and $z_3$ is $\infty$, the cross ratio is computed by \eqref{eq:LFT2}, \eqref{eq:LFT3} or \eqref{eq:LFT4}.
	\end{definition}
\end{svgraybox}
When $z_0$, $z_1$, $z_2$ and $z_3$ are not distinct, then the cross-ratio is not defined, or defined to be some junk value.


An important property of the cross ratio is that it is invariant under linear fractional transformation.

\begin{svgraybox}
	\begin{theorem}
		Let $f(z)$ be a linear fractional transformation mapping $z_k$ to $z_k'$, for $k=0,1,2,3$, then
		$$
		[z_0,z_1,z_2,z_3] = [z_0', z_1', z_2', z_3'].
		$$
	\end{theorem}
\end{svgraybox}

\begin{proof} (sketch)
	Since a linear fractional transformation can be obtained by composing an affine transformation and a complex inversion (see \eqref{eq:bilinear_injective}), we can study the effect of applying a translation, a scaling and an inversion to the cross ratio.
	
	
	Let $z_0$, $z_1$, $z_2$ and $z_3$ be four complex numbers. 
	If we add a complex number $\zeta$ to all of them, the cross ratio of the resulting points is
	\begin{align*}
		[z_0+\zeta, z_1+\zeta, z_2+\zeta, z_3+\zeta ]&=  \frac{z_0+\zeta - (z_1+\zeta)}{z_0+\zeta-(z_3+\zeta)} \cdot \frac{z_2+\zeta-(z_3+\zeta)}{z_2+\zeta-(z_1+\zeta)}\\
		&= \frac{z_0 - z_1}{z_0-z_3} \cdot \frac{z_2-z_3}{z_2-z_1} \\
		&= [z_0, z_1,z_2,z_3].
	\end{align*}
	By similar calculation, we can easily verify that, for any nonzero complex number $\alpha$,
	$$
	[\alpha z_0, \alpha z_1, \alpha z_2, \alpha z_3] = 
	[z_0, z_1,z_2,z_3].
	$$
	
	When $z_0$, $z_1$, $z_2$, and $z_3$ are all nonzero complex number in $\mathbb{C}$, we have
	\begin{align*}
		\big[ \frac{1}{z_0}, \frac{1}{z_1}, \frac{1}{z_2}, \frac{1}{z_3}\big] &=
		\frac{\frac{1}{z_0} - \frac{1}{z_1}}{\frac{1}{z_0}-\frac{1}{z_3}} \cdot \frac{\frac{1}{z_2}-\frac{1}{z_3}}{\frac{1}{z_2}-\frac{1}{z_1}} \\
		&=  \frac{z_1 - z_0}{z_3-z_0} \cdot \frac{z_3-z_2}{z_1-z_2} \\
		&= [z_0,z_1,z_2,z_3].
	\end{align*}
	There are few more cases to check when $z_0$, $z_1$, $z_2$ and $z_3$ can $\infty$ as the value. They are left as exercise.
\end{proof}


Suppose $f(z)$ is an LFT that map $z_k$ to $w_k$ for $k=1,2,3$. From Theorem~\ref{thm:LFT3points0} we know that such a function exists and is unique. We can compute $f(z)$ using of cross ratio. For any complex number $z$ not equal to $z_1$, $z_2$ and $z_k$, suppose $f(z) =w$. Since the cross ratio of $[z,z_1,z_2,z_3]$ is equal to $[w,w_1,w_2,w_3]$, we obtain
\begin{equation}
	\frac{z - z_1}{z-z_3} \cdot \frac{z_2-z_3}{z_2-z_1} =
	\frac{w - w_1}{w-w_3} \cdot \frac{w_2-w_3}{w_2-w_1}.
	\label{eq:LFT_equation}
\end{equation}
We can solve for $w$ as a function of $z$. This gives an explicit form for the function~$f(z)$.

\begin{example}
	We construct a linear fractional transformation that maps three points $z_1=1$, $z_2=2$, and $z_3=3$ on the real axis to three points $w_1 = i$, $w_2=3i$, and $w_3=5i$ on the imaginary axis.
	%, where $a_1$, $a_2$ and $a_3$ are distinct real numbers. When $a_1=1$, $a_2=2$ and $a_3=3$, we can simply select the linear fractional transformation that corresponds to rotation by $\pi/2$. For general $a_1$, $a_2$ and $a_3$,
	We  apply the equation in~\eqref{eq:LFT_equation},
	$$
	\frac{z-1}{z-3}\cdot \frac{2-3}{2-1} = \frac{w-i}{w-5i} \cdot \frac{3i-5i}{3i-i}.
	$$
	After some algebraic calculations, we can simplify it to
	$$
	w = 2iz-i.
	$$
	The desired linear fractional transformation is
	$$
	f(z) = 2iz - i = i(2z-1).
	$$
\end{example}



Using the fact that linear fractional transformation maps circle/straight line to circle/straight line, we have a concise condition for four points lying on a circle. We say that a set of points are \df{concyclic} if they are on the same circle. 

Recall that $[z,z_1,z_2,z_3]$ is a linear fractional transformation that maps $z_1$ to 0, $z_2$ to 1, and $z_3$ to the point of infinity. Since 0, 1 and the point at infinity lies on a straight line, namely the real axis, any point $z_0$ that is concyclic with $z_1$, $z_2$ and $z_3$ must also be mapped to the real axis by the same fractional linear transformation. Hence, 
\begin{equation}
	[z_0,z_1,z_2,z_3] \text{ is real iff $z_0$, $z_1$, $z_2$ and $z_3$ are concyclic.}
	\label{eq:concyclic} 
\end{equation}

This gives a convenient test to determine whether four points are located on a circle.

\begin{example}
	Derive the equation of the circle that passes through $(0,0)$, $(1,0)$ and $(0,2)$ in the $x$-$y$ plane.
	
	We take $z_1=0$, $z_2=1$ and $z_3=2i$. Let $z=x+iy$ be a point that lie on the circle that contains $z_1$, $z_2$ and $z_3$. Applying the criterion in \eqref{eq:concyclic}, we obtain the condition
	$$
	\imag\Big( \frac{x+iy}{x+iy-2i} \cdot \frac{1-2i}{1}\Big) = 0,
	$$
	which can then be simplified to
	$$
	x^2 +y^2 - x - 2y =0,
	$$
	or equivalently
	$$
	(x-1/2)^2 + (y-1)^2 = \frac{5}{4}.
	$$
	This is the equation of a circle with center at $(1/2,1)$ and radius $\sqrt{5}/2$.
\end{example}







\section{Python Program}

We implement a test of concyclicity in Python. This program takes for points with real coordinates as inputs, and returns 1 if the four points lie on a circle. Even though this geometrical question itself has nothing to do with complex numbers, we can use complex numbers and complex arithmetic to get the answer easily. The calculation is done in one line.


\begin{lstlisting}[language=Python, frame=single]
# Determine whether four distinct points lie on a circle
	
def test_concyclic(pt0, pt1, pt2, pt3):
"""
  Input four tuples:  pt1= (x1,y1), pt2= (x2,y2), pt3= (x3,y3), pt4=(x4,y4)
  return True if the four points pt1, pt2, pt3, pt4 are concyclic, False if not concyclic
"""
	z0 = complex(pt0[0],pt0[1])   # convert to complex numbers
	z1 = complex(pt1[0],pt1[1])
	z2 = complex(pt2[0],pt2[1])
	z3 = complex(pt3[0],pt3[1])
	return abs(((z0-z1)*(z2-z3)/((z0-z3)*(z2-z1))).imag)< 1e-8
\end{lstlisting}

For example, we can use this Python function to check that four points $(0,0)$, $(1,0)$, $(1,2)$, and $(0,2)$ lie on the circle in the previous example.


\begin{lstlisting}[language=Python, frame=single]
In [1]: test_concyclic((0,0), (1,2), (0,2), (1,0))
Out[1]: True
\end{lstlisting}







\section{Appendix: Details of stereographic projection}
\label{sec:stereographic_projection}
Stereographic projection can be explicitly computed using the formula below. Recall that the sphere $S$ is centered at $(0,0,r)$ with radius $r$.

\begin{greenbox}
	\begin{theorem}
		The functions $P :\mathbb{R}^2 \rightarrow S\setminus \{(0,0,2r)\}$ described in the previous paragraph is given by
		$$
		P(x,y) = \Big( \frac{4r^2x}{4r^2+x^2+y^2},\ \frac{4r^2y}{4r^2+x^2+y^2},\ \frac{2r(x^2+y^2)}{4r^2+x^2+y^2} \Big).
		$$
		This gives a bijection between the complex plane $\mathbb{R}^2$ and the punctured sphere $S\setminus \{(0,0,2r)\}$. \label{thm:stereographic_projection}
	\end{theorem}
\end{greenbox}

We may alternately write the projection function $P$  with $z=x+iy$ as input,
$$
P(x+iy) = \Big( \frac{4r^2 \real{z}}{4r^2 + |z|^2}, \frac{4r^2 \imag{z}}{4r^2 + |z|^2}, \frac{2r|z|^2}{4r^2+|z|^2} \Big).
$$


\begin{proof}
	Let $N$ denote the north pole $(0,0,2r)$ of $S$. The triangle with vertices $N$, $(0,0,0)$ and $(x,y,0)$ and the triangle with vertices $N$, $(0,0,\gamma)$, $P(x,y)$ are similar. This gives
	\begin{equation}\frac{x}{\alpha} = \frac{y}{\beta} = \frac{2r}{2r-\gamma} .
		\label{eq:A}
	\end{equation}
	Substitute $\alpha= \frac{2r-\gamma}{2r} x$ and $\beta = \frac{2r-\gamma}{2r}y$ into the equation of sphere $S$ (Equation~\ref{eq:sterographic_projection_proof1}), we get
	\begin{align*}
		\big(\frac{2r-\gamma}{2r}\big)^2(x^2+y^2) + (\gamma - r)^2 & = r^2 \\
		(2r-\gamma) (x^2+y^2) &= 4r^2\gamma .
	\end{align*}
	In the last step, we cancel the factor $2r-\gamma$ from both sides. This is possible because $\gamma<2r$.
	We can express $\gamma$ in terms of $x$ and $y$
	$$
	\gamma = \frac{2r(x^2+y^2)}{4r^2+x^2+y^2}.
	$$
	Hence,
	$$
	2r-\gamma = \frac{8r^3}{4r^2+x^2+y^2}.
	$$
	Substituting the expression of $2r-\gamma$ back to $\eqref{eq:A}$ gives 
	$$\alpha = \frac{4r^2x}{4r^2+x^2+y^2} \text{ and }\beta = \frac{4r^2y}{4r^2+x^2+y^2}.
	$$
	This proves that $P(x,y)$ can be written as in the proposition.
	
	We now verify that $P$ is a bijection by checking that $P$ has an inverse function
	\begin{equation}
		Q(\alpha,\beta,\gamma) \triangleq \big( \frac{2r\alpha}{2r-\gamma}, \frac{2r\beta}{2r-\gamma} \big) .
		\label{eq:reverse_stereographical_projection}
	\end{equation}
	For any point $(\alpha,\beta,\gamma)$ on $S\setminus \{N\}$, we can compute
	\begin{align}
		P(Q(\alpha,\beta,\gamma)) &= P\Big(\frac{2r\alpha}{2r-\gamma}, \frac{2r\beta}{2r-\gamma}\Big) \notag\\
		&=\Big( \alpha \frac{2r(2r-\gamma)}{(2r-\gamma)^2+\alpha^2+\beta^2}, \beta \frac{2r(2r-\gamma)}{(2r-\gamma)^2+\alpha^2+\beta^2},
		\frac{2r(\alpha^2+\beta^2)}{(2r-\gamma)^2+\alpha^2+\beta^2}
		\Big) .  \label{eq:B}
	\end{align}
	From the equation of $S$, we have
	\begin{equation}
		\alpha^2 +  \beta^2  = 2r\gamma -  \gamma^2 .
		\label{eq:C}
	\end{equation}
	The demoninators in \eqref{eq:B} can be simplified to
	$$(2r-\gamma)^2 + \alpha^2 + \beta^2 = 4r^2 - 4r\gamma + \gamma^2 + \alpha^2 + \beta^2 = 4r^2-2r\gamma.
	$$
	The point in \eqref{eq:B} can then be simplified to $(\alpha, \beta, \gamma)$. This proves that the composite function $P \circ Q$ is the identity function on the punctured sphere.
	
	In the other direction, we want to show that 
	\begin{align*}
		Q(P(x,y)) &= Q  \Big( \frac{4r^2x}{4r^2+x^2+y^2},\ \frac{4r^2y}{4r^2+x^2+y^2},\ \frac{2r(x^2+y^2)}{4r^2+x^2+y^2} \Big) 
	\end{align*}
	is equal to the point $(x,y)$. To this end, we compute
	\begin{align*}
		\frac{2r\alpha}{2r - \gamma} &=	\frac{2r \frac{4r^2x}{4r^2+x^2+y^2} }{2r - \frac{2r(x^2+y^2)}{4r^2+x^2+y^2} } = x
	\end{align*}
	and
	\begin{align*}
		\frac{2r\beta}{2r - \gamma} &=	\frac{2r \frac{4r^2y}{4r^2+x^2+y^2} }{2r - \frac{2r(x^2+y^2)}{4r^2+x^2+y^2} } = y.
	\end{align*}
	This proves that the composite function $Q\circ P$ is the identity function on the complex plane.
	
	Because $Q\circ P$ and $P\circ Q$ are identity functions, we conclude that $P$ is a bijection.
\end{proof}


\medskip

The stereographic projection is not merely a bijection, it has the feature of preserving circle and straight line. Suppose that a plane with equation
\begin{equation}
	a \alpha + b \beta + c \gamma = d
	\label{eq:stereographic_projection_proof2}
\end{equation}
for some constants $a$, $b$, $c$ and $d$ intersects with the sphere $S$. We know that the cross-section is circular in shape. We derive the equation that describes the stereographic projection of the cross-section as follows.




Consider a plane $\Pi$ that contains the north pole of $S$. This plane intersects the sphere $S$ at a circle as long as it is not a horizontal plane, i.e., when $a$ and $b$ are not both zero. 
Using the fact that the vector $(a,b,c)$ is a normal vector of $\Pi$, we can write the equation of $\Pi$ as
\begin{align*}
	[(\alpha,\beta,\gamma) - (0,0,2r)] \cdot (a,b,c) &= 0 \\
	a \alpha + b \beta + \gamma c   &= 2rc.
\end{align*}
We note that when the plane $\Pi$ passes through the north pole, the constant $d$ in \eqref{eq:stereographic_projection_proof2} is equal to $d=2rc$. The intersection of $\Pi$ and the horizontal plane is the straight line described by the equation
$a \alpha + b \beta = 2rc$.



Next, suppose $a\alpha+b\beta+c\gamma=d$ is the equation of a plane $\Pi$ that cut through the sphere $S$ but does not pass through the north pole $N$. 
Let $(x,y)$ be a point on the horizontal plane such that the point $(\alpha,\beta,\gamma)=P(x,y)$ lies on the plane $\Pi$. We  obtain the equation
\begin{align*}
	a \frac{4r^2x}{4r^2+x^2+y^2} + b \frac{4r^2y}{4r^2+x^2+y^2} + c \frac{2r(x^2+y^2)}{4r^2+x^2+y^2} &= d \\
	4r^2ax + 4r^2by + 2rc(x^2+y^2) & = d(4r^2+x^2+y^2) \\
	(d-2rc)(x^2+y^2) - 4r^2ax - 4r^2by + 4r^2d  &=0.
\end{align*}
This equation defines a circle provided that $d\neq 2rc$, which happens if and only if the plane $\Pi$ does not contain the north pole. By completing the squares we see that the center of the circle is located at
\[
(x,y) = \frac{2r^2}{d-2rc}(a,b)
\]
and the radius is
\begin{equation}
	\frac{2r}{|d-2rc|} \sqrt{r^2(a^2+b^2)-d(d-2rc)}.
	\label{radius}
\end{equation}

The number $r^2(a^2+b^2)-d(d-2rc)$ should be positive, because the radius is supposed to be positive. We can show this algebraically. We are assuming that the perpendicular distance between the plane and the center of $S$ is less than~$r$, otherwise the plane $\Pi$ and the sphere $S$ has no intersection, or intersect at a single point. We  express this condition using the dot product 
$$
\Big| (\alpha,\beta,\gamma-r) \cdot \frac{(a,b,c)}{\sqrt{a^2+b^2+c^2}} \Big| < r.
$$
This condition can be simplified to
$$
(\alpha a + \beta b + \gamma c - rc)^2 < r^2 (a^2+b^2+c^2)
$$
which can be further simplified to
$$
d^2 - 2drc < r^2(a^2+b^2).
$$
This is equivalent to the condition that $r^2(a^2+b^2)-d(d-2rc)$  is positive.




\section{Appendix: Complex projective line} \label{sec:complex_projective_geometry}
There is alternate description of the extended complex plane using projective geometry. We use the notation $\mathbb{C}^2$ to represent the complex vector space consisting of column vector $\begin{bmatrix}z_1 \\ z_2\end{bmatrix}$, with  $z_1$ and $z_2$ not both zero. Define an equivalence relation by declaring  $\begin{bmatrix}z_1 \\ z_2\end{bmatrix} \sim \begin{bmatrix}z_1' \\ z_2'\end{bmatrix}$ if there exists a complex number $\lambda$ such that $z_1' = \lambda z_1$ and $z_2' = \lambda z_2$. Two pairs are thus equivalent if they lie on the same complex line  that passes through the origin in $\mathbb{C}^2$. 


The {\em complex projective line} is defined as the set that consists of all equivalence classes defined in the previous paragraph. Common notation for the complex projective line are $\mathbb{CP}_1$ or $\mathbb{C}P^1$. 
In an equivalence class, the exact values of $z_1$ and $z_2$ are not relevant, but
only the fraction $z_1/z_2$ is relevant.  When $z_2$ is nonzero, the fraction $z_1/z_2$ is a complex number, and we can pick the vector $\begin{bmatrix}z \\ 1 \end{bmatrix}$, with $z=z_1/z_2$, as the representative of this equivalence class. If $z_2=0$, we can select the vector $\begin{bmatrix} 1 \\ 0 \end{bmatrix}$ as the representative.

We have  a bijection between the extended complex plane  and the complex projective line.
A finite complex number $z$ corresponds to the equivalence class that contains $\begin{bmatrix} z \\ 1 \end{bmatrix}$. The point at infinity is the equivalence class that contains $\begin{bmatrix} 1 \\ 0 \end{bmatrix}$. 
We can thus regard the extended complex number system as a set consisting of elements in the form $\begin{bmatrix} z \\ 1 \end{bmatrix}$, for $z\in\mathbb{C}$, and a special element $\begin{bmatrix} 1 \\ 0 \end{bmatrix}$.





We can perform a linear fractional transformation by multiplying $\begin{bmatrix} z_1 \\  z_2 \end{bmatrix}$ from the left by a $2\times 2$ matrix,
$$
\begin{bmatrix} a & b \\ c & d \end{bmatrix}
\begin{bmatrix} z_1 \\ z_2 \end{bmatrix}
$$
One can show that it does not depend on the choice of $\begin{bmatrix} z_1 \\ z_2 \end{bmatrix}$ in an equivalence class. If $z_2\neq 0$, the resulting vector is equivalent to 
$$
\begin{bmatrix} \frac{az+b}{cz+d} \\ 1\end{bmatrix}
$$
where $z = z_1/z_2$.

If we multiply with the point at infinity $\begin{bmatrix} 1 \\ 0 \end{bmatrix}$, we obtain
$$\begin{bmatrix} a & b \\ c & d \end{bmatrix}\begin{bmatrix} 1 \\ 0 \end{bmatrix} \sim \begin{bmatrix} a \\ c \end{bmatrix}
$$

If we multiply with the projective point that corresponds to the complex number $-d/c$,
$$\begin{bmatrix} a & b \\ c & d \end{bmatrix}\begin{bmatrix} -d/c \\ 1 \end{bmatrix} \sim \begin{bmatrix} a(-d/c)+b \\ 0 \end{bmatrix},
$$
which is equivalent to the point at infinity (using the assumption that $ad-bc\neq 0$.) Therefore, the definition using matrix multiplication is isomorphic to Definition~\ref{def:LFT}.



\begin{comment}
	
	\section{Complex number in LEAN}
	
	We can encode the mathematical statements and proofs using LEAN and the Mathlib library. In LEAN, real numbers are constructed by Cauchy sequences. There is no precision issue, and hence, when we say real number $x$ is equal to real number $y$, we mean exact equality. However, the real number type is not computable in LEAN, which is in contrast to Python and Matlab. 
	
	Given that real numbers are available, complex number is represented by a data structure that contains two real numbers: the real part and the imaginary part.
	
	
	\begin{lstlisting}[language=Lean, frame=single]
		theorem funext {f₁ f₂ : ∀ (x : α), β x} (h : ∀ x, f₁ x = f₂ x) : f₁ = f₂ := by
		show extfunApp (Quotient.mk f₁) = extfunApp (Quotient.mk f₂)
		apply congrArg
		apply Quotient.sound
		exact h
	\end{lstlisting}
\end{comment}

\section*{Summary}


\begin{itemize}
\item 	A {linear fractional transformation}, also known as {M\"obius} transformation, is a function in the form 
$	f(z) = \frac{az+b}{cz+d}$, 
where $a$, $b$, $c$ and $d$ are complex numbers satisfying $ad-bc\neq 0$. This class of functions include translation, rotation, expansion, and inversion as special cases.

\item The stereographic projection maps the points on the complex plane to the Riemann sphere injectively. The extended complex number system $\bar{\mathbb{C}}$ is obtained by adjoining a point at infinity to $\mathbb{C}$.

\item  Straight lines and circles on the complex plane are mapped to circles on the Riemann sphere.

\item We can extend the domain and range of the inversion function $f(z) = 1/z$ by defining the point at infinity as the inverse of~0.

\item Properties of linear fractional transformation:

\begin{enumerate}
\item Composition of two LFT's is an LFT. The inverse function of LFT is  an LFT.

\item The class of linear fractional transformation can map any three distinct points in the extended complex plane to three other distinct points.
\end{enumerate}

\item 	The cross ratio of four complex numbers $z_1$, $z_2$, $z_3$ and $z_4$ is defined as
$$
[z_0,z_1,z_2,z_3] \triangleq \frac{z_0 - z_1}{z_0-z_3} \cdot \frac{z_2-z_3}{z_2-z_1}
$$
when $z_0,\ldots, z_3$ are all finite complex numbers. When one of $z_0,\ldots, z_3$ is $\infty$, the value can be computed by applying the convention $\infty/\infty = 1$.

\item The value of cross ratio is invariant under linear fractional transformation.
\end{itemize}

\section*{Exercises}

\renewcommand{\labelenumi}{\arabic{enumi}.}
\begin{enumerate}



\item 
Let $A,B \in \mathbb{C}$ and let $k>0$ be a constant.  
Consider the locus of points $P$ satisfying
\[
\frac{|P-A|}{|P-B|} = k.
\]

\begin{enumerate}
	\item Show that this is the equation of a circle (or a line when $k=1$) by rewriting the condition in complex form as $|z-a| = k|z-b|$.
	\item Find the center and radius of the circle explicitly.
\end{enumerate}



\item Find the image of the half plane $x>c$, where $c$ is a positive real number, under the inversion function $f(z)=1/z$.



\item  Find the image of the following geometric shapes under the map
$$
f(z) = \frac{z-i}{z+i}.
$$

\begin{enumerate}
\item the closed disk with radius $1$ and center at $i$, i.e., the set of points that satisfy $|z-i| \leq  1$.
\item  the closed unit disk $|z|\leq 1$.
\end{enumerate}

\item Show that the linear fractional transformation
$$
w = \frac{z - z_0}{z_0^* z - 1}
$$
where $|z_0| < 1$, maps the unit circle $|z|=1$ onto the unit circle $|w|=1$. 
%(Hint: Let $z$ be a complex number with modulus 1, calculate $|w|^2 = ww^*$, and simplify. Indicate at which step you use the assumption $|z_0|<1$.)



\item  Find the linear fractional transformation $w = \frac{az+b}{cz+d}$ that maps:
\begin{flalign*}
	(a) \ & 
	z_1 = 1 \text{ to } w_1= i, 
	z_2 = i \text{ to } w_2= -1, 
	z_3 = -1 \text{ to } w_3= -i; &&&& \\ 
	(b) \ & 
	z_1 = \infty \text{ to } w_1= 1, 
	z_2 = 0 \text{ to } w_2= i, 
	z_3 = 1 \text{ to } w_3= \infty. &&&& 
\end{flalign*}

\item Prove that the solutions to $z^n=(z+1)^n$ lie on the vertical line $\real(z) = -1/2$ in the complex plane.

\item Describe geometrically the transformation in parts (a) and (b) by decomposing them into elementary transformations (translation, rotation, dilation, inversion).
\begin{flalign*}
	(a)\ & w = (1+i)z + 2 & (b)\ & w = \frac{z+1}{z-1} &&&&
\end{flalign*}


\item Write the linear fractional transformation 
$$
f(z) = \frac{z-2}{z+2}
$$
as a composition of translation, scaling, rotation, and inversion.

\item Show that a linear fractional transformation has either 1 or 2 fixed points in the extended complex plane. Suppose $f(z)=(az+b)/(cz+d)$. Derive a condition on coefficients $a$, $b$, $c$ and $d$ under which $f(z)$ has two distinct fixed points.

\item Let $C_1$ and $C_2$ be two distinct circles in the complex plane. Prove that there exists a linear fractional transformation $f(z) = (az+b)/(cz+d)$ that maps $C_1$ and $C_2$ onto two circles $C_1'$ and $C_2'$ with the same finite radius. 

(Hint: Analyze the problem in four cases based on the intersection of the circles: (i) one includes the other in the interior, (ii) disjoint and exterior, (iii) intersecting at two points, and (iv) tangent at one point.)

\item  Let $C_1$ and $C_2$ be two disjoint circles in the complex plane (one strictly inside the other). Prove that there exists a linear fractional transformation that maps $C_1$ and $C_2$ onto two concentric circles centered at the origin.

\item Let $C_1$ and $C_2$ be two circles that intersect at two distinct points $z_a$ and $z_b$ with an angle of intersection $\alpha$ ($0 < \alpha < \pi$).
Find a linear fractional transformation that maps the intersection region (the ``lens'' shape) onto an infinite angular sector (a wedge) of angle $\alpha$ with the vertex at the origin.


\item Let \(z_1 = 0\), \(z_2 = 1\), \(z_3 = \infty\), and \(z_4 = x \in \mathbb{C}\setminus\{0,1\}\).  
Compute the cross ratio
\[
[0,1,\infty,x].
\]
What is the geometric meaning of your answer.

\item Let \((z_1,z_2,z_3,z_4)\) and \((w_1,w_2,w_3,w_4)\) be two quadruples of distinct points in $\bar{\mathbb{C}}$.  
Show that if
\[
[z_1,z_2,z_3,z_4] = [w_1,w_2,w_3,w_4],
\]
then there exists a linear fractional transformation \(T\) such that \(T(z_i) = w_i\) for all \(i\).

\item  Two points $z$ and $z^*$ are said to be {\em symmetric with respect to a circle} $C$ if the line passing through the center of circle $C$ and $z$ also passes through $z^*$, and $|z-c||z^*-c| = R^2$ (where $c$ is the center and $R$ is the radius).
Prove that linear fractional transformations preserve symmetry. That is, if $z$ and $z^*$ are symmetric with respect to a circle $C$, and $T$ is an LFT, show that $T(z)$ and $T(z^*)$ are symmetric with respect to the image circle $T(C)$.

(Hint: Use the cross-ratio definition of symmetry: $z, z^*$ are symmetric with respect to $C$ passing through $z_1, z_2, z_3$ if and only if the cross-ratio $(z^*, z_1, z_2, z_3) = \overline{(z, z_1, z_2, z_3)}$. Use the fact that LFTs preserve cross-ratios.)


\item Using the cross ratio and linear fractional transformation, prove Ptolemy's theorem: four points $A$, $B$, $C$ and $D$ lie on a circle if and only if
$$
\overline{AB} \cdot \overline{CD} + \overline{AD} \cdot \overline{BC} = \overline{AC} \cdot \overline{BD}.
$$

\item Draw a circle with radius $r$ centered at $O$. Pick two points $A$ and $B$ inside that circle so that $A\neq 0$ and $A\neq B$. On the line containing $O$ and $A$, draw a point $A^*$  so that $OA\cdot OA^* = r^2$, and $A$ is between $O$ and $A$. Similarly, on the line containing $O$ and $B$, draw a point $B^*$ so that $OB \cdot OB^* = r^2$, and $B$ is between $O$ and $B$. Show that $A$, $B$, $A^*$ and $B^*$ lie on a circle. 




\end{enumerate}








%%%%%%%%%%%%%%%%%%%%%%%%%%%%%%%


\chapter{Complex Exponential Function and Euler's formula}


In this chapter we introduce the complex exponential function by defining it as an infinite series, extending its domain from the real number line to the complex plane. 



\section{Convergent sequences} 

\noindent Notation: We denote a \df{sequence} of complex numbers $z_1, z_2, z_3,\ldots$ is denoted by $(z_k)_{k=1}^\infty$ or simply $\{z_k\}$.

\begin{svgraybox}
	\begin{definition} \index{Sequence}
		\index{Sequence!Convergent}
		\index{Sequence!Divergent}
		Given a complex sequence $(z_n)_{n=1}^\infty$, we say that $z_n$ \df{converges} to a complex number $w \in \mathbb{C}$ if 
		$$
		\forall \epsilon>0, \ \exists N, \ s.t.\ |z_n-w|< \epsilon, \ \forall n\geq N.
		$$		
		We write 
		$$
		\lim_{n\rightarrow \infty} z_n = w
		$$
		if $(z_n)_{n=1}^\infty$ converges to the limit $w$.
		
 A sequence that is \df{divergent} if it is not converging to any complex number in $\mathbb{C}$.
		\label{def:sequence_converge}
	\end{definition}
\end{svgraybox}

Geometrically, if a sequence  is convergent, it is the same as saying that for any small $\epsilon$, the points in the complex sequence will eventually stay inside the open disk $D(w,\epsilon)$.

\begin{example} The complex sequence $\{z_n\}$ defined by
	$$
	z_n = \frac{1}{n} + \frac{i}{n^2}
	$$
	converges to 0, because $|z_n|\rightarrow 0$ as $n\rightarrow \infty$.
\end{example}


\begin{example}
	Compute $\lim_{n\rightarrow\infty} \frac{n}{n+i}$.
	
	We can first make a guess that the limit should be 1, because when $n$ is large, adding $i$ to $n$ has negligible effect. To make the argument rigorous, we write
	$$
	\frac{n}{n+i} = \frac{n+i-i}{n+i} = 1 - \frac{i}{n+i}.
	$$
	Then
	$$
	\Big| \frac{n}{n+i} - 1 \Big| = \Big| \frac{i}{n+i} \Big| = \frac{1}{\sqrt{n^2+1}} \rightarrow 0, \qquad \text{ as } n\rightarrow\infty.
	$$
	Therefore $n/(n+i) \rightarrow 1$ as $n\rightarrow\infty$.
\end{example}

The convergence of a complex sequence reduces to the real case through the following theorem.

\begin{svgraybox}
	\begin{theorem}
 A complex sequence  $(z_n)_{n=1}^\infty$ converges if and only if both $(\real(z_n))_{n=1}^\infty$ and $(\imag(z_n))_{n=1}^\infty$ converge. \label{thm:complex_converge_real_imag}
	\end{theorem}
\end{svgraybox}

It is easy to show that a complex sequence is Cauchy if and only if the real and imaginary parts are real Cauchy sequences. As a result, the basic property of Cauchy sequence for real numbers extends to the complex case readily.

\begin{svgraybox}
	\begin{definition} \index{Cauchy sequence}
		A sequence $(z_n)_{n=1}^\infty$ is called a \df{Cauchy sequence} if for all $\epsilon>0$, there exists an integer $N$ such that
		$$
		|z_m - z_n| < \epsilon \ \ \text{ whenever } m,n\geq N.
		$$
	\end{definition}
\end{svgraybox}

The set of complex numbers inherits the completeness property of real numbers.


\begin{svgraybox}
	\begin{theorem}
		A complex sequence  $(z_n)_{n=1}^\infty$ converges if and only if $(z_n)_{n=1}^\infty$ is Cauchy.
		\label{prop:Cauchy_sequence}
	\end{theorem}
\end{svgraybox}

\begin{proof}
The proof is a reduction to the real parts and the imaginary parts of the sequence by Theorem~\ref{thm:complex_converge_real_imag}.
\end{proof}


If we consider the complex sequence as points on the Riemann sphere, then it may converge to the point at infinity, even though it does not converge to any point in the complex plane.

\begin{svgraybox}
	\begin{definition} \index{Convergence!at the point at infinity}
		A sequence of complex numbers $\{z_k\}_{k=1}^\infty$ is said to \df{converge to the point at infinity} if for any positive number $r$, there exists a sufficiently large integer $N$ such that $|z_k| > r$ for all $k\geq N$.
	\end{definition}
\end{svgraybox}

We say that a sequence ${z_k}$ is \df{bounded} if there exists a real number $M$ such that $|z_k|\leq M$ for all $k\geq 1$. A sequence that converges to $\infty$ cannot be bounded.


When the sequence $\{z_k\}_{k=1}^\infty $ converges to the point at infinity, we write $z_n\rightarrow \infty$ as $n\rightarrow \infty$, or
$$
\lim_{n\rightarrow \infty} z_n = \infty.
$$



\begin{example}
	The sequence $(n(1+i))_{n=1}^\infty$ is divergent, but is converging to the point at infinity on the Riemann sphere.
\end{example}



\begin{example} This example highlight a difference between sequence on the extended complex plane and sequence on the extended real number system.
	Let $\{z_k\}_{k=1}^\infty$ be a sequence of complex numbers such that
	$$
	\lim_{n\rightarrow \infty} z_n = \infty.
	$$
	Consider another sequence of complex numbers defined by $w_k = -z_k$. The sequence $\{w_k\}_{k=1}^\infty$ converges to the point at infinity as well,
	$$
	\lim_{n\rightarrow \infty} w_n = \lim_{n\rightarrow \infty} (-z_n)  =  \infty.
	$$
	The two sequences $(w_n)_{n\geq 1}$ and $(z_n)_{n\geq 1}$ have the same limit, even though they proceed in opposite direction.
\end{example}

We record some basic properties of complex sequences in the following theorem. The proof is the same as in the real case, and hence is omitted.

\renewcommand{\labelenumi}{(\alph{enumi}).}
\begin{svgraybox}
	\begin{theorem}
		Let $(z_n)_{n=1}^\infty$ and $(w_n)_{n=1}^\infty$ be convergent complex sequences, with limits $\alpha$ and $\beta$, respectively.
		
		\begin{enumerate}
			\item $\lim_{n\rightarrow\infty} (az_n + b w_n) = a\alpha + b \beta$, for complex constants $a$ and $b$.
			\item  $\lim_{n\rightarrow\infty} z_n w_n = \alpha \beta$.
			\item  $\lim_{n\rightarrow \infty} z_n/w_n = \alpha/\beta $, if $\beta\neq 0$.
			\item $\lim_{n\rightarrow\infty} |z_n| = \big| \alpha \big|$ .
			\item $\lim_{n\rightarrow\infty} (z_n^*) = \alpha^*$.
		\end{enumerate}
	\end{theorem}
\end{svgraybox}




\section{Complex series}

\begin{svgraybox}
	\begin{definition} \index{Infinite series} Given a sequence of complex numbers $(z_k)_{k=1}^\infty$, we define the value of the \df{infinite series} $\sum_{k=1}^\infty z_k$ by
		$$
		\sum_{k=1}^\infty z_k \triangleq \lim_{n\rightarrow\infty}\sum_{k=1}^n z_k
		$$
		if the limit exists. 
	\end{definition}
\end{svgraybox}


As in calculus, the convergence of infinity series is the same as the convergence of sequence of partial sums.


\begin{example}(Complex geometric series)
	Evaluate $\sum_{k=1}^\infty (0.5i)^k$.
	
	For any finite $n$, we have
	$$
	\sum_{k=1}^n (0.5i)^k = \frac{(0.5i)^{n+1} - 0.5i}{0.5i-1}.
	$$
	We take limit as $n\rightarrow\infty$,
	$$
	\sum_{k=1}^\infty (0.5i)^k = \lim_{n\rightarrow\infty} \frac{(0.5i)^{n+1} - 0.5i}{0.5i-1} =
	\frac{-0.5i}{0.5i-1} = \frac{-1+2i}{5}.
	$$
\end{example}

We record a few elementary properties of complex series in the next theorem.

\begin{svgraybox}
	\begin{theorem} \label{thm:conjugate_power_series}
Let $\sum_{k=1}^\infty z_k$ denote a convergent complex series.
\begin{enumerate}
	\item $\sum_{k=1}^\infty (z_k^*) = \big(\sum_{k=1}^\infty z_k \big)^*$.
	\item $\real\big( \sum_{k=1}^\infty z_k\big) = \sum_{k=1}^\infty \real(z_k)$.
	\item $\imag\big( \sum_{k=1}^\infty z_k\big) = \sum_{k=1}^\infty \imag(z_k)$.
\end{enumerate}
	\end{theorem}
\end{svgraybox}

We give a proof of the first part of this theorem. The proof of the remaining parts are similar.

\begin{proof}
	Let $L$ be the limit of $\sum_{k=1}^\infty z_k$. By definition,
	for any positive real number $\epsilon$, there exists a positive integer $N$ such that
	$$
	\Big| \sum_{k=1}^n z_k - L \Big| \leq \epsilon
	$$
	for all $n\geq N$. Because a complex number and its conjugate have the same absolute value, we can write
	$$
	\Big| \sum_{k=1}^n z_k^* - L^* \Big| \leq \epsilon
	$$	
	for any  $n \geq N$. Hence by the definition of limit, 
	$$
	\sum_{k=1}^\infty z_k^* = L^* = \big( \sum_{k=1}^\infty z_k \big)^* .
	$$
\end{proof}

As in the real case, we have the linearity property for infinite series.
\begin{svgraybox}
	\begin{theorem}
If $\sum_{k=1}^\infty z_k$ and $\sum_{k=1}^\infty w_k$ are convergent complex series, and $\alpha$ and $\beta$ are complex constants, then
$$
\sum_{k=1}^\infty (\alpha z_k + \beta w_k) =
\alpha \sum_{k=1}^\infty z_k + \beta \sum_{k=1}^\infty w_k.
$$
	\end{theorem}
\end{svgraybox}


The next two proposition provides two tests for convergence. The first one is a test for divergence, and the second one is a test for convergence.

\begin{svgraybox}
	\begin{proposition} ($n$-th term test)
		If $\sum_{k=1}^\infty z_k$ converges, then $|z_n|\rightarrow 0$ as $n\rightarrow \infty$. In other words, if $(|z_n|)_{n=1}^\infty$ does not converge to 0, then $\sum_{k=1}^\infty z_k$ does not converge.\label{prop:n_term_test}
	\end{proposition}
\end{svgraybox}



\begin{proof}
	Suppose $\sum_{k=1}^\infty z_k$ converges to~$w$. The sequence of partial sums $(\sum_{k=1}^n z_k)_{n=1}^\infty$ converges, and is thus a Cauchy sequence (Prop.~\ref{prop:Cauchy_sequence} part (ii)). By the definition of Cauchy sequence, given any $\epsilon >0$, there exists an integer $N$ such that
	$$
	\big|\sum_{k=1}^m z_k - \sum_{k=1}^n z_k \big| < \epsilon
	$$
	for any $n,m \geq N$. By picking $m=n+1$, we obtain $|z_{n+1}| < \epsilon$ for all $n\geq N$.
\end{proof}

\begin{svgraybox}
	\begin{proposition} (Absolute convergence test)
		Given a sequence of complex numbers $(z_k)_{k=1}^\infty$, if the real infinite series $\sum_{k=1}^\infty |z_k|$ converges, then $\sum_{k=1}^\infty z_k$ also converges.\label{prop:abs_convergence}
	\end{proposition}
\end{svgraybox}


\begin{proof}
	Suppose  $\sum_{k=1}^\infty |z_k|$ converges. Since $|\real(z_k)| \leq |z_k|$ for all $k\geq1$, the real series  $\sum_{k=1}^\infty |\real(z_k)|$ converges.
	By the absolute convergence property for real series, $\sum_{k=1}^\infty \real(z_k)$ converges. Similarly, $\sum_{k=1}^\infty \imag(z_k)$  is convergent. The convergence of $\sum_{k=1}^\infty z_k$ then follows by applying part (i) of Theorem.~\ref{prop:Cauchy_sequence}.
\end{proof}


\begin{example}
	The series 
	\begin{equation}
		\sum_{k=1}^\infty \frac{2+i/k}{k^2+2ki}
		\label{ex:convergence_series}
	\end{equation}
	is convergent because we can upper bound the magnitude of each term
	$$\Big|\frac{2+i/k}{k^2+2ki}\Big| \leq c/k^2$$
	for some real constant $c$. Because the series $\sum_{k=1}^\infty k^{-2}$ is convergent, by Prop.~\ref{prop:abs_convergence}, the series in~\eqref{ex:convergence_series} converges.
\end{example}

In view of Prop.~\ref{prop:abs_convergence}, we make the following definition.

\begin{svgraybox}
	\begin{definition} \index{Convergence!absolute} \index{Convergence!conditional} \label{def:absolute_convergence}
		We say that a series $\sum_{k=1}^\infty z_k$ is \df{absolutely convergent} if $\sum_{k=1}^\infty |z_k|$ is convergent. If  $\sum_{k=1}^\infty z_k$ converges but $\sum_{k=1}^\infty |z_k|$ does not converge, then we say that it \df{converges conditionally}.
	\end{definition}
\end{svgraybox}



\begin{svgraybox}
	\begin{proposition} (limit ratio test) \index{Limit ratio test} 		\label{prop:ratio_test}
		Suppose $z_k$, for $k=1,2,3,\ldots$ are complex numbers and $\lim_{k\rightarrow\infty} |z_{k+1}/z_k|$ exists and is equal to $L$. 
		\begin{enumerate}
			\item If $L>1$, then $\sum_{k=1}^\infty z_k$ diverges;
			\item If $L<1$, then $\sum_{k=1}^\infty z_k$ converges absolutely.
		\end{enumerate}
	\end{proposition}
\end{svgraybox}


\begin{proof}
	First consider the case $L>1$. By the meaning of converging to $L$, which is larger than 1 by assumption, there exists an integer $N$ such that $|z_{k+1}/z_k|>1$ for all $k \geq N$. Therefore
	$$
	|z_N| < |z_{N+1}| < |z_{N+2}| < \cdots
	$$
	By applying the $n$-th term test, we see that $\sum z_k$ is divergent.
	
	Next consider the case $L<1$. Let $\epsilon$ be a small positive real number such that $L<1-\epsilon <1$. By the definition of convergence, there exists an integer $N$ such that
	$$
	\Big| \frac{z_{k+1}}{z_k} \Big| < 1-\epsilon <1
	$$
	for all $k\geq N$. Then, we have
	\begin{align*}
		|z_{N+1}| &< |z_N|(1-\epsilon) \\
		|z_{N+2}| &< |z_{N+1}|(1-\epsilon) < |z_N|(1-\epsilon)^2,
	\end{align*}
	and in general
	$$
	|z_{N+k}| < |z_N|(1-\epsilon)^k
	$$
	for $k\geq 1$. Therefore, for each positive number $m$, we have
	$$
	\sum_{j=0}^m |z_{N+j}| \leq |z_N| \sum_{j=0}^m (1-\epsilon)^j.
	$$
	Since $\sum_{j=0}^\infty (1-\epsilon)^j$ is finite, the sequence
	$$
	\big(\sum_{j=0}^m |z_{N+j}| \big)_{m=0}^\infty
	$$
	is a bounded and monotonically increasing sequence. Hence the series $\sum_{j=0}^\infty |z_j|$ is convergent. By Prop.~\ref{prop:abs_convergence}, the infinite series $\sum_{j=0}^\infty z_{N+j}$ is also convergent. This proves that $\sum_{k=1}^\infty z_k$ is convergent.
\end{proof}

We quote two properties of absolutely convergent series.

\begin{svgraybox}
	\begin{proposition} (Cauchy product) \index{Cauchy product}	\label{prop:series_convolution}
		If $\sum_{k=0}^\infty a_k$ and $\sum_{k=0}^\infty b_k$ are absolutely convergent series, then
		$$
		\Big( \sum_{k=0}^\infty a_k\Big) \Big( \sum_{k=0}^\infty b_k \Big)= \sum_{k=0}^\infty c_k,
		$$
		where
		$$
		c_k = a_0b_k + a_1 b_{k-1} + \cdots + a_{k-1} b_1 + a_k b_0.
		$$
	\end{proposition}
\end{svgraybox}

\begin{svgraybox}
	\begin{proposition}
		If a series converges absolutely, then a series obtained by re-arranging the terms converges and converges to the same limit.
	\end{proposition}
\end{svgraybox}


The proofs can be found in \cite[Theorem 3.50]{Rudin} and \cite[Theorem 3.55]{Rudin}. Although the corresponding theorems in \cite{Rudin} are stated for real series, but the proofs works for the complex case in a straightforward manner.




\section{Convergence of power series}




We study the convergence of power series in this section. For ease of notation, we assume that the center of the power series is the origin. All results below hold in the general case when the center is not the origin.

\begin{svgraybox}
	\begin{proposition}
		Suppose $\sum_{k=0}^\infty a_k z^k$ converges at $z_0\neq 0$, then it converges absolutely for all $z$ with $|z|<|z_0|$.
		\label{prop:absolute_convergence}
	\end{proposition}
\end{svgraybox}

\begin{proof}
	By the $n$-th term test (Prop.~\ref{prop:n_term_test}) we know that 
	$|a_k z_0^k| \rightarrow 0$ as $k\rightarrow\infty$.
	Fix $\epsilon >0$.  there is a sufficiently large $N$ such that
	$$
	|a_k z_0^k| < \epsilon, \qquad \forall k\geq N.
	$$
	Hence, for all $k\geq 0$, we have
	$$
	|a_k z_0^k| \leq \max(\epsilon, |a_0|, |a_1z_0|, |a_2 z_0^2|, \ldots, |a_{N-1} z_0^{N-1}|) \triangleq C.
	$$
	We fix a complex number $z$ with $|z|<|z_0|$. Let $\rho$ be the ratio $|z|/|z_0|$.
	For each $k$, re-write $|a_k z^k|$ as
	$$
	|a_k z^k| = |a_k z_0^k| \frac{|z|^k}{|z_0|^k} \leq C \rho ^k,
	$$
	We note that the constant $C$ depends on the value of $|z|$.
	
	Because $\rho<1$ and $\sum_k C \rho^k$ is a convergent geometric series, we can apply comparison test and conclude that $\sum_k a_k z^k$ is convergent absolutely (See Prop.~\ref{prop:abs_convergence}).
\end{proof}

We can re-write Prop.~\ref{prop:absolute_convergence} as

\begin{svgraybox}
	\begin{proposition}
		If $\sum_k a_k z^k$ diverges at a point $z_1$, then $\sum_{k} a_k z^k$ diverges for all $|z| > |z_1|$.
	\end{proposition}
\end{svgraybox}

\begin{proof}
	Suppose $\sum_k a_kz^k$ diverges at $z_1$. We want to show that if $\sum_k a_k z^k$ converges at any point $z_2$ with modulus $|z_2|>|z_1|$, we would obtain a contradiction.	By the previous proposition, if such a $z_2$ exists, the power series must converge at $z=z_1$, and this directly contradicts the assumption that $\sum_k a_k z^k$ diverges at~$z_1$.
\end{proof}






In view of the previous propositions, we can see that if a power series converges in a region, the region has a circular shape.

\begin{svgraybox}
	\begin{definition} Given a power series $\sum_{k=0}^\infty a_k z^k$, we will call
		\begin{equation} 			\label{eq:radius_of_convergence}
			R \triangleq \sup\big\{|z|:\, \sum_k a_kz^k \text{ converges} \big\}
		\end{equation}
		the  \df{radius of convergence} of $\sum_{k=0}^\infty a_k z^k$. The region $|z|<R$ is called the \df{region of convergence}. 
	\end{definition}
\end{svgraybox}

By Theorem~\ref{prop:absolute_convergence}, the power series is convergent at all points $z$ in the open disk $|z|<R$, and $R$ is the largest radius such that the series converges for all points in the disk $|z|<R$. If $\sum_k a_k z^k$ converges for all $z$, then we write $R=\infty$, and say that the radius of convergence is infinite. If $\sum_k a_k z^k$ converges only at $z=0$, then $R=0$.



\begin{example}
	The power series that define the complex functions $\exp(z)$, $\sin(z)$, $\cos(z)$, $\sinh(z)$ and $\cosh(z)$ have radius of convergence $R=\infty$. 
\end{example}

%\begin{example}
%	If we expand the function $1/(1+z^2)$ as a power series, the resulting power series is
%	$$
%	1 - z^2 + z^4 - z^6 + \cdots
%	$$
%	The radius of convergence is 1.
%\end{example}

\begin{example}
	The power series $\sum_{k=0}^\infty z^k$ diverges for all $z$ on the unit circle, but converges for all $|z|<1$. Specifically, the $k$-th term in the series is $z^k$. When $z$ lies on the unit circle, $|z|$ is equal to 1, had hence $|z|^k$ is also equal to 1. The $k$-th term test says that the power series $\sum_k z^k$ does not converge whenever $|z|=1$.
	
	The the convergence of the series  $\sum_{k=0}^\infty z^k$ in $|z|<1$ is not uniform. We note that the partial sum has a closed-form expression:
	$$
	\sum_{k=0}^n z^k = \frac{1-z^{n+1}}{1-z},
	$$
	which converges to $\frac{1}{1-z}$ for $|z|<1$.  However, the difference
	$$
	\frac{1-z^{n+1}}{1-z} - \frac{1}{1-z} = \frac{z^{n+1}}{1-z}.
	$$
	For each integer $n$, the supremum of $z^{n+1}/(1-z)$ over $|z|<1$ equals $\infty$. By the criterion in \eqref{eq:alternate_uniform_convergence}, the convergence of $\sum_{k=0}^\infty z^k$ is not uniform in $|z|<1$. But the convergence is uniform if we restrict the domain to $|z|<r$ for any positive constant $r$ that is strictly less than 1.
\end{example}


\begin{example}
	The power series $\sum_{k=1}^\infty \frac{1}{k}z^k$ diverges at $z=1$, but converges for all other points on the unit circle. An outline of proof is in the exercise.
\end{example}

\begin{example}
	The power series $\sum_{k=1}^\infty \frac{1}{k^2}z^k$ converges at all points on the unit circle.
	
	When $z=1$, the series $\sum_{k=1}^\infty 1/k^2$ is convergent (converging to $\pi^2/6$). For other point $z$ on the unit circle, the power series $\sum_{n=1}^\infty z^2/k^2$ converges by Theorem~\ref{thm:uniform_convergence_circle}.
\end{example}


\begin{example} (Lacunary series)
	The function $f(z)$ defined by
	$$
	f(z) = \sum_{k=0}^\infty z^{2^k}
	$$
	diverges when $z$ is a $2^n$-th root of unity, for all $n\geq 0$.
\end{example}



An explicit expression for calculating the radius of convergence is given by the Hadamard formula. 

\begin{svgraybox}
	\begin{theorem}[Hadamard formula for radius of convergence] 		\label{thm:Hadamard_formula}
		
		Given a complex power series $\sum_{k=0}^\infty a_k z^k$, the radius of convergence can be computed by
		$$
		R = \frac{1}{\limsup_k |a_k|^{1/k}}.
		$$
	\end{theorem}
\end{svgraybox}

\begin{proof}
	Suppose $z$ is a complex number with $|z|<R$. Pick a real number $R_1$ between $|z|$ and $R$, i.e.,
	$|z| < R_1 < R$. By the definition of $R$ in the theorem,
	$$
	\frac{1}{R_1} > \limsup |a_k|^{1/k}.
	$$
	Using the property of limsup, there exists a sufficiently large $N$ such that for all $k\geq N$,
	$$
	|a_k|^{1/k} < 1/R_1 \qquad \Rightarrow \qquad |a_k|< 1/R_1^k.
	$$
	This yields
	$$
	|a_k z^k| < \frac{|z|^k}{R_1^k}.
	$$
	Since $|z|^k/R_1^k<1$, the geometric series $\sum_k |z|^k/R_1^k$ is convergent. By comparison test, $\sum_k a_kz^k$ is convergent.
	
	Now suppose that $|z|>R$. Pick a real number $R_2$ such that $|z| > R_2 > R$. Taking reciprocal,
	$$\frac{1}{R_2} < \limsup |a_k|^{1/k}.
	$$
	By the defining property of limsup,
	$$
	\frac{1}{R_2} < |a_k|^{1/k} \quad \text{ for infinitely many } k.
	$$
	Hence $|a_kz^k| > \big| \frac{z}{R_2} \big|^k > 1$ for infinitely many $k$. By the $n$th-term test (Prop.~\ref{prop:n_term_test}), the power series diverges whenever $|z|>R$.
\end{proof}





\section{Complex exponential function}
\begin{svgraybox}
	\begin{definition} \index{Complex power series}
		A \df{complex power series} centered at the origin is a series in the form $\sum_{k=0}^\infty a_k z^k$, where $a_k \in \mathbb{C}$.
	\end{definition}
\end{svgraybox}


The power series in the next theorem is the most important power series.

\begin{svgraybox}
	\begin{theorem}
		For any complex number $z\in\mathbb{C}$, the power series
		$\sum_{k=0}^\infty z^k/k!$ converges. Moreover, it converges absolutely.
	\end{theorem}
\end{svgraybox}

\begin{proof}
	It follows from the limit ratio test. The ratio of two consecutive terms has absolute value
	$$
	\frac{\Big| \frac{z^{k+1}}{(k+1)!} \Big|}{\Big| \frac{z^k}{k!}\Big|} =
	\frac{|z|}{k+1}. 
	$$
	For any fixed complex number $z$, the absolute value $|z|$ is fixed, and hence the above ratio has limit zero. By part (b) of Proposition~\ref{prop:ratio_test}, the power series converges absolutely.
\end{proof}

We can also use Theorem~\ref{thm:Hadamard_formula} to show that the radius of convergence is infinite. By Stirling's approximation, we have
$$
k! \sim \sqrt{2\pi k} k^k e^{-k},
$$
so that
$$
( k!)^{1/k} \sim k/e.
$$
The $k$-th root of the $k$-th coefficient approaches 0. By Hadamard formula, the radius of convergence is infinite. The power series converge for all points in the complex plane.

\begin{svgraybox}
	\begin{definition} \index{Exponential function}
		Define the \df{complex exponential function} $\exp(z)$ by the complex power series
		$$
		 \exp(z) \triangleq \sum_{k=0}^\infty \frac{z^k}{k!} = 1 + \frac{z}{1!} + \frac{z^2}{2!} + \frac{z^3}{3!} + \cdots 
		$$
		for $z\in\mathbb{C}$. An alternate notation is $e^z$.
		\label{def:exp}
	\end{definition}
\end{svgraybox}


When $z$ is a real number, the above definition is the same as the power series of the exponential function with real domain. For example, we have $\exp(0)=1$, and $\exp(1)$ is the constant $e$. This power series has infinite radius of convergence.

We establish below a fundamental property of the function $\exp(z)$.

\begin{svgraybox}
	\begin{theorem} 	\label{thm:exp_additive}
		For any complex numbers $z_1$ and $z_2$, 
		$$
		\exp(z_1+z_2) = \exp(z_1)\exp(z_2).
		$$
	\end{theorem}
\end{svgraybox}


\begin{proof}
	We know that both series $\sum_{k=0}^\infty z_1^k/k!$ and $\sum_{k=0}^\infty z_2^k/k!$ converge absolutely. We can apply Prop.~\ref{prop:series_convolution} to $a_k = z_1^k/k!$ and $b_k = z_2^k/k!$. This gives
	$$
	\exp(z_1)\exp(z_2) = \sum_{k=0}^\infty c_k
	$$	
	where $c_k$ is defined by
	$$
	c_k \triangleq  \sum_{\ell=0}^k \frac{z_1^\ell}{\ell !} \frac{z_2^{k-\ell}}{(k-\ell)!} 
	$$
	for $k\geq 0$. By binomial theorem, we can simplify it to
	\begin{align*}
		c_k &= \sum_{\ell=0}^k \frac{z_1^\ell}{\ell !} \frac{z_2^{k-\ell}}{(k-\ell)!} = \frac{1}{k!} \sum_{\ell=0}^k  \binom{k}{\ell} z_1^\ell z_2^{k-\ell} = \frac{1}{k!} (z_1 + z_2)^k.
	\end{align*}
	Hence,
	$$
	\exp(z_1)\exp(z_2) = \sum_{k=0}^\infty \frac{(z_1+z_2)^k}{k!} \triangleq \exp(z_1+z_2).
	$$	
	
\end{proof}

Using this theorem we can derive some immediate properties

\begin{svgraybox}
	\begin{theorem} \label{thm:exp_properties} \ 
		
		\begin{enumerate}
			\item For any $z\in \mathbb{C}$, $\exp(z)\neq 0$, and
$e^{-z} = (e^z)^{-1}$.
			\item If $z=x+iy$, then $\exp(z) = e^x e^{iy}$.
			\item For real number $y$, we have $|e^{iy}| = 1$.			
		\end{enumerate}
	\end{theorem}
\end{svgraybox}


\begin{proof}
	To see that $e^z\neq 0$ for all $z\in\mathbb{C}$, we  take $z_1 = z$ and $z_2 = -z$ in Theorem~\ref{thm:exp_additive},
	$$
	e^{z} e^{-z} = e^{z-z} = e^0 = 1.
	$$
	Since $0$ times any complex number $w$ is equal to 0, we cannot have $e^z=0$, otherwise it would contradict the above equality, which says that there exists a complex number, namely $e^{-z}$, such that $e^z$ time $e^{-z}$ is equal to a nonzero value. We can also deduce from this equality that the multiplicative inverse of $e^z$ is equal to $e^{-z}$.
	
	For any real numbers $x$ and $y$, we let $z_1=x$ and $z_2=iy$ and apply Theorem~\ref{thm:exp_additive} to get
	$$e^{x+iy} = e^x e^{iy}.
	$$
	The factor $e^x$ is a real number because $x$ is real. On the other hand, $e^{iy}$ lies on the unit circle, because 
	$$
	|e^{iy}|^2 = e^{iy} (e^{iy})^* = e^{iy}  e^{-iy} = e^{iy-iy}= e^0 = 1.
	$$
	We have used part 1 of Theorem~\ref{thm:conjugate_power_series} to replace $(e^{iy})^*$ by $e^{-iy}$.
	This proves that $e^{iy}$ has modulus 1.
\end{proof}

%Although the complex exponential function is available from the \textsf{numpy} library, we can implement it from scratch using power series. The function \textsf{my\_exp} below add the first 50 terms in the power series that defines the complex  exponential function, starting from the higher degree to the lower degree.



\section{Euler's formula}

From Theorem~\ref{thm:exp_properties}, we know that $e^{iy}$ lies on the unit circle for any real number $y$. The precise location is given by Euler's formula.

\begin{svgraybox}
	\begin{theorem}[Euler's formula] 	\label{thm:Euler1} \index{Euler's formula}
		For any real number $\theta$, we have
		\begin{equation}
			e^{i\theta} = \cos(\theta) + i \sin(\theta).
			\label{eq:Euler_formula1}
		\end{equation}
	\end{theorem}
\end{svgraybox}

\begin{proof}
	The proof follows from re-arranging the terms in the power series that defines $e^{i\theta}$. This is permissible because 
	$$
	e^{i\theta} = 1 + \frac{i\theta}{1} + \frac{i^2\theta^2}{2!} + \frac{i^3\theta^3}{3!}
	+ \frac{i^4\theta^4}{4!} + \frac{i^5\theta^5}{5!} + \cdots$$
	is an absolutely convergent power series. By separating the real and imaginary parts by terms re-arrangement, we obtain
	$$
	e^{i\theta} = \Big(1 - \frac{\theta^2}{2!}  + \frac{\theta^4}{4!} - \cdots\Big) + i 
	\Big(\theta - \frac{\theta^3}{3!} +\frac{\theta^5}{5!}-\cdots  \Big)
	$$
	the first term is precisely the power series expansion of $\cos(\theta)$, and the second term is $i$ times the power series expansion of $\sin(\theta)$.
\end{proof}

Combining Theorems~\ref{thm:exp_properties} and \ref{thm:Euler1}, we can evaluate the complex exponential function in terms of the real exponential function and the real sine and cosine functions
\begin{svgraybox}
	\begin{theorem}
\begin{equation}
	e^{x+iy} = e^x (\cos(y) + i \sin(y)).
	\label{eq:exponential_function}
\end{equation}
	\end{theorem}
\end{svgraybox}

\begin{remark}
	Some books, such as \cite{BrownChurchill} defines the complex exponential function by~\eqref{eq:exponential_function}. This approach usually defines  holomorphic function first, and then verify that \eqref{eq:exponential_function} is an example of holomorphic function.
	A drawback of this order of definition is that Euler's formula is then true by definition.
	In this book, we took the path of defining the complex exponential function by power series. With this approach, Euler's identity is a magical formula. Nevertheless, we will see later (in Example~\ref{example:exponential_function_uniqueness}) that there is only one complex function that is compatible with the real exponential function. Any definition of complex exponential function will yield the same result.
\end{remark}

\begin{remark}
	The Euler's formula is a generalization of DeMoivre formula. We can recover DeMoivre formula from Euler's formula by considering $e^{in\theta}$, which can be written as
	$$
	e^{in \theta} = \cos(n\theta)+i\sin(n\theta).
	$$
	By Prop.~\ref{thm:exp_additive}, we can also write it as
	$$
	e^{in \theta} = (e^{i\theta})^n = (\cos\theta+i\sin\theta)^n.
	$$
\end{remark}



\begin{example}
	We can compute $e^{\sqrt{2}+i}$ by
	$$
	e^{\sqrt{2}+i} = e^{\sqrt{2}}(\cos(1) + i \sin(1)).
	$$
\end{example}

\begin{example}
	When $z=\pi i$, we obtain
	$$
	e^{\pi i} = \cos(\pi)+i\sin(\pi) = -1.
	$$
\end{example}

From Euler's formula, we see that the complex exponential function is a periodic function with complex period $2\pi i$.
\begin{svgraybox}
	\begin{theorem}
		For any $z\in \mathbb{C}$, we have $\exp(z+2\pi i k) = \exp(z)$ for all integers $k$.
	\end{theorem}
\end{svgraybox}

\begin{example}
	Consider the infinite series
	$$
	\sum_{k=0}^\infty \frac{\cos k\theta}{2^k}
	$$
	as a function of real number $\theta$. This series is convergent by applying the comparison test and comparing with $\sum_{k=0}^\infty 1/2^k$. We can evaluate the exact value by defining $z \triangleq \cos \theta +i \sin \theta$ and  considering the series as the real part of the complex series
	$$
	\sum_{k=0}^\infty \frac{z^k}{2^k}.
	$$
	This is a complex geometric series, and the limit is
	$$
	\frac{1}{1- \frac{z}{2}}.
	$$
	The remaining task is to obtain the real part of $(1-z/2)^{-1}$. To this end, we write $(1-z/2)^{-1}$ as
	\begin{align*}
		\frac{2}{2 - \cos \theta - i\sin\theta} & = \frac{2}{2 - \cos \theta - i\sin\theta}  \cdot \frac{2 - \cos \theta + i\sin \theta}{2 - \cos \theta + i \sin\theta}\\
		&= \frac{4 -2\cos \theta +2 i\sin \theta}{5 -4 \cos \theta}.
	\end{align*}
	Therefore,
	$$
	\sum_{k=0}^\infty \frac{\cos k\theta}{2^k} = \frac{4-2\cos \theta}{5-4\cos\theta}.
	$$
	This is a periodic function of $\theta$.
\end{example}



\section{Application to alternating current}



In an alternating-current (AC) circuit, voltages and currents vary sinusoidally in time. A typical voltage source can be written as
\[
V(t) = V_0 \cos(\omega t + \phi),
\]
where $V_0$ represents the amplitude, $\omega$ is the angular frequency, and $\phi$ is the phase. Using Euler's formula,
\[
e^{i(\omega t + \phi)} = \cos(\omega t + \phi) + i \sin(\omega t + \phi),
\]
we can regard this voltage as the real part of a complex exponential:
\[
V(t) = \real\bigl(V_0 e^{i(\omega t + \phi)}\bigr).
\]
The complex quantity
\[
 V_0 e^{i\phi}
\]
is called a \emph{phasor}. Although the physical voltage is always a real-valued function, the complex exponential form is often easier to manipulate.

%As a simple example, consider a capacitor with capacitance $C$. The current through the capacitor is related to the voltage by
%\[
%I(t) = C \frac{dV}{dt}.
%\]
%If we represent the voltage by its phasor $\tilde{V}(t) = V_0 e^{i\omega t}$, then
%\[
%\frac{d}{dt} \tilde{V}(t) = i\omega V_0 e^{i\omega t},
%\]
%so the corresponding current phasor is
%\[
%\tilde{I}(t) = C \frac{d}{dt} \tilde{V}(t)
%= i\omega C\, \tilde{V}(t).
%\]
%Thus, in the phasor representation, the capacitor behaves like a multiplication by $i\omega C$. We can rewrite this as
%\[
%\tilde{V} = Z_C \tilde{I}, \qquad Z_C = \frac{1}{i\omega C},
%\]
%where $Z_C$ is the \emph{impedance} of the capacitor. A differential equation in time has been turned into a simple algebraic relation between complex numbers.

Figure~\ref{fig:phasor} illustrates the corresponding voltage phasor in the complex plane.



\begin{figure}[h]
	\centering
	\begin{tikzpicture}[>=Latex, scale=1.2]
		% Axes
		\draw[->] (-0.2,0) -- (3.2,0) node[right] {$\real$};
		\draw[->] (0,-0.2) -- (0,2.2) node[above] {$\imag$};
		
		% Phasor
		\draw[->, thick, blue] (0,0) -- (2,1.2) node[midway, above left] {$\tilde{V}$};
		
		% Angle
		\draw (1,0) arc[start angle=0, end angle=31, radius=1];
		\node at (1.1,0.35) {$\phi$};
		
		% Projection on real axis
		\draw[dashed] (2,1.2) -- (2,0);
		\node[below] at (2,0) {$V_0 \cos\phi$};
		
		% Projection on imaginary axis
		\draw[dashed] (2,1.2) -- (0,1.2);
		\node[left] at (0,1.2) {$V_0 \sin\phi$};
		
		% Magnitude
		\node[right] at (2,1.2) {$|\tilde{V}| = V_0$};
	\end{tikzpicture}
	\caption{Voltage phasor $\tilde{V} = V_0 e^{i\phi}$ in the complex plane.}
	\label{fig:phasor}
\end{figure}

\begin{example}
	Consider $n$ sinusoidal functions with the same frequency $\omega$,
	$$
	A_k \cos(\omega t + \theta_k),
	$$
	where amplitudes $A_k$ and phase $\theta_k$, for $k=1,2,\ldots, n$, are real numbers.  The sum of them is another sinusoidal function with the same frequency $\omega$. To find the amplitude and phase of the sum, we can simply add the $n$ complex numbers $A_k e^{i\theta_k}$, for $k=1,2,\ldots, n$, 
	$$
	\sum_{k=1}^n A_k e^{i\theta_k} \triangleq z_0.
	$$
	The modulus and argument of $z_0$ is then the amplitude and the phase of the resulting sinusoidal function
	$$
	\sum_{k=1}^n A_k \cos(\omega t + \theta_k) = \real\big(e^{i\omega t}\sum_{k=1}^n A_ke^{i\theta_k} \big) = \real(e^{i\omega t}z_0) =
	|z_0| \cos(\omega t + \arg(z_0)).
	$$
\end{example}





\section*{Summary}

\begin{itemize}
	\item The basic properties of complex sequence and series are the same as real sequence and real series.

	\item (Absolute convergence of power series) If a power series $\sum_{k=0}^\infty a_k z^k$ converges at $z_0\neq 0$, then it converges absolutely for all $z$ with $|z|<|z_0|$.

\item The region of convergence of power series is circular in shape. The radius of convergence is defined as
$$	R \triangleq \sup\big\{|z|:\, \sum_k a_kz^k \text{ converges} \big\}.$$

\item Hadamard formula for the radius of convergence: 
$$R = \frac{1}{\limsup_k |a_k|^{1/k}}.$$

\item ``A function $f(z)$ converges to $\infty$'' means that $1/f(z)$ converges to 0.

	\item Complex exponential function is defined by extending the power series of the real exponential function to complex variable. This power series is absolutely convergent for all complex numbers.
	
	\item A formula for computing complex exponential function: $e^{x+iy} = e^x(\cos y + i \sin y)$.
	
	\item Euler's formula: $e^{i\theta} = \cos \theta + i \sin \theta$, for all real number $\theta$.
\end{itemize}





\section*{Exercises}

\renewcommand{\labelenumi}{\arabic{enumi}.}
\begin{enumerate}
	

\item Determine whether the following sequences are convergent (converging to the point at infinity is allowed). If yes, what is the limit?
\begin{flalign*}
(a)\ & \lim_{n\rightarrow\infty} \frac{n^2+i}{i n^2+ 2}, & 
(b)\ & \lim_{n\rightarrow\infty} \sin(i\pi n), & 
(c)\ & \lim_{n\rightarrow\infty} \Big( \frac{z}{n} \Big)^n.& 
\end{flalign*}


	
\item Show that
$$
\sum_{k=0}^n \cos^2(k \theta ) = \frac{n+1}{2} + \frac{\sin((n+1)\theta)}{\sin \theta } \cos(n\theta) 
$$	
for real number $\theta$ that is not an integral multiple of $\pi/2$.

\item Evaluate the followings and express your answer in the form $a+bi$.
\begin{flalign*}
	\mathrm{(a)}\ & e^{3i\pi}, & \mathrm{(b)}\ & e^{1+\pi i /6}, &
	\mathrm{(c)}\ & \exp(\exp(\pi i/8)), & \mathrm{(d)}\ & e^{412\pi i / 6}. &&&
\end{flalign*}	

\item Let $(z_n)_{n\geq 1}$ and $(w_n)_{n\geq 1}$ be two convergent sequences of complex numbers. Prove that
$$
\lim_{n\rightarrow \infty} \frac{z_{n-1}w_0 + z_{n-2}w_1 + \cdots + z_0w_{n-1}}{n} = \big( \lim_{n\rightarrow \infty} z_n \big)  \big( \lim_{n\rightarrow\infty} w_n \big).
$$


\item  (A problem from Polya) Given two series $\sum_{n=0}^\infty a_n$ and $\sum_{n=0}^\infty b_n$, their Cauchy product is defined as the series $\sum_{n=0}^\infty$, with coefficient $c_n$ given by
$$
c_n = \sum_{k=0}^n a_k b_{n-k}.
$$

\begin{enumerate}
	\item Suppose that $\sum_{n=0}^\infty a_n$ is absolutely convergent. Prove that if $\sum_{n=0}^\infty$ is convergent (but not necessarily converging absolutely), then the series $\sum_{n=0}^\infty c_n$ converges.
	\item (hard) Show that if the Cauchy product of $\sum_{n=0}^\infty a_n$ and $\sum_{n=0}^\infty b_n$ is convergent whenever  $\sum_{n=0}^\infty b_n$ is convergent, then $\sum_{n=0}^\infty a_n$ is necessarily an absolutely convergent series.
\end{enumerate}

\item Prove the following identities about binomial coefficients.
\begin{align*}
\sum_{\substack{0\leq k \leq n \\
	k \equiv 0 \bmod 4}} \binom{n}{k} &= 2^{n-2}+2^{(n-2)/2} \cos \big( \frac{n\pi}{4}\big)\\
\sum_{\substack{0\leq k \leq n \\
		k \equiv 1 \bmod 4}} \binom{n}{k}&= 2^{n-2}+2^{(n-2)/2} \sin \big( \frac{n\pi}{4}\big)\\
\sum_{\substack{0\leq k \leq n \\
		k \equiv 3 \bmod 4}} \binom{n}{k}&= 2^{n-2}-2^{(n-2)/2} \sin \big( \frac{n\pi}{4}\big)
\end{align*}

\item ($\pi$ formulae) 
\begin{enumerate} 
	\item Compute the product $(2+i)(3+i)$ and hence prove
	$$
	\frac{\pi}{4} = \tan^{-1} \frac{1}{2} + \tan^{-1} \frac{1}{3}.
	$$
	
	\item Using similar technique, show that $\pi/4$ can be computed by
	$$
	\frac{\pi}{4} = 4 \tan^{-1} \frac{1}{5} - \tan^{-1} \frac{1}{239}.
	$$
	(Hint: use $(5-i)^4(1+i)$ )
	
	\item Use the result in part 2 and the power series expansion
	Use the power series expansion 
	$$
	\tan^{-1}(x) = x - \frac{x^3}{3} + \frac{x^5}{5} - \cdots
	$$	 
	to estimate the value of~$\pi$.
\end{enumerate}


\item The objective of this exercise is to prove 
$$
\sum_{k=0}^{n-1} \frac{1}{\cos^2 (\pi k /n)} = n^2.
$$
for odd~$n$.
\begin{enumerate}
	\item Suppose that $x_1$, $x_2,\ldots, x_n$ are the roots of polynomial $p(z) = a_0+a_1x+a_2x^2+\cdots + a_nx^n$. Assuming all the roots are nonzero, show that
	$$
	\sum_{k=1}^{n} \frac{1}{x_k^2} = \frac{a_1^2 - 2a_0a_2}{a_0^2}.
	$$	
	\item For integer $m\geq 1$, define the $m$-th Chebyshev polynomial $T_m(x)$ by the relation
	$$
	\cos (m\theta) = T_m(\cos \theta).
	$$
	Show that $\cos(\pi k/n)$, for $0\leq k \leq 2n-1$ are the roots of $T_{2n}(x)-1$.
	\item 	Using parts (a) and (b), derive
	$$
	\sum_{k=0}^{2n-1} \frac{1}{\cos^2 (\pi k /n)} = 2n^2.
	$$
\end{enumerate}


\end{enumerate}







%%%%%%%%%%%%%%%%%%%%%%%%%%%%%%%%%%



\chapter{Complex Trigonometric Function and Log Function}

We complexify the elementary functions in calculus.



\section{Trigonometric functions}

We extend the definition of trigonometric function and hyperbolic trigonometric functions to the complex plane.

\begin{svgraybox}
	\begin{definition} 	\label{def:cos_sin} \index{Trigonometric functions} For $z\in\mathbb{C}$,
		define the \df{complex sine}  and \df{complex cosine} functions by
		\begin{align*}
			\sin(z) & \triangleq \sum_{n=0}^\infty (-1)^n \frac{z^{2n+1}}{(2n+1)!} ,\\
			\cos(z) & \triangleq \sum_{n=0}^\infty (-1)^n \frac{z^{2n}}{(2n)!} .
		\end{align*}
	\end{definition}
\end{svgraybox}


By the limit ratio test (Prop.~\ref{prop:ratio_test}), both of them  converge absolutely on the whole complex plane. The power series reduce to the real trigonometric and real when $z$ is a real number. We can express the Euler's formula as

\begin{svgraybox}
	\begin{theorem}  	\label{thm:Euler2} \index{Euler's formula} For $z\in\mathbb{C}$,
		\begin{equation}e^{iz} = \cos(z) +i \sin(z).
			\label{eq:Euler_formula}
		\end{equation}
	\end{theorem}
\end{svgraybox}


\begin{proof} The proof is basically the same as the proof of Theorem~\ref{thm:Euler1}. We just need to apply the definition of $\cos(z)$ and $\sin(z)$ for complex number $z$.
	
	Starting from the power series expansion of $e^{iz}$
	$$
	e^{iz} = 1 + iz -\frac{z^2}{2!} - \frac{iz^3}{3!} + \frac{z^4}{4!} + \frac{iz^5}{5!} - \frac{z^6}{6!} + \cdots
	$$
	because of absolute convergence, we can re-arrange the terms and separate the terms without $i$ and the terms with $i$,
	\begin{align*}
		e^{iz} & \triangleq \sum_{n=0}^\infty \frac{(iz)^n}{n!} \\
		&= \sum_{k=0}^\infty (-1)^k \frac{z^{2k}}{(2k)!} + i
		\sum_{k=0}^\infty (-1)^k \frac{z^{2k+1}}{(2k+1)!} \\
		&= \cos(z) + i \sin(z).
	\end{align*}
	The last equality is true by the definition of complex sine and complex cosine functions.
\end{proof}



Using Theorem~\ref{thm:Euler2}, we can express sine function and cosine function in terms of the complex exponential function.

\begin{svgraybox}
	\begin{theorem} 	\label{thm:complex_cos_sin} For any $z\in \mathbb{C}$,
		\begin{align*}
			\cos(z) &= \frac{e^{iz} + e^{-iz}}{2 } \\
			\sin(z) &= \frac{e^{iz} - e^{-iz}}{2i } .
		\end{align*}
	\end{theorem}
\end{svgraybox}


\begin{proof}
	The complex cosine function is an even function, i.e., $\cos(-z) = \cos(z)$, because the power series that define $\cos(z)$ consists of even powers. Similarly, because all terms in the power series that define $\sin(z)$ have odd powers, the complex sine function is an odd function, i.e., $\sin(-z) = -\sin(z)$. This gives
	\begin{equation}
		e^{-iz} = \cos(-z) + i \sin(-z) = \cos(z) - i \sin(z)
		\label{eq:Euler_proof}
	\end{equation}
	
	By adding \eqref{eq:Euler_proof} to \eqref{eq:Euler_formula}, we get
	$$
	e^{iz} + e^{-iz} = {2} \cos(z).
	$$
	By subtracting \eqref{eq:Euler_proof} from \eqref{eq:Euler_formula}, we get
	$$
	e^{iz} - e^{-iz} = 2i \sin(z).
	$$
\end{proof}

The result in Theorem~\ref{thm:complex_cos_sin} holds for any complex numbers. In particular, when we restrict $z$ to a real number $\theta$, this yields
\begin{align}
	\cos(\theta) &= \frac{e^{i\theta} + e^{-i\theta}}{2 } , \\
	\sin(\theta) &= \frac{e^{i\theta} - e^{-i\theta}}{2i } .
\end{align}


\begin{remark}
	Some people take the two identities in Theorem~\ref{thm:complex_cos_sin} as the definition of complex sine and complex cosine, and then deduce Definition~\ref{def:cos_sin} as a corollary. Which definitions we adopt is a matter of taste. The resulting theory is the same.
\end{remark}


\begin{example} Evaluate $\sin(i)$.
	By Theorem~\ref{thm:complex_cos_sin},
	$$
	\sin(i) = \frac{e^{i\cdot i} - e^{-i \cdot i}}{2i} = \frac{e^{-1}- e^1}{2i} = i \frac{e^1 - e^{-1}}{2} = i \sinh(1).
	$$
	The trigonometric functions in the \textsf{numpy} library can compute complex number. We can verify the calculations in this example as follows.
\end{example}


\begin{example} For any real number $a$, we can evaluate $\cos(\pi+ai)$ by using the formula in Theorem~\ref{thm:complex_cos_sin}, 
	$$
	\cos(\pi+ai) = \frac{e^{i(\pi+ ai)} + e^{-i(\pi+ ai)}}{2} = \frac{e^{-a+i\pi}+ e^{a-i\pi}}{2} = - \frac{e^{-a} + e^{a}}{2} = -\cosh(a).
	$$
\end{example}






Some trigonometric identities in the real case continue to hold when the variables are complex numbers.

\renewcommand{\labelenumi}{(\alph{enumi}).}
\begin{svgraybox}
	\begin{theorem}
		For any complex numbers $z$ and $w$, the following hold:
		\begin{enumerate}
			\item $\cos^2 z + \sin^2 z = 1$.
			\item $\sin(z+w) = \sin z \cos w + \sin w \cos z$.
			\item $\cos(z+w) = \cos z \cos w - \sin z \sin w$.
		\end{enumerate}
	\end{theorem}
\end{svgraybox}

\begin{proof}
	We prove the first part. The proofs of the rest are similar.
	\begin{align*}
		L.H.S. & = \cos^2 z + \sin^2 z \\
		&= \Big( \frac{e^{iz} + e^{-iz}}{2}\Big)^2 + 
		\Big( \frac{e^{iz} - e^{-iz}}{2i}\Big)^2 \\ 
		&= \frac{1}{4} [(e^{iz}+e^{-iz})^2 -(e^{iz}-e^{-iz})^2] \\
		&= \frac{1}{4} [ e^{2iz} + 2 + e^{-2iz} - e^{2iz} + 2 - e^{-2iz}]\\
		&= 1 = R.H.S.
	\end{align*}
\end{proof}


The complex hyperbolic sine and cosine functions are defined similarly by power series.

\begin{svgraybox}
	\begin{definition} 	\label{def:cosh_sinh} \index{Hyperbolic trigonometric functions} For $z\in\mathbb{C}$,
		The \df{complex hyperbolic sine}  and \df{complex hyperbolic cosine} functions  are defined as
		\begin{align*}
%			\sin(z) & \triangleq \sum_{n=0}^\infty (-1)^n \frac{z^{2n+1}}{(2n+1)!} ,\\
%			\cos(z) & \triangleq \sum_{n=0}^\infty (-1)^n \frac{z^{2n}}{(2n)!} .
			\sinh(z) & \triangleq \sum_{n=0}^\infty  \frac{z^{2n+1}}{(2n+1)!}, \\
			\cosh(z) & \triangleq \sum_{n=0}^\infty  \frac{z^{2n}}{(2n)!} .
		\end{align*}
	\end{definition}
\end{svgraybox}



The hyperbolic trigonometric functions are closely related to the complex sine and cosine functions.
\begin{svgraybox}
	\begin{theorem}	\label{thm:sinh_sin}  For all complex number $z\in\mathbb{C}$,
		\begin{align*}
			\sin(iz) & = i \sinh(z) ,\\
			\cos(iz) &= \cosh(z).
		\end{align*}
	\end{theorem}
\end{svgraybox}

\begin{proof}
	The first identity can be proved by considering the power series expansion of $\sin(iz)$,
	\begin{align*}
		\sin(iz) &\triangleq  \sum_{k=0}^\infty (-1)^k \frac{(iz)^{2k+1}}{(2k+1)!} 
		=  i\sum_{k=0}^\infty (-1)^k \frac{i^{2k}z^{2k+1}}{(2k+1)!} 
		=  i\sum_{k=0}^\infty \frac{z^{2k+1}}{(2k+1)!} 
		= i \sinh(z).
	\end{align*}
	The proof of the second identity is similar.
\end{proof}

In particular, for any real number $x$, we have
$$
\sin(ix) =  i \sinh(x), \ \ \text{ and } \cos(ix) = \cosh(x).
$$

\begin{svgraybox}
	\begin{theorem} For any complex number $z=x+iy$, we have
\begin{align*} 
\cos(z) &= \cos x \cosh y - i \sin x \sinh y, \\
\sin(z) &= \sin x \cosh y + i \cos x \sinh y.
\end{align*}
	\end{theorem}
\end{svgraybox}


\section{Log function}

The complex exponential function has a complex period $2\pi i$; that is, for any $z\in\mathbb{C}$,
$$
e^{z+2\pi k i} = e^{z},\ \ \text{ for } k \in \mathbb{Z}.
$$
As a result, the inverse function of $e^z$ is multi-valued. 

\noindent {\em Notation:} We will use the notation $\ln x$ for the natural log function with positive real number as input.

\begin{example}
	Find $z \in \mathbb{C}$ such that $e^z = 1+i$.	
	
	The complex number $1+i$ in polar coordinates is
	$$
	\sqrt{2}(\cos(\pi/4)+i\sin(\pi/4)).
	$$ 
	We want to find all complex numbers $x+iy$ such that
	$$
 e^x(\cos y + i \sin y) = \sqrt{2}(\cos(\pi/4)+i\sin(\pi/4)).
	$$
	By comparing the moduli and arguments, we have
	\begin{align*}
		x &= \frac{1}{2} \ln(2) \\
		y &= \pi/4 + 2\pi k, \ \ \text{for } k=0, \pm1, \pm2, \ldots
	\end{align*}
	Hence
	$$
	\log(1+i) = \frac{1}{2} \ln(2) + i (\pi/4 + 2\pi k)
	$$
	for $k\in\mathbb{Z}$.
\end{example}


In general, given a complex number $r(\cos\theta+i\sin\theta)$ in polar form, the complex log of $r(\cos\theta+i\sin\theta)$ can take values
$$
\ln r + i(\theta + 2\pi k)
$$
for any integer $k$. In general, we have the following

\begin{svgraybox}
	\begin{definition} \index{Log function}
		For nonzero $w\in\mathbb{C}$, the \df{complex log function} is a multi-function defined as
		$$
		\log(w) \triangleq \ln |w| + i(\arg(w)+2\pi k)
		$$
		where $k=0$, $\pm 1$, $\pm 2, \ldots$ 
		\label{def:complex_log}
	\end{definition}
\end{svgraybox}


The function $\arg(w)$ in Definition \ref{def:complex_log} is the multi-valued argument function. If we specify a range for the angle by take the principal argument (Definition \ref{def:principal_argument}), we have a uniquely defined function.

\begin{svgraybox}
	\begin{definition} \index{Log function!principal branch}
		The function
		$$
		\Log(w) \triangleq \ln |w| + i\Arg(w),
		$$
		is called the \df{principal value} or the \df{principal branch} of complex log function. We will denote the principal log function by $\Log(z)$.
		\label{def:log_prinicipal_value}
	\end{definition}
\end{svgraybox}



The complex log function extends the domain of the real log function to the whole complex plane except the origin. We note that it is not possible to define the log(0) even if we extend the number system to complex numbers, because the complex exponential function is never equal to zero. One way to see it is by the following general fact
$$
\forall z \in \mathbb{C},\ \ \ e^z e^{-z} = e^{z-z} = e^0 = 1.
$$
If we assume $e^z=0$ for some $z$, then we would obtain $0\cdot 0 = 1$, which cannot be true. This means that $e^z\neq 0$ for all $z\in \mathbb{C}$.


Using the complex log function we can compute the log of a negative number.

\begin{example}
	Compute the principal log of $-10$, using a principal argument function with range $[0,2\pi)$.
	
	We first write $-10$ as
	$$
	-10 = 10 (\cos \pi + i \sin \pi)
	$$
	If $e^{x+iy} = 10 (\cos \pi + i \sin \pi)$, we must have $10 = e^x$, giving $x= \ln 10$. The imaginary part $y$ equals $\pi$, which is the unique argument in the range $[0,2\pi)$. Therefore, the principal log of $-10$ is
	$$
	\Log(-10) = \ln(10) + i \pi.
	$$
	


\section{Power function}

We can define the complex power function via the complex log function.

\begin{svgraybox}
	\begin{definition} \index{Complex power}
		Given a nonzero complex number $z$ and a complex number $w$, we define the \df{complex power} $z^w$ by
		$$
		z^w \triangleq e^{w \log(z)}
		$$
		where $\log(z)$ is the complex log function.
		\label{def:complex_power}
	\end{definition}
\end{svgraybox}


\begin{example} Compute $2^i$.
	\begin{align*}
		2^i \triangleq \exp(i \log 2) &= \exp(i (\ln 2 + i2\pi k))\\
		&= \exp(-2\pi k + i \log 2) \\
		&= e^{-2\pi k}(\cos(\ln 2)+i\sin(\ln 2)),
	\end{align*}
	where $k$ can take any integer as its value. We can consider the value when $k=0$ as the principal value. If we want to write down just one answer, we can write
	$$
	2^i = \cos(\ln 2)+ i \sin(\ln 2).
	$$
\end{example}

We can use the complex power function to compute the $n$-th roots of a complex numbers. We can compare this method with the DeMoivre formula in \eqref{eq:DeMoivre}.

\begin{example} Calculate the cube roots of $2i$ by $(2i)^{1/3} \triangleq \exp(\log(2i)/3)$.
	
	We compute
	\begin{align*}
		\log(2i) &= \ln 2 + i \big( \pi/2 + 2\pi k\big), \quad \text{ for } k \in \mathbb{Z} \\
		\frac{\log(2i)}{3} &= \frac{\ln 2}{3 }+ i \big( \pi/6 + 2\pi k/3 \big).
	\end{align*}
	Exponentiating both sides, we get
	\begin{align*}
		\exp\big( \frac{\log(2i)}{3}\big) &=
		\sqrt[3]{2} \Big( \cos \big( \frac{\pi}{6}+ \frac{2\pi k}{3} \big)+ i  \sin \big( \frac{\pi}{6}+ \frac{2\pi k}{3} \big) \Big) 
	\end{align*}
	for $k=0,1,2$. There are three answers, all with magnitude $\sqrt[3]{2}$.
\end{example}

The next example has complex base and complex exponents at the same time.

\begin{example}
	Compute $i^i$ by calculating $e^{i\log i}$.
	\begin{align*}
		\exp(i \log i) &= \exp\Big(i (\log 1 + i(\frac{\pi}{2}+2\pi k))\Big)\\
		&= \exp\big(-(\frac{\pi}{2}+2\pi k)\big) ,
	\end{align*}
	where $k \in \mathbb{Z}$.  All possible values are real numbers. 
	By taking $k=0$, we can say that the principal value of $i^i$ is $e^{-\pi/2}$.
\end{example}


\section{Inverse of complex trigonometric functions}

Inverse trigonometry function can be computed by solving quadratic equation.

\begin{example}
	Find all complex numbers $z$ such that $\cos(z)=2$.
	
	Recall from Theorem~\ref{thm:complex_cos_sin} that we can express $\cos(z)$ as
	$$
	\cos(z) = \frac{e^{iz}+e^{-iz}}{2}.
	$$
	Let $w=e^{iz}$. This amounts to solve
	$$
	\frac{w +w^{-1}}{2} = 2.
	$$
	This is a quadratic equation
	$w^2-4w+1=0$. The roots are $2 \pm \sqrt{3}$. Hence 
	$$e^{iz} = 2\pm \sqrt{3}.$$
	By taking the complex log function, we obtain
	$$
	z = -i \big[\ln(2\pm\sqrt{3})+2\pi i k \big] = -i\ln(2\pm \sqrt{3})+2\pi k,
	$$
	where $k\in\mathbb{Z}$.
	

\end{example}


\begin{example}
	Solve $\cot(z)=2i$ for $z\in\mathbb{C}$.
	
	This is equivalent to solving
	$$
	\frac{e^{iz}+e^{-iz}}{2} \frac{2i}{e^{iz}-e^{-iz}} = 2i.
	$$
	By letting $w=e^{iz}$, we can simplify it to
	\begin{align*}
		\frac{w+w^{-1}}{w-w^{-1}} &= 2 \\
		\frac{w^2+1}{w^2-1} &= 2.
	\end{align*}
	The values for $w$ is $\pm \sqrt{3}$. Then, we can obtain $z$ by taking complex log
	\begin{align*}
		z &= -i\log(\pm\sqrt{3}) 
		=-i(\ln\sqrt{3} + i \pi k) 
		= \pi k - i \ln \sqrt{3},
	\end{align*}
	where $k$ is any integer.
\end{example}







\section{Visualization of complex function using phase portrait}
\index{Phase portrait} 
Phase portrait, which is also known as domain coloring, plots the argument of a complex function on a plane, by associating the numbers in the range $[0,2\pi)$ with the color spectrum. The brightness corresponds to the magnitude of the complex function. A dark point means the function is zero. A bright color indicates the area where function magnitude is large. We take the phase portrait of the identity function as the reference (Fig.~\ref{fig:phase_portrait1}).

The phase portrait of some complex functions are shown in Fig.~\ref{fig:phase_portraits}. The picture for the function $f(z)=1/z$ has a bright point near the origin, because this function has large magnitude near the origin. We also note that as we go around the origin, we see the color in the reverse order. It is because the function $f(z)=1/z$ multiply the argument by $-1$. In the phase portrait of $f(z)=z^2$, we see a zero of order 2 at the origin. As we go around the origin, we meet all the colors twice.

\begin{figure}
	\begin{center}
		\includegraphics[width=6cm]{phase_portrait_1.png}
	\end{center}
	\caption{Phase portrait of $f(z)=z$}
	\label{fig:phase_portrait1}
\end{figure}

\begin{figure}
	\begin{subfigure}[b]{0.4\textwidth}
		\includegraphics[width=6cm]{phase_portrait_2.png}
		\caption{Phase portrait of $f(z)=1/z$}
	\end{subfigure}
	\hfill
	\begin{subfigure}[b]{0.4\textwidth}
		\includegraphics[width=6cm]{phase_portrait_3.png}
		\caption{Phase portrait of $f(z)=z^2$}
	\end{subfigure}
	\bigskip
	
	\begin{subfigure}[b]{0.4\textwidth}
		\includegraphics[width=6cm]{phase_portrait_4.png}
		\caption{Phase portrait of \\ $f(z)=(z^2+9)/(z-2)$}
	\end{subfigure}
	\hfill
	\begin{subfigure}[b]{0.5\textwidth}
		\begin{center}
			\includegraphics[width=6cm]{phase_portrait_5.png}
		\end{center}
		\caption{Phase portrait of $f(z)=\tan(z)$}
	\end{subfigure}
	\caption{Phase portraits.}
	\label{fig:phase_portraits}
\end{figure}







\section{Numerical calculations}
Python support complex trigonometric functions. For example, we can check that $\sin(i) = i \sinh(1)$ by

\begin{lstlisting}[language = Python, frame=single]
In [0]: from numpy import sin, sinh
In [1]: sin(1j)
Out[1]: 1.1752011936438014j

In [2]: sinh(1)
Out[2]: 1.1752011936438014
\end{lstlisting}


The \textsf{arccos} function in Python is able to compute the complex arc cosine function. We can check that it returns $-i\log(2+\sqrt{3})$ as the answer.
\begin{lstlisting}[language=Python, frame=single]
In [0]: from numpy import arccos, log
In [1]: arccos(2+0j)
Out[1]: -1.3169578969248166j

In [2]: -log(2+sqrt(3))*1j
Out[2]: (-0-1.3169578969248166j)
\end{lstlisting}

The other solution is $-i\log(2-\sqrt{3})$.

\begin{lstlisting}[language=Python, frame=single]
In [3]: -log(2-sqrt(3))*1j
Out[3]: 1.3169578969248164j
In [4]: cos(-log(2-sqrt(3))*1j)
Out[4]: (1.9999999999999996-0j)
In [5]: cos(-log(2+sqrt(3))*1j)
Out[5]: (1.9999999999999998-0j)
\end{lstlisting}


\begin{comment}
\section{Appendix: A unified approach to limit and convergence using filter}

There are four types of limits in the extended complex number system:
\begin{flalign}
	& \lim_{z\rightarrow z_0} f(z) = w_0, & 
	& \lim_{z\rightarrow z_0} f(z)= \infty, & 
	& \lim_{z\rightarrow \infty} f(z)= w_0 ,& 
	& \lim_{z\rightarrow \infty} f(z)= \infty. & &&&&
	\label{eq:four_limits}
\end{flalign}
The variable $z_0$ and $w_0$ are finite complex numbers. The four limits share a lot of similarity, but are not exactly the same. One way to understand these four notions of limit is to use the inversion function, and reduce the last three to the first one.

The meaning of ``$f(z)$ converges to $w_0$ as $z$ converges to $z_0$'' is defined formally as
$$
\forall \epsilon>0, \exists \delta>0 \ s.t.\ |f(z)-w_0| \leq \epsilon \text{ for all } 0< |z-z_0|< \delta.
$$
We note that we do not require ``$0= |z-z_0|$'' in the above definition, and hence $f$ may not be defined at the point~$z_0$.

The other three combinations are listed below.
\begin{itemize}
	\item $\lim_{z\rightarrow z_0} f(z)= \infty$ means $\lim_{z\rightarrow z_0} \frac{1}{f(z)}= 0$.
	\item $\lim_{z\rightarrow \infty} f(z)= w_0$ means $\lim_{z\rightarrow 0} f(1/z) = w_0$.
	\item $\lim_{z\rightarrow \infty} f(z)= \infty$ means $\lim_{z\rightarrow 0} \frac{1}{f(1/z)} = 0$.
\end{itemize}



%While this approach is logically valid, it is not convenient when we come to proof. 
Suppose $f(z)$ and $g(z)$ are functions whose domains and ranges are the extended complex system. If $g(z)$ converges to $w_0$ as $z\rightarrow z_0$, and $f(w)$ converges to $u_0$ as $w\rightarrow w_0$, then we expect that the composite function $f(g(z))$ converges to $u_0$, as $z\rightarrow z_0$, i.e.,
\begin{equation}
	\lim_{z\rightarrow z_0} f(g(z)) = u_0.
	\label{eq:limit_compose}
\end{equation}
If we use the above definitions of limits, a proof of this fact requires the consideration of eight cases, depending on whether $z_0$, $w_0$ and $u_0$ are the point at infinity or not. This is certainly not the best way to get things done.

\begin{remark}
	The same problem is even more severe in calculus. We also have convergence from the left, convergence from the right, convergence to infinity, and convergence to negative infinity. Basically the idea is the same, but an elementary book on calculus would list all combinations one by one, give one proof in the simplest case, and then say that the other cases can be proved similarly.
\end{remark}



A better way is to have a single definition of limit that includes all possibilities in \eqref{eq:four_limits}. Then we just need one high-level proof to cover all cases. We introduce the notion of ``filter'' from topology. This is relatively more abstract, but it is the precise level of abstraction that is required to simplify things. We recall that a topological space is a pair $(X, O(X))$, consisting of a ground set $X$ and the collection $O(X)$ of all open subsets in $X$. In our application, we are primarily interested in the complex plane $\mathbb{C}$ as the ground set. The collection of all open sets in $\mathbb{C}$ is denoted by $O(\mathbb{C})$. It is known that the intersection of finitely many open sets is an open set, and an arbitrary union of open sets is also open.

In the following definition,
we abstract the crucial properties of the collection of all open neighborhoods of a point. The abstract concept is called {\em filter}. We refer the readers to~\cite{Berge} for a treatment of topological spaces based on filter. 


\begin{svgraybox}
	\begin{definition} \label{def:filter} Given a collection of open sets $O(X)$ in a topological space $X$,
		a \df{filter} $\mathcal{F}$ is a nonempty collection of subsets of $X$ that is closed under taking finite intersection and going upward. This is, $\mathcal{F}$ is a nonempty subset of the power set $\mathcal{P}(X)$ of $X$, satisfying the following requirements:
		\begin{enumerate}
			\item If $A$ and $B$ are in $\mathcal{F}$, then $A\cap B$ is in $\mathcal{F}$, and
			\item If $A$ is in  $\mathcal{F}$ and $B$ is in $O(X)$ and $A\subseteq B$, then $B\in \mathcal{F}$.
		\end{enumerate}
	\end{definition}
\end{svgraybox}

%
%\begin{remark}
%	We note that conditions 1 and 2 in the above definition involve set operation $\cap$ and $\subseteq$. In general, the notion of filter is defined with $O(X)$ replaced by a lattice $(L, \wedge, \vee)$. However, in this note, we only discuss filter in the context of complex plane.
%\end{remark}



\begin{example}
	On the real number line $\mathbb{R}$, all open sets that contained a given point $x_0$ form a filter.
\end{example}

\begin{example}
	Given a set $B$ in a topological space $X$, the collection of sets
	$$
	\{S \subseteq X:\, \text{$S$ is open, } S \supseteq B\}
	$$
	is called the {\em principal filter} of $B$.
\end{example}


\begin{example}
	On the complex plane $\mathbb{C}$, all open sets that contained a given complex number $z_0$ form a filter. We will call this filter $\mathcal{F}_{z_0}$. 
\end{example}




\begin{example}
	Given complex number $z_0$, the collection 
	$$
	\mathcal{F}_{z_0}' \triangleq \{A \setminus \{z_0\}:\, \textrm{$A$ is an open set in $\mathbb{C}$ containing $z_0$} \}
	$$	
	is a filter in the relative topology on the set $\{z_0\}^c$. We call this filter a {\em punctured filter}, and use the symbol $\mathcal{F}_{z_0}'$. This filter capture the behavior of a function near~$z_0$ but not equal to~$z_0$.
\end{example}


\begin{example}
	On the complex plane $\mathbb{C}$, the collection of sets
	$$
	\{A \subseteq \mathbb{C}:\, \text{$A$ is open, } \exists R>0 \text{ s.t. } ( |z|>R \Rightarrow z \in A ) \}
	$$
	is a filter. We remark that the radius $R$ may depend on the set $A$. A set $A$ is in this collection of sets if the complement $A^c$ is bounded. The sets in this filter are the open neighborhoods of the point at infinity. We denote this filter by $\mathcal{F}_{\infty}'$. This is the collection of all open neighborhoods of the point at infinity in $\mathbb{C}$. Since the point at infinity does not belong to $\mathbb{C}$, this is a punctured filter.
\end{example}

%\begin{remark} Given a principal filter $\mathcal{F}$ of a point $z_0$, the intersection of all subsets in $\mathcal{F}$ is the corresponding point in~$\mathbb{C}$. Hence, we can recover the point $z_0$ from the sets in the principal filter. In this sense, we can regard a filter as a ``generalized element''.
%\end{remark}

\begin{svgraybox}
	\begin{definition}
		Suppose $f:\mathbb{C}\rightarrow \mathbb{C}$ is a function in complex variable, $\mathcal{F}$ and $\mathcal{G}$ are filters in $\mathbb{C}$. We say that \df{$f$ converges to filter $\mathcal{G}$ along filter $\mathcal{F}$} if for all $V \in \mathcal{G}$, the pre-image $f^{-1}(V)$ of $V$ under the function $f$ is a member of $\mathcal{F}$.
	\end{definition}
\end{svgraybox}

Using this terminology, we can interpret the four types of limits in a unified manner:
\begin{enumerate}
	\item  $\lim_{z\rightarrow z_0} f(z) = w_0$ is the same as ``$f(z)$ converges to filter $\mathcal{F}_{w_0}'$ along filter $\mathcal{F}_{z_0}'$''.
	\item  $\lim_{z\rightarrow z_0} f(z) = \infty$ is the same as ``$f(z)$ converges to filter $\mathcal{F}_{\infty}'$ along filter $\mathcal{F}_{z_0}'$''.
	\item  $\lim_{z\rightarrow \infty} f(z) = w_0$ is the same as ``$f(z)$ converges to filter $\mathcal{F}_{w_0}'$ along filter $\mathcal{F}_{\infty}'$''.
	\item  $\lim_{z\rightarrow \infty} f(z) = \infty$ is the same as ``$f(z)$ converges to filter $\mathcal{F}_{\infty}'$ along filter $\mathcal{F}_{\infty}'$''.
\end{enumerate}

We note that the filters above are all punctured filters. We need to use punctured filters because the function $f$ has type $\mathbb{C}\rightarrow\mathbb{C}$, and hence does not take $\infty$ as its value. For example, in the second limit above, we statement ``$f(z_0)=\infty$'' has no meaning as $\infty$ is not a complex number in $\mathbb{C}$. The filter $\mathcal{F}_\infty'$ captures the behavior of function $f(z)$ when $z$ is close to the point at infinity.

We can now give a short proof of \eqref{eq:limit_compose}. The proof only requires some argument with set operations.

\begin{svgraybox}
	\begin{proposition}
		Let $f$ and $g$ denote two complex functions from $\mathbb{C}$ to $\mathbb{C}$. Let $\mathcal{F}$, $\mathcal{G}$ and $\mathcal{H}$ be three filters on $\mathbb{C}$. If $g$ converges to $\mathcal{G}$ along $\mathcal{F}$, and $f$ converges to $\mathcal{H}$ along $\mathcal{G}$, then the composite function $f(g(z))$ converges to $\mathcal{H}$ along $\mathcal{F}$.
	\end{proposition}
\end{svgraybox}

\begin{proof}
	Let $V$ be a set in $\mathcal{H}$. Since $f$ converges to $\mathcal{H}$ along $\mathcal{G}$, we have $f^{-1}(V) \in \mathcal{G}$. Now, by the assumption that  $g$ converges to $\mathcal{G}$ along $\mathcal{F}$, we have $g^{-1}(f^{-1}(V))\in\mathcal{F}$. Finally, we note that the inverse of the composition $f\circ g$ is $g^{-1} \circ f^{-1}$ and hence $g^{-1}(f^{-1}(V)) = (f\circ g)^{-1}(V)$. This proves that $(f\circ g)^{-1}(V)$ is a set in $\mathcal{F}$.
\end{proof}


Filter becomes increasingly important in recent years, due to its use in computer formalization of mathematics. See 
\href{https://leanprover-community.github.io/theories/topology.html}{Maths in Lean: Topological, uniform and metric spaces} 

\end{comment}



If we directly ask Python to compute $\log(-10)$, we will get a warning, saying that there is a type mismatch.
\begin{lstlisting}[language=Python, frame=single]
In [0]: from numpy import log
In [1]: log(-10)
RuntimeWarning: invalid value encountered in log
Out[1]: nan
\end{lstlisting}


One way to compute it works is to manually convert the input to a complex number. 
\begin{lstlisting}[language=Python, frame=single]
In [2]: log( complex(-10) )
Out[2]: (2.302585092994046+3.141592653589793j)
\end{lstlisting}
\end{example}

An alternate way is to use the complex log function in the \textsf{cmath} library. It automatically converts the input value to the complex number data type.
\begin{lstlisting}[language=Python,frame=single]
In [1]: from cmath import log
In [2]: log(-10)
Out[2]: (2.302585092994046+3.141592653589793j)
\end{lstlisting}

The following Python commands compute the principal value of $2^i$.
\begin{lstlisting}[language=Python,frame=single]
In [0]: from numpy import cos, sin, log
In [1]: 2**1j
Out[1]: (0.7692389013639721+0.6389612763136348j)

In [2]: cos(log(2))+1j*sin(log(2))
Out[2]: (0.7692389013639721+0.6389612763136348j)	
\end{lstlisting}




The value of $i^i$ is a real number $e^{-\pi/2}$.
\begin{lstlisting}[language=Python,frame=single]
In [0]: from numpy import exp, pi
In [1]: 1j**1j
Out[1]: (0.20787957635076193+0j)

In [2]: exp(-pi/2)
Out[2]: 0.20787957635076193
\end{lstlisting}


Relation between complex log function and arctan function.
\begin{lstlisting}[language=Python, frame=single]
In [0]: from numpy import log, arctan

In [1]: -1j*log(3**(1/2))
Out[1]: -0.5493061443340548j

In [2]: arctan(1/2j)
Out[2]: -0.5493061443340549j
\end{lstlisting}






\section{Summary}


\begin{itemize}
	\item The complex version of $\sin(z)$, $\cos(z)$, $\sinh(z)$, and $\cosh(z)$ are defined by power series, extending the domain to the whole complex plane.
	
	\item The Euler's formula $e^{iz} = \cos(z) + i \sin(z)$ holds for all complex number $z\in\mathbb{C}$.
	
	\item The complex exponential function is related to the complex trigonometric functions by
	$$
	\cos(z) = \frac{e^{iz}+ e^{-iz}}{2}, \quad \text{ and } 
	\sin(z) = \frac{e^{iz}- e^{-iz}}{2i}.
	$$
	
	\item The complex log function is a multi-function:
	$$
	\log(z) \triangleq \ln |z| + i(\arg(z)+2\pi k),
	$$
	defined for $z\neq 0$, where $k=0$, $\pm 1$, $\pm 2, \ldots$
	
	\item The complex power is defined in terms of log function,
	$$
	z^w \triangleq e^{w \log(z)},
	$$	
	and hence is a multi-function.
\end{itemize}





\section*{Exercises}


\renewcommand{\labelenumi}{\arabic{enumi}.}
\begin{enumerate}
	
\item Evaluate $\sin(z)$, $\cos(z)$, $\tan(z)$ for the following values of $z$:
\begin{flalign*}
\mathrm{(a)} \ & 1-i,  & \mathrm{(b)}\ & 1+i & \mathrm{(c)} \ &\frac{\pi}{4}i. &&&&&
\end{flalign*}	

\item In each part, find all complex number $z$ that satisfy the equation.
\begin{flalign*}
(a)\ & z^{\sqrt{2}} = 2i & (b)\ & \cos(z) = -4 & (c)\ & e^z = 1+\sqrt{3}i &&&&&\\
(d)\ & z^i = 3 & (e)\ & \ln(z) = \pi i &&&&&
\end{flalign*}

\item Find a solution to the following equation
$$
i^z = 1+i,
$$
and verify your answer with Python.


\item Define
\begin{flalign*}
\tanh z &  \triangleq \frac{\sinh z}{\cosh z} \quad (\text{for } \cosh z\neq 0), &
\coth z &  \triangleq \frac{\cosh z}{\sinh z} \quad (\text{for } \sinh z\neq 0), &&&& 
%\sech z &  \triangleq \frac{1}{\cosh z} \quad (\text{for } %\cosh z\neq 0), &
%\csch z &  \triangleq \frac{\cosh z}{\sinh z} \quad (\text{for } \sinh z\neq 0). &&&& 
\end{flalign*}
Prove that
\begin{enumerate}
	\item $\tanh(iz) = i \tan z$,
	\item $\tan(iz) = i \tanh z$,
	\item $\coth(iz) = -i \cot z$,
\item $\cot(iz) = -i \coth z$.
\end{enumerate}

\item Solve
$$
\cos z + i \sin z = 1/\sqrt{3}.
$$



\item In this exercise we derive some trigonometric identities using complex roots of unity.
\begin{enumerate}
\item By expressing $\tan(\pi/5)$ and $\sin(\pi/10)$ in terms of the 5-th root of unity $e^{2\pi i/5}$, show that
\begin{align*}
	\tan\big( \frac{\pi}{5}\big) &= \sqrt{5 - 2\sqrt{5}} \\
	\sin\big( \frac{\pi}{10}\big) &= \frac{\sqrt{5}-1}{4}.
\end{align*}

\item By using the $7$-th root of unity $ e^{2\pi i / 7}$, prove
$$
\tan\big( \frac{\pi}{7} \big) + 
\tan\big( \frac{2\pi}{7} \big) + 
\tan\big( \frac{3\pi}{7} \big) = \sqrt{7}.
$$
\end{enumerate}

\item Derive the following formulae for inverse trigonometric functions.
\begin{enumerate}
	\item $\arcsin(z) = -i \log(\sqrt{1-z^2} +iz)$.
\item $\arccos(z) = -i \log(z+i\sqrt{1-z^2} )$.	
\item $\arctan (z) = \frac{1}{2i} \log  {\frac{1+iz}{1-iz}}$.
\end{enumerate}

\item
Show that the zeros of $\sin(z)$ and $\cos(z)$ are located on the real axis.
\begin{enumerate}
	\item  $\sin z=0$ if and only if $z=k\pi$ for some integer $k$.
	\item $\cos z=0$ if and only if $z= (\pi/2) + k\pi$ for some integer $k$.
\end{enumerate}


	
%\item (Convergence of power series on the boundary) \begin{enumerate} 
%	\item Show that the power series $\sum_{n=0}^\infty z^n$ diverges for all $z$ on the unit circle
%	
%	\item Show that the power series $\sum_{n=1}^\infty \frac{1}{n}z^n$ diverges at $z=1$, but converges for all other points on the unit circle.
%	
%	\item Show that the power series $\sum_{n=1}^\infty \frac{1}{n^2}z^n$ converges at all points on the unit circle.
%	\end{enumerate}
\end{enumerate}



